% concatenated from Arxiv.tex, Introduction.tex, Background.tex, Compression.tex, main-datageneration.tex, main-scalinglaw.tex, Discussion.tex, conclusion.tex, App-LLLcompressor.tex, App-kolmogorovstructure.tex, Data_Model.tex, Scaling_law.tex, Instruction_tuning.tex, Experiment.tex, Appendix.tex
\documentclass[11pt]{article}
\usepackage{setspace}
\usepackage{url}            % simple URL typesetting
\usepackage{booktabs}       % professional-quality tables
\usepackage{amsfonts}       % blackboard math symbols
\usepackage{nicefrac}       % compact symbols for 1/2, etc.
\usepackage{microtype}      % microtypography
\usepackage{xcolor}         % colors
\usepackage{amsmath}
\usepackage{amsthm}
\usepackage{graphicx}
\usepackage{subcaption}
\usepackage{bbm}
\usepackage[round]{natbib}
\usepackage{hyperref}   
\usepackage[capitalize,noabbrev]{cleveref}
\usepackage{thmtools}
\usepackage{thm-restate}
\usepackage{amssymb}
\usepackage{comment}
\usepackage{mathtools}
\usepackage[normalem]{ulem}

\usepackage{environ}
\usepackage{ifthen}
\newboolean{nips}
\setboolean{nips}{false}
\NewEnviron{nips}{
  \ifthenelse{\boolean{nips}}{
    \BODY
  }{}
}
\NewEnviron{arxiv}{
  \ifthenelse{\not\boolean{nips}}{
    \BODY
  }{}
}
\usepackage{geometry}
\geometry{
     left=1in,
     right=1in,
     top=1in,
     bottom=1in,
     footskip=0.5in
}

\setlength{\floatsep}{5pt plus 1pt minus 1pt}

\usepackage{amsmath}
\usepackage{amssymb}
\usepackage{mathtools}
\usepackage{amsthm}
\usepackage{ulem}

\theoremstyle{plain}
\newtheorem{theorem}{Theorem}[section]
\newtheorem{hypothesis}[theorem]{Hypothesis}
\newtheorem{proposition}[theorem]{Proposition}
\newtheorem{lemma}[theorem]{Lemma}
\newtheorem{corollary}[theorem]{Corollary}
\newtheorem{remark}[theorem]{Remark}
\theoremstyle{definition}
\newtheorem{definition}[theorem]{Definition}
\newtheorem{assumption}[theorem]{Assumption}
\newtheorem{example}[theorem]{Example}
\newenvironment{proof_sketch}{\noindent{\textbf{Proof Sketch}.}}{\hfill$\square$}

\crefname{assumption}{Assumption}{Assumptions}



%% spacing around \paragraph
\makeatletter
\renewcommand{\paragraph}{%
  \@startsection{paragraph}{4}%
  {\z@}{1ex \@plus 1ex \@minus .2ex}{-1em}%
  {\normalfont\normalsize\bfseries}%
}
\makeatother

\newenvironment{subproof}[1][\proofname]{%
  \renewcommand{\qedsymbol}{$\square$}%
  \begin{proof}[#1]%
}{%
  \end{proof}%
}

\newenvironment{ssubproof}[1][\proofname]{%
  \renewcommand{\qedsymbol}{$\blacksquare$}%
  \begin{proof}[#1]%
}{%
  \end{proof}%
}

\newcounter{mycounter}
\setcounter{mycounter}{1}





\newcommand{\dd}{\mathrm{d}}

\newcommand{\tokenset}{\mathbb{V}}
\newcommand{\Voc}{\mathbb{V}}
\newcommand{\Vocsyn}{\mathbb{V}_{\text{syn}}}
\newcommand{\Vocknow}{\mathbb{V}_{\text{knw}}}
\newcommand{\numsyntax}{n_s}
% \newcommand{\synPrior}{\mathcal{P}_{\text{syn}}}

\newcommand{\capacity}{C}

\newcommand{\DKL}{D_{\mathsf{KL}}}
\newcommand{\redundency}{\mathsf{Red}}
\newcommand{\bayered}{\mathbf{Red}}
\newcommand{\crossentropy}{H}
\newcommand{\entropyrate}{\hat{H}}
\newcommand{\KLrate}{\hat{D}_{\mathsf{KL}}}

\newcommand{\knowledge}{v^*}
\newcommand{\Gknw}{G_{\text{knw}}}
\newcommand{\phiI}{\boldsymbol{\phi_{\text{in}}}}

\newcommand{\source}{\phi}

\newcommand{\maxlength}{l_{\max}}

\newcommand{\codelength}{L}

\newcommand{\llm}{\boldsymbol{M}}
\newcommand{\Model}{\boldsymbol{M}}

%\newcommand{\dataModel}{\boldsymbol{\Phi}}
\newcommand{\dataModel}{\phi_{\text{data}}}
\newcommand{\knwModel}{\phi_{\text{knw}}}
\newcommand{\knwModelNew}{\phi_{\text{nknw}}}
\newcommand{\synModel}{\phi_{\text{syn}}}
\newcommand{\instructionModel}{\phi_{\text{ins}}}
\newcommand{\knwfamily}{\mathbf{\Phi_{\text{knw}}}}
\newcommand{\synfamily}{\mathbf{\Phi_{\text{syn}}}}
\newcommand{\datafamily}{\mathbf{\Phi_{\text{data}}}}
\newcommand{\knwprobability}{\mathcal{P}_{\boldsymbol{\Phi_{\text{knw}}}}}
\newcommand{\synprobability}{\mathcal{P}_{\boldsymbol{\Phi_{\text{syn}}}}}
\newcommand{\dataprobability}{\mathcal{P}_{\boldsymbol{\Phi_{\text{data}}}}}
\newcommand{\knwcluster}{\phi_{\text{knw}}}
\newcommand{\knwdimension}{d_{\text{knw}}}
\newcommand{\syndimension}{d_{\text{syn}}}
\newcommand{\knwLip}{L_{\text{knw}}}
\newcommand{\synLip}{L_{\text{syn}}}
\newcommand{\syncover}{n_{\text{syn}}}
\newcommand{\knwcover}{n_{\text{knw}}}
\newcommand{\knwset}{\mathbb{K}}
\newcommand{\knw}{\mathsf{k}}
\newcommand{\syn}{\mathsf{s}}

\newcommand{\knwPredictor}{\boldsymbol{\Theta}_{\text{k}}}
\newcommand{\synPredictor}{\boldsymbol{\Theta}_{\text{s}}}

\newcommand{\PYP}{\mathsf{PYP}}
\newcommand{\PYMM}{\mathsf{PYMM}}
\newcommand{\PYCRP}{\mathsf{PYCRP}}
\newcommand{\basemeasure}{\pi_{\text{knw}}}
\newcommand{\synprior}{\pi_{\text{syn}}}
\newcommand{\question}{\psi}
\newcommand{\questionset}{\Psi}
\newcommand{\answer}{\omega}
\newcommand{\ratedistortion}{\mathcal{D}}
\newcommand{\distortionrate}{\mathcal{R}}

\newcommand{\Len}{\mathsf{Len}}

\newcommand{\testerror}{\mathsf{Red}}
\newcommand{\ferror}{\mathsf{Red}}
\newcommand{\dataerror}{\mathsf{Red}_D}
\newcommand{\modelerror}{\mathsf{Red}_M}
\newcommand{\testdataerror}{\mathsf{Red}_D}
\newcommand{\testmodelerror}{\mathsf{Red}_M}

\newcommand{\Knwelem}{\boldsymbol{\kappa}}
\newcommand{\Synelem}{\boldsymbol{\xi}}

\DeclareMathOperator*{\E}{\mathbb{E}}


%%%%%%%%%%%%%%%%%%%%%%%%%%%%%%%%%


\newcommand{\jian}[1]{{\it \color{red}{Jian: #1}}}
\newcommand{\zhixuan}[1]{{\it \color{blue}{zhixuan: #1}}}
\newcommand{\sw}[1]{{\it \color{orange}{shaowen: #1}}}

\newcommand{\eat}[1]{}


% \input{math_commands}
% \input{symbols}

%\title{Understanding LLM Behaviors via Kolmogorov and Shannon Theories: Compression, Knowledge Acquisition, and Scaling Laws}

\title{Understanding LLM Behaviors via Compression: 
Data Generation, Knowledge Acquisition and Scaling Laws}

%\title{Understanding LLM Behaviors: Compression, 
%Data Generation, Scaling Laws}


%\title{Understanding LLM Behaviors
%via Kolmogorov and Shannon's Information Theory
%}

\author{%
  Zhixuan Pan$^{1,}$\thanks{Equal contribution.} \quad Shaowen Wang$^{1,}$\footnotemark[1] \quad Pengfei Liao$^{2}$ \quad Jian Li$^{1,}$\thanks{Corresponding author.} \\
  $^{1}$Institute for Interdisciplinary Information Sciences, Tsinghua University \\
  $^{2}$School of Computer Science and Engineering, Beihang University \\
  \texttt{panzx24@mails.tsinghua.edu.cn}, \texttt{wangsw23@mails.tsinghua.edu.cn} \\ \texttt{liaopf22@buaa.edu.cn}, \texttt{lapordge@gmail.com}
}

\newcommand{\fix}{\marginpar{FIX}}
\newcommand{\new}{\marginpar{NEW}}

\begin{document}


% \begin{document}

% \date{}
\maketitle


\begin{abstract}
Large Language Models (LLMs) have demonstrated remarkable capabilities across numerous tasks, yet principled explanations for their underlying mechanisms and several phenomena, such as scaling laws, hallucinations, and related behaviors, remain elusive. In this work, we revisit the classical relationship between compression and prediction, grounded in Kolmogorov complexity and Shannon information theory, to provide deeper insights into LLM behaviors. By leveraging the Kolmogorov Structure Function and interpreting LLM compression as a two-part coding process, we offer a detailed view of how LLMs acquire and store information across increasing model and data scales -- from pervasive syntactic patterns to progressively rarer knowledge elements. Motivated by this theoretical perspective and natural assumptions inspired by Heap’s and Zipf’s laws, we introduce a simplified yet representative hierarchical data-generation framework called the Syntax-Knowledge model. Under the Bayesian setting, we show that prediction and compression within this model naturally lead to diverse learning and scaling behaviors observed in LLMs. In particular, our theoretical analysis offers intuitive and principled explanations for both data and model scaling laws, the dynamics of knowledge acquisition during training and fine-tuning, factual knowledge hallucinations in LLMs. The experimental results validate our theoretical predictions.
\end{abstract}
\newpage
\begin{spacing}{0.9} 
    \tableofcontents
\end{spacing}
\newpage



\section{Introduction}

% \begin{arxiv}
% \textit{This is arxiv.}
% \end{arxiv}
% \begin{nips}
% \textit{This is nips. }
% \end{nips}

Large Language Models (LLMs) have emerged as one of the most influential breakthroughs in modern artificial intelligence, achieving impressive performance across a multitude of tasks, ranging from fluent text generation, translations, and summarization to answering factual queries, and even performing complex reasoning. Despite these remarkable achievements, a theoretical understanding of what enables LLMs to generalize so effectively remains limited. Traditional learning theory frameworks have not yet fully explained why certain scaling laws hold, why in-context learning emerges, or when and why hallucinations arise in the output of these models.

One promising lens for gaining deeper insight into LLM behavior is the intrinsic connection between prediction and compression. Kolmogorov complexity and Shannon information theory have long established that optimal prediction of a data sequence is intimately tied to the most efficient compression of that sequence (see e.g.,
\cite{cover2006elements,li2008introduction, rissanen1984universal, hutter2005universal}).
From this perspective, a predictive model, particularly a large language model (LLM), can be viewed as a practical approximation of the Kolmogorov compressor (see e.g., \citep{Ilyatalk2023,deletanglanguage,li2025lossless}) of training data. 
See \Cref{sec:kolmogorov} and \Cref{app:prediction_compression} for more details.

\eat{
Following the above idea, we go one step further.
We draw inspiration from the {\em Kolmogorov structure function} \citep{li2008introduction,rissanen2005kolmogorov} and we view LLM training as constructing a {\em two-part (data-to-model) code} for the training data. In the first part (the model part, which corresponds to the LLM), the model encodes regularities or structures by adjusting its parameters to efficiently compress frequently recurring syntactic patterns, followed by relatively common knowledge, and subsequently progressing to increasingly rare knowledge elements. In the second part (the data part, which corresponds to the compressed code of the data, using LLM as the compressor
\footnote{
See \citep{rissanen1976generalized,cover2006elements} or \citep{deletanglanguage}
for the details of using an LLM (or any auto-regressive predictor) to compress
data using Arithmetic code.
}
), the residual from the model’s predictions, such as factual knowledge beyond the model’s capacity, or unpredictable intrinsic noise, are explicitly encoded. Consequently, when viewed as sophisticated compression engines, LLMs systematically learn to acquire and store more `compressible" patterns within their capacity constraints, explicitly encoding rarer information and unpredictable noise separately.
}

Following this line of thought, we interpret LLM training as constructing a {\em two-part (data-to-model) code} for the training data. The first part (the model compressor part) corresponds to the LLM itself, which adjusts its parameters to learn patterns and structural regularities for more efficient compression. The second part (the data part) is the compressed code of the data, generated by using the LLM as the compressor.
\footnote{For details on using an LLM or any auto-regressive predictor to compress data via Arithmetic coding, see \citep{rissanen1976generalized,cover2006elements,deletanglanguage}
or Appendix~\ref{app:llmcompressor}.
} 
An important concept for quantifying the compression capability of a model or compressor is the Kolmogorov structure function, originally introduced by Kolmogorov (see, e.g., \citep{li2008introduction,rissanen2005kolmogorov}). The structure function characterizes the compressibility of data as a function of the model’s capacity, providing a formal framework for analyzing how model complexity influences achievable compression.
As already indicated by prior work on the structure function (see e.g., \citep{vereshchagin2004kolmogorov,rissanen2005kolmogorov,lee2022relaxing}), a model of low complexity captures only the most prominent regularities in the data, whereas allowing model of higher complexity captures more nuanced structures. In the context of compressing language corpus under capacity constraints, the most efficient LLM-based compressor should initially focuses on compressing frequently recurring regularities such as syntactic patterns, then integrates relatively common knowledge, and eventually handles increasingly rare knowledge elements. Any “residual” (such as factual knowledge exceeding the model’s capacity or unpredictable noise) must be left out of the model and explicitly encoded in the second part. 
%Consequently, when viewed as sophisticated compression engines, LLMs learn to store “compressible” patterns within their capacity constraints, while leaving rarer information and unpredictable noise out of the model. 
See \Cref{sec:kolmogorov} for formal definition
of Kolmogorov structure function
and Figure~\ref{fig:kol}(a) for a schematic illustration.

Motivated by the insight from the study of Kolmogorov Structure Function and natural assumptions inspired by {\em Heap’s law} \citep{heaps1978information} 
and {\em Zipf’s law} \citep{zipf2013psycho,zipf2016human}
(see \Cref{sec:relatedwork} for more details), 
we propose the Syntax-Knowledge model, a (simplified) hierarchical data-generative framework that decomposes the generation process into two components (see Figure~\ref{fig:kol}(b)). The first component, a {\em parametric Syntax Model}, captures the syntactic structures of language, allowing for random syntactic variations. The second component, a {\em  Knowledge Model}, encodes relevant world knowledge and is modeled using the nonparametric {\em Pitman-Yor Chinese Restaurant Process}. This choice reflects the growing nature of human knowledge and captures the fact that certain pieces of information occur disproportionately more frequently than others in the data. By examining how efficiently data generated by this model can be compressed (specifically, by minimizing perplexity, or equivalently coding redundancy), we clarify the learning behaviors of LLMs: highlighting that syntax model is learned first at a faster rate, and the (factual) knowledge elements are acquired according to the order of their frequency. Furthermore, we theoretically demonstrate how model performance scales with both data size and model capacity, thus providing intuitive explanations for the scaling laws observed in real-world LLM training.
Our contributions are summarized as follows:




\begin{enumerate}
    \item \textbf{Kolmogorov Structure Function Perspective.}
   By viewing LLM training through the lens of compression, in particular a two-part coding process in the Kolmogorov Structure Function, we propose a principled framework that interprets the scaling laws of LLMs as compression under varying model capacities. In particular, we observe a natural connection between the Kolmogorov Structure Function and model scaling laws in LLMs. This perspective enables the use of theoretical tools (such as universal coding, see \cref{subsec:universalcoding_codinggame} for a brief introduction) and insights from Kolmogorov complexity and Shannon information theory, and provides an theoretical basis for distinguishing between structural regularities and incompressible randomness of language. See \Cref{sec:kolmogorov}.

    \item \textbf{A Non-Parametric Hierarchical Data Generation Model.}
    In Section~\ref{sec:datageneration}, we introduce the \textbf{Syntax-Knowledge model}, a generative model of language. The model 
    is a hierarchical Bayesian mixture model, that separates language syntax (captured by a parametric model) from (factual) knowledge (represented by a nonparametric \emph{Pitman-Yor Chinese Restaurant Process}). This design naturally accommodates the growing nature and  power-law distributions of knowledge elements, motivated by Heap’s law \citep{heaps1978information} and Zipf’s laws 
    \citep{zipf2016human}, reflecting their disproportionate frequency in real-world data.

    \item \textbf{Data Scaling Law.}
    Within a Bayesian coding framework (defined in the universal coding context), we show that perplexity minimization applied to data generated by our Syntax-Knowledge model naturally leads to data scaling laws observed in real world LLMs. In particular, one can bound Bayesian redundancy of the optimal compressor/predictor 
    using the mutual information between the prior and
    the data, and provide an upper bound of data scaling law
    by    
    $\widetilde{O}\left(C_{\text{knw}}/N^{1-\alpha}+ C_{\text{syn}}/N\right)+H$, 
    (see \cref{cor:datascalinglaw} for details)
    where $H$ is the entropy of the model (incompressible part of the loss), $N$ is the size of the training data, $C_{\text{knw}}$ and $C_{\text{syn}}$
    are constants depending on the knowledge model and syntax model respectively, and $\alpha$ is the discount parameter of the Pitman-Yor Chinese Restaurant Process employed in our knowledge model.
    See Section~\ref{sec:baye data law}. 
    
    \item \textbf{Model Scaling Law.} 
    In Section~\ref{sec:model_scaling_law}, we extend our theoretical models to account for model scaling laws, under a slightly different set of assumptions. As a consequence, we offer a more fine-grained understanding of model scaling by decomposing the test loss according to the frequency of knowledge elements. This allows us to accurately predict which knowledge elements LLMs can acquire under capacity constraints during training, and which ones they are more likely to hallucinate(even though the model may have encountered them many times) as demonstrated in Figure~\ref{fig:accuracy_freq_modelsize}.
    See Figure~\ref{fig:loss_modelscaling_decomposition} 
    for both our theoretical prediction and the corresponding experimental results on model scaling behaviors. 
    \footnote{
    We follow the experimental setting of \cite{AllenZhu-icml2024-tutorial,allen2024physics}, but use a power law distribution of individuals. The experimental details can be found in Appendix~\ref{app:experiment}.
    }
    \begin{arxiv}
    \item \textbf{Fine-Tuning.}
    In Section~\ref{sec:instruction}, we provide an intuitive explanation of the learning behaviors during fine-tuning (for instruction-following or/and knowledge injection) using the theoretical insight from our Syntax-Knowledge framework. In particular, minimizing fine-tuning loss on instruction-specific datasets primarily leads the model to learn {\em new syntax model} first while retaining previously learned knowledge. 
    Our theoretical model also provides practical recommendations that are in line with the established best practices in the literature.
    \end{arxiv}
    
     %Moreover, our model suggest that injecting new knowledge may be less effective and lead to more severe forgetting when the syntax format is drastically different from that of pretraining, especially when the pretrained model is more capacity constrained.
\end{enumerate}


\begin{figure}[t]
    \centering
    \begin{minipage}[b]{0.45\textwidth}
        \centering
        \includegraphics[width=\textwidth,scale=0.95]{loss_decomposition_fixed.jpg}
%        \caption*{(a)}
    \end{minipage}
    \hfill
    \begin{minipage}[b]{0.45\textwidth}
        \centering
        \includegraphics[width=\textwidth,scale=0.95]{theory_loss.jpg}
%        \caption*{(b)}
    \end{minipage}
    \caption{
    Decomposition of validation loss on knowledge tokens by frequency class, as model size increases.
    (a) Empirical results on a power-law-distributed dataset: knowledge tokens are grouped into four frequency classes (from most to least frequent) and colored accordingly. We observe the following trend: smaller models capture only the most frequent knowledge (the loss of the most frequent class decreases the first), while larger models gradually acquire less frequent knowledge. Each vertical dashed line marks the model size beyond which further loss reduction for a given frequency class becomes negligible, indicating the irreducible part of the loss. (b) Theoretical prediction of the same loss decomposition 
    (the optimal solution of the constrained optimization problem~\eqref{eq:constrainedoptimization}
    in Section~\ref{sec:model_scaling_law}) with irreducible loss part
    (i.e., the $H(P_\phi)$ term in \eqref{eq:reddef1}), which reproduces this frequency-dependent acquisition order and plateauing behavior. 
    }  
    \label{fig:loss_modelscaling_decomposition}
    % \vspace{-0.4cm}
\end{figure}


%By grounding LLM behavior in the frameworks of Kolmogorov complexity and Shannon information theory, our work provides a theoretical lens for understanding how these models learn, store, and retrieve knowledge. 



\section{Background}
\label{sec:background}

In this section, we review the necessary background on information theory, coding, perplexity minimization, and scaling laws in LLMs. 
Let $\tokenset$ denote the set of all possible tokens. 
We denote the training corpus as $X_{1:N} := X_1, X_2, \ldots, X_N$, where each $X_i = x_1^{(i)} x_2^{(i)} \cdots x_{l_i}^{(i)}$ is a sentence represented as a sequence of tokens, with each token $x_j^{(i)} \in \tokenset$. 
% We further assume that $l_i \leq \maxlength$ for all $i$, meaning that all sentences have a maximum length of $\maxlength$.




\paragraph{Perplexity Minimization in LLMs.}
In LLMs, the cross-entropy loss (or log-loss) serves as the metric for 
measuring how well the model predicts a given sequence of tokens. 
A model \( \Model \), which induces a predictive distribution \( P_{\Model} \), estimates the conditional probability of each token \( x_t \) given the preceding context \( x_{1:t-1} \) for each sentence.
% which is a sequence of at most $\maxlength$ tokens.
We assume that all sentences in the training corpus $X_{1:N}$
are independently drawn from the source distribution $P_{\dataModel}$.
Given the training corpus $X_{1:N}$, the training objective is to minimize the \emph{empirical averaged cross-entropy loss}:
\begin{align}
\label{eq:perplexity_minimizatio}
    \mathcal{L}(\Model) 
     &   = -\frac{1}{N}\log P_{\Model}(X_{1:N}) =  - \frac{1}{N}
    \sum_{i=1}^N \log \, P_{\Model} (X_i) 
    = \mathbb{E}_{X\sim \widehat{P}_\phi}\left[-\log P_{\Model}(X)\right] \notag
    \\    & = \crossentropy(\widehat{P}_{\source} \| P_{\Model})
    =
    - \frac{1}{N}\sum_{i=1}^N \sum_{t}
    \log \, P_{\Model}(x_t^{(i)} \mid x^{(i)}_{1 : t-1}),
\end{align}
where $\widehat{P}_{\source}$ is the empirical measure of 
the underlying source distribution $P_{\dataModel}$,
and $\crossentropy(P\|Q) = \E_{X\sim P}[-\log Q(X)]$ 
is the standard cross-entropy loss.
The \emph{perplexity} (PPL) is then defined as 
\(\mathrm{PPL} := \exp\bigl(\mathcal{L}(\Model)\bigr)\), capturing the 
effective "number of choices" the model has at each token. 
%Minimizing perplexity thus encourages the LLM to align its predictions more 
%closely with the true distribution of the training data.

\paragraph{Lossless Compression and Redundancy.}
The goal of lossless compression is to encode a sequence of tokens 
\(X=\displaystyle x_{1:n}\) (sampled from the source distribution $P_{\phi}$ parametrized by
\(\displaystyle \phi\)) into a binary sequence \(\displaystyle y_{1:m}\) of 
minimal \emph{expected} length, such that \(\displaystyle x_{1:n}\) can be 
recovered perfectly from \(\displaystyle y_{1:m}\). We use a binary source 
code \(\displaystyle c : \tokenset^* \to \{0, 1\}^*\), 
which encodes each possible 
sequence \(\displaystyle x_{1:n}\) into a binary codeword 
\(\displaystyle c(x_{1:n})\) of length \(\displaystyle \ell\bigl(c(x_{1:n})\bigr)\) bits. 
The objective is to minimize the expected code length
\begin{align}
\label{eq:codelength_crossentropy}
    \codelength(Q_c) \;:=\; 
    \mathbb{E}_{X \sim P_{\phi}}\bigl[\ell\bigl(c(X)\bigr)\bigr]
    \;=\; 
    \mathbb{E}_{X \sim P_{\phi}}\bigl[-\log Q_c(X)\bigr]
    \;=\; 
    \crossentropy(P_{\phi}\,\|\, Q_c).
\end{align}
where $Q_c(x) \propto 2^{-\ell(c(x))}$ is the {\em predictive probability corresponding to the code} $c$.
% ($Q_c$ is a valid probability distribution by Kraft inequality \citep{cover2006elements}).
According to Shannon's source coding theorem~\cite{Shannon1948}, the average length of any lossless code is bounded below by the entropy \(H(P_\phi)\), and one can design an encoding scheme whose code length is \(H(P_\phi) + O(1)\) if $P_\phi$ is known. However, when the source distribution \(P_\phi\) is unknown, universal coding becomes necessary (see Appendix~\ref{subsec:universalcoding_codinggame} for more details). In this setting, the extra code length beyond \(H(P_\phi)\), referred to as the {\em redundancy} of the code \(c\), is defined as follows:
\begin{align}
\label{eq:reddef1}
    \redundency(Q_c, P_\phi) 
    = L(Q_c)-H(P_\phi) 
     = \crossentropy(P_\phi\,\|\,Q_c)-H(P_\phi) 
    = \DKL(P_\phi\,\|\,Q_c).
\end{align}
where $\DKL(P_{\phi}\,\|\,Q_c)
=\mathbb{E}_{x\sim P_{\phi}}\left[\log (P_{\phi}(x)/Q_c(x)) \right]$ is the Kullback–Leibler (KL) divergence between source distribution $P_{\phi}$ and predicted distribution $Q_c$.
Comparing equations \eqref{eq:perplexity_minimizatio} and \eqref{eq:codelength_crossentropy} reveals that obtaining a predictive model that minimizes perplexity (or cross-entropy loss) is essentially equivalent to finding a code with minimal expected length. (See Appendix~\ref{app:llmcompressor} for further details on the {\em equivalence between prediction and compression}.) Moreover, one can see from  \eqref{eq:reddef1} that \(H(P_\phi)\) constitutes the irreducible part of the loss, corresponding to the minimum achievable code length under Shannon’s source coding theorem. Therefore, the goal of minimizing the expected code length is essentially equivalent to minimizing the redundancy.


\paragraph{Scaling Laws in LLMs.}
The performance of LLMs, particularly the cross-entropy loss $\mathcal{L}$, has been observed to improve predictably with increases in model size, dataset size, and computational resources\citep{kaplan2020scaling,hoffmann2022training}. These empirical relationships are known as \emph{scaling laws}. A common formulation for the loss as a function of dataset size $D$ (e.g., number of tokens) and model size $M$ (e.g., number of parameters) $\mathcal{L}(D, M) \approx (D/D_0)^{-\alpha} + (M/M0)^{-\beta} + E$.
In these formulations, $D_0$ and $M_0$ represent characteristic scales for the dataset and model size respectively. The exponents $\alpha > 0$ and $\beta > 0$ determine how quickly the loss decreases as the dataset size and model size increase. The term $E$ represents the \emph{irreducible loss}, which is the minimum achievable loss that cannot be reduced by further scaling, potentially due to factors like the inherent entropy of the data.



\begin{figure}[t]
  \centering
  \begin{minipage}[t]{0.48\linewidth}
    \centering
    \includegraphics[width=\linewidth,scale = 0.95]{accuracy_frequency.jpg}
  \end{minipage}
  \hfill
  \begin{minipage}[t]{0.48\linewidth}
    \centering
    \includegraphics[width=\linewidth,scale = 0.95]{accuracy_model_size.jpg}
  \end{minipage}
  \caption{
    (a) Accuracy of sufficiently trained models with different sizes across varying input frequencies. When the frequency falls below a model-specific threshold, small models inevitably hallucinate and fail to learn the corresponding facts. 
    (b) Accuracy of different frequency classes (split into four quantiles) under varying model sizes. As model size increases, the model progressively learns the more frequent data first, while infrequent data becomes learnable only at larger scales.
  }
  \vspace{-0.4cm}
  \label{fig:accuracy_freq_modelsize}
\end{figure}

\section{Kolmogorov Structure Function and Scaling Laws}
\label{sec:kolmogorov}
In this section, we briefly introduce
the concepts of Kolmogorov complexity
$K(\cdot)$ and Kolmogorov structure function
$h_X(\cdot)$
(see the classic book \cite{li2008introduction} for more details),
and how these concepts are connected to LLMs.
Although $K(\cdot)$ and $h_X(\cdot)$
are not computable, their properties directly 
motivate our theoretical modeling in latter sections.

\noindent
{\bf Kolmogorov complexity:}
The Kolmogorov complexity $K(X)$ of a string $X$ is defined as the length of the shortest binary program $p$ that outputs $X$ when run on a universal Turing machine $U$. Formally:
$
K_U(X) := \min\nolimits_{p} \{|p| : U(p) = X\}
$
where $U$ is a fixed universal Turing machine, $p$ is a binary program and $|p|$ denotes the length of the program $p$ in bits.
When there is no confusion, we omit the subscript $U$.
The Kolmogorov complexity $K(X)$ measures the information content of $X$ and is sometimes referred to as the {\em algorithmic complexity} of $X$.
$K(X)$ has been considered by many as {\em the ``ultimate" notion for the compression } 
of $X$
\footnote{
Kolmogorov complexity is {\em universal} in the following sense:
If $U$ is a universal TM, for any other TM $A$, there is a 
constant $c_A$ such that $K_U(X)\leq K_A(X)+c_A$ for all string $X$.
See \cite{li2008introduction}.
}
and modern LLMs can be thought as approximations of the Kolmogorov compressor (e.g., Ilya Sutskever's talk \citep{Ilyatalk2023}).

\noindent
{\bf Two-part description:}
One can describe the data $X$ by a two-part description:
the model description, and the {\em data-to-model} code describing
$X$ using the model. In particular, 
for any given lossless compressor $\llm$, one can see that
\begin{equation}
    K(X)\le K(\llm)+\codelength_{\llm}(X)+O(1), \label{universal optimal}
\end{equation}
where $\codelength_{\llm}(X)$ represents the compressed length of $X$ when using a lossless compressor $\llm$.  
When viewing an LLM as the compressor, 
the two-part description of the data $X$
consists of $K(\llm)$, the description of the LLM (the architecture and the parameters), and the data-to-model code of length
$\codelength_{\llm}(X)$.
Specifically, for an LLM that predicts the probability
$P_{\llm}(x_n\mid x_{1:n-1})$ of the
next token auto-regressively, 
the code length $\codelength_{\llm}(X)$ can be 
bounded by $-\log P_{\llm}(X)+O(1)$ since  
one can encode $X$ token by token using this many bits (via Arithmetic coding \citep{rissanen1976generalized,deletanglanguage}).
See Appendix~\ref{app:llmcompressor} for  details.
 

%As $|x|\rightarrow \infty$, the terms $K(C)$ and $O(1)$ on the right-hand side of the inequality \eqref{eq: universal optimal} become negligible compared to $C(x)$. Thus Kolmogorov compressor is universally nearly optimal.

\noindent
{\bf Kolmogorov structure function:}
To better understand \eqref{universal optimal},
the structure of the data $X$,
and the interplay between the model size $K(\llm)$ and code length 
$\codelength_{\llm}(X)$,
we need the concept of {\em Kolmogorov structure function}, also introduced by \cite{Kolmogorov1974}.
\footnote{
See \cite{vereshchagin2004kolmogorov} for a detailed historic account.
}
It describes how well data $X$ can be “explained” or “compressed” by models of varying complexity.

\begin{definition}[\textbf{Kolmogorov structure function}]
In this paper, we need the probabilistic version of 
Kolmogorov structure function.
(see \cite{grunwald2003kolmogorov}).
If \(X\) is genuinely “typical” under certain probabilistic model 
\(P_{\Model}\) (parametrized by $\Model$), 
\(-\log P_{\Model}(X)\) represents the code length
of $X$ under model $\Model$, ignoring the integer constraint.
The Kolmogorov structure function is defined as
$$
h_X(\alpha) = \min\nolimits_{\Model}
\{
-\log P_{\Model}(X) : P_{\Model}(X) >0, K(P_{\Model})\leq \alpha.
\}
$$
\end{definition}

\noindent
Recall from Eq.~\eqref{eq:perplexity_minimizatio}
that $-\log P_{\Model}(X)$ is the empirical cross-entropy loss, which is also code length of data $X$ using predictive model $M$(up to an additive constant).
See Figure~\ref{fig:kol}
for a schematic diagram of Kolmogorov structure function
in the context of LLMs (or any probabilistic language models).
See Appendix~\ref{app:kolstructurefunc} for more details
about the set definition of Kolmogorov structure functions, 
and related concepts
including {\em random deficiency}, {\em minimum sufficient statistics}.
We explain several key aspects that motivate
our theoretical modeling in this paper.

\begin{enumerate}
\item 
\noindent
{\bf  Kolmogorov Structure Function and Scaling Laws}:  
The Kolmogorov Structure function measures the performance of an optimal compressor under a given model-size constraint. It provides deep insights into the structure of the data from the compression perspective as the complexity of compressor increases. 
\uline{Such relationship directly mirrors the {\em model scaling laws} observed in the context of LLMs} \citep{kaplan2020scaling,hoffmann2022training}, \uline{and $h_X(\alpha)$ can be viewed a {\em theoretical lower bound} 
of the empirical scaling laws (for all possible LLM architectures).}


\item 
{\bf Structure of the Data:}  
From the perspective of the Kolmogorov structure function \(h_X(\alpha)\), 
\uline{the structure of the data $X$ (its underlying regularities and randomness)
is progressively revealed as the model size (or complexity) $\alpha$ increases, from easily compressible patterns (syntax, widely used common knowledge) to increasingly incompressible components (rare knowledge, randomness).}
More concretely, with relatively small 
\(\alpha\), a small model can captures the most pervasive regularities, such as syntax in a linguistic dataset, because such patterns recur across every sentences (and therefore yield the greatest immediate reduction in coding cost or cross-entropy loss). As \(\alpha\) grows, the model gains enough capacity to encode less frequent regularities, including widely shared factual knowledge that appears relatively frequently in the data, albeit less frequently than the universal syntax. Next, it captures rarer forms of knowledge that appear in smaller subsets of the data. Finally, for very large \(\alpha\) exceeding the {\em minimal sufficient statistics}, the model starts to encode  “random noise".
We note that similar high-level ideas appeared in the study of
the shape of the structure function
\citep{vereshchagin2004kolmogorov,rissanen2005kolmogorov,li2008introduction}.


\item 
{\bf Minimal Sufficient Statistics and Randomness: }
The \emph{irreducible test loss} (e.g., the entropy term in \eqref{eq:reddef1} or \eqref{redundancy}) 
corresponds to the lowest possible test loss, ideally achieved at the \emph{minimal sufficient statistics}. Beyond this point, any further drop in \emph{training} loss merely reflects memorizing only pure randomness, such as different re-phrasings of the same fact. In our data generation model (see Section~\ref{sec:datageneration}), we adopt a probabilistic syntax model to capture such pure randomness in the language.


\item 
{\bf Redundancy and Scaling Laws:}  
A central quantity in Figure~\ref{fig:kol} is the 
{\em redundancy} of the code, given by the KL divergence \(\DKL(P_{\dataModel}\|P_{\llm}) = \crossentropy(P_{\dataModel}\|P_{\llm}) - H(P_{\dataModel})\), where \(P_{\llm}\) denotes the distribution predicted by the model and \(P_{\dataModel}\) the true data distribution. 
Hence, \uline{characterizing the redundancy curve $\DKL(P_{\dataModel}\|P_{\llm})$ (varying the description length of the model $\llm$) provides the idealized model scaling law.}


\item 
{\bf Growing Knowledge Models:} 
    Imagine that the data (consisting of sentences) has been generated sequentially. As the data size increases, it is natural to assume that the minimal sufficient statistics also grows, reflecting the expansion of “world knowledge” (e.g., newly discovered species, proteins, facts etc.). 
    % See Figure~\ref{fig:kol_3d} in the appendix for a three‑dimensional diagram (incorporating the $K(X)$ axis) to illustrate the shape of $h_X(\alpha)$ with increasing data sizes.
    This insight motivates our subsequent data model, in which the knowledge model grows as the data size (modeled using a {\em nonparametric} model). 
%    This indicates that continually enlarging both the dataset and the model capacity leads to further performance improvements, consistent with observed scaling laws in modern LLMs. 
\end{enumerate}


\noindent
{\bf Remarks:}
Due to the non-computable nature of $K(X)$ and $h_X(\alpha)$ the quantitative bounds presented in later sections are developed within a probabilistic framework grounded in Shannon’s information theory. Nevertheless, Kolmogorov complexity remains conceptually essential to our approach for two key reasons.
(1) Insights from Kolmogorov complexity and Solomonoff’s model
(which is based on Kolmogorov complexity) fundamentally motivate our generative data model.
(2) Kolmogorov complexity and Shannon information theory are deeply interconnected: results in one domain often translate to the other under appropriate data assumptions. For example, for a sequence $X$ typical with respect to a source $P_{\phi}$, the Kolmogorov complexity $K(X)$ asymptotically approximates the Shannon code length $-\log P_{\phi}(X)$, implying that $\mathbb{E}_{X \sim P_{\phi}}[K(X)] \approx H(P_{\phi})$. Likewise, the Kolmogorov structure function serves as a non-probabilistic analogue of the {\em rate–distortion function}  in information theory~\citep{grunwald2003kolmogorov}
(we leverage rate–distortion function in 
deriving model scaling laws, see~\cref{appsec:model scaling}). 




% \begin{figure}[t]
%     \centering
%     % \begin{minipage}[width=0.48\textwidth]
%         \includegraphics[width=\textwidth]{Arxiv/Tikz_Figure.pdf}
%         \caption{Kolmogorov Structure Function View of LLMs:
%     The $\alpha$-axis represents model size, while the y-axis represents the loss of the model, corresponding to code length of data given the model. The anti-diagonal solid straight line ($x+y=K(X)$) represents the lower bound of the code length of all possible two-part codes. The upper solid curve represents $h_X(\alpha)$.
%     \label{fig:kol} 
%     \jian{need to draw a test loss.}}
%     % \end{minipage}
        
%     % \begin{minipage}[b]{0.48\textwidth}
%     %     \centering
%     %     \includegraphics[width=\textwidth]{Arxiv/data model.pdf}
%     %     \caption*{(b)}
%     % \end{minipage}
       
% \end{figure}
\begin{arxiv}
    \begin{figure}[t]
    \centering
    \begin{minipage}[b]{0.48\textwidth}
        \centering
        \includegraphics[width=\textwidth]{Tikz_Figure.pdf}
    \end{minipage}
    \hfill
    \begin{minipage}[b]{0.48\textwidth}
        \centering
        \includegraphics[width=\textwidth]{data_model.pdf}
    \end{minipage}
    \caption{(a) Kolmogorov Structure Function View of LLMs:
    The $\alpha$-axis represents model size, while the y-axis represents the loss of the model, corresponding to code length of data given the model. The anti-diagonal solid straight line is the sufficiency line ($x+y=K(X)$), which is the lower bound of the code length of all possible two-part codes. The upper solid curve represents $h_X(\alpha)$.
    The dashed blue curve is the test loss of some LLM. 
    % See more details about the test loss curve in Figure (b) in \Cref{fig:kol_3d}.
    % \jian{need to draw a test loss.
    % Zhixuan pls handles this.} \zhixuan{111111}
    (b) Illustration of the hierarchical 
    Syntax-Knowledge model.}
    \label{fig:kol}
    \vspace{-0.3cm}
\end{figure}
\end{arxiv}

\begin{nips}
    \begin{figure}
        \centering
        \includegraphics[width=0.5\linewidth]{Tikz_Figure.pdf}
        \caption{Kolmogorov Structure Function View of LLMs:
    The $\alpha$-axis represents model size, while the y-axis represents the loss of the model, corresponding to code length of data given the model. The anti-diagonal solid straight line is the sufficiency line ($x+y=K(X)$), which is the lower bound of the code length of all possible two-part codes. The upper solid curve represents $h_X(\alpha)$.
    The dashed blue curve is the test loss of some LLM. See more details about the test loss curve in Figure (b) in \Cref{fig:kol_3d}.}
    \label{fig:kol}
    \end{figure}
\end{nips}

% \jian{mark h in the figure. and change notations in the figure..important , do not forget}

% \zhixuan{Minimal sufficient statistic, figure}


% Apart from being convenient for the technical analysis in this work, \(\lambda_X(\alpha)\) is the celebrated two-part Minimum Description Length code length (Section \textbf{V-B}) as a function of \(\alpha\), with the model class restricted to models of code length at most \(\alpha\).

% First, we consider whether LLM can achieve good compression for the training data. Below, we approach this from the perspective of the two-part code in MDL (Minimum Description Length), which is related to the structure function in Kolmogorov's framework.

% \begin{align}
%      &K(LLM)+K(X|LLM^*) \nonumber
%      \\=& K(X,LLM) \nonumber
%      \\=& K(X)     \nonumber
%      \\ \le& K(P(\cdot|LLM))-\log P(X|LLM)  \nonumber
%      \\ \le& M -\log P(X|LLM).
%      \label{Inequality}
% \end{align}
% Here, $K(X)$ represents the Kolmogorov complexity, $LLM^*$ denotes the shortest program that can output $LLM$, $LLM(X)$ is the length of the code obtained by compressing $X$ using $LLM$, and $M$ represents the model size of $LLM$. For all computations involving Kolmogorov complexity here and in the future, we will omit the $O(1)$ term.

% How can we say that LLM is a good compressor for the training data $X$? According to the inequality above, we know that
% \[
% -\log P(X|LLM) \ge K(X) - M.
% \]
% Note that the right-hand side of the equation is a fixed value. Therefore, we can say that when equality holds in the inequality, LLM is a good compressor because it has achieved optimal compression. 
% Now, let's examine how equality can be achieved step by step.

% \subsection{Structure Function——Model Scaling Law}
% First, consider the first inequality, which holds when
% \[
% K(X) = K(P(\cdot | LLM)) - \log P(X | LLM).
% \]
% In other words, LLM is a maximum likelihood estimator, which imposes two requirements:

% 1. The model size $M$ needs to be sufficiently large, at least not smaller than the minimal sufficient statistic.

% In this context, the limitation on model size corresponds to the constraint on $K(S)$ in the Minimal Description Length (MDL) function.




% \begin{figure}
%     \centering
%     \includegraphics[width=0.95\linewidth]{icml2025/Sufficient Statistic.png}
%     \caption{Caption}
%     \label{Sufficient statastic}
% \end{figure}


% {\color{blue}Figure \ref{Sufficient statastic} v.s. Scaling Law(可以做一些相关的实验验证吗？)

% synthetic实验：

% 1.不同阶的Markov Chain（可能能有相应的理论）

% 2.不同depth或者是width的transformer

% 3.更一般地distill setting}

% 2.在model size足够大的条件下，LLM确实是一个maximum likelihood estimator。


% \begin{align*}
%     P(Y|X) = \dfrac{P(X|Y)P(Y)}{P(X)}
% \end{align*}

% \subsection{Data Scaling Law}

% \jian{figure, discuss scaling law
% and kol structure function.}

% Figure \ref{Sufficient statistic} shows the scaling law corresponding to model size. We can also consider scaling law-related issues from the perspective of dataset size.

% On this front, we are concerned with how the loss changes as the dataset size increases when the model size is infinitely large. In this case, we can consider the compression rate reaching $K(X)$, where the PPL loss corresponds to the compression rate, i.e., $K(X) / |X|$.

% The corresponding curve should have $|X|$ on the x-axis and $K(X) / |X|$ on the y-axis. If we follow this formulation, we would encounter the problem of Kolmogorov Complexity being uncomputable. We will consider the following data model for a more concrete setting.




\vspace{-0.3cm}
\section{A Hierarchical Data Generation Model}
\label{sec:datageneration}
\vspace{-0.2cm}

In this section, we introduce a hierarchical data generation framework, termed the {\em syntax–knowledge model}. The design of this model reflects several key insights derived from our analysis of the Kolmogorov structure function in \cref{sec:kolmogorov} and \cref{app:kolstructurefunc}. 
Our model is also deeply inspired by Solomonoff’s model (see \cite{Solomonoff1964a,cover2006elements}), 
which can be interpreted as a Bayesian mixture model with a countably infinite set of components—each corresponding to a distinct Turing machine. Similarly, our model adopts a Bayesian mixture model to represent knowledge; however, instead of treating each component as a Turing machine, we employ parametric models as the mixture components, avoiding computability and learnability issues.

Now we describe our syntax-knowledge model.
In this model, where each sentence in the training corpus is generated by a {\em syntax encoder} that encode a (factual) knowledge element, sampled from the {\em  knowledge model}. 
The syntax model (encoder) is parameterized by $\synModel$, the knowledge model is denoted as $\knwModel$, and the entire data model is denoted as 
$\dataModel=\{\synModel,\knwModel\}$. 

In our model, the syntax model $\synModel$
(e.g., a probabilistic CFG, an English/code grammar, or even a template-based format) does not grow with the size of the dataset and can be modeled using a finite-dimensional parameterized model $\synModel$. On the other hand, the knowledge model employs a {\em non-parametric} stochastic process to account for two empirically observed phenomena: 1) the unbounded growth of factual information as datasets grow (mirroring Heap's Law in lexical growth patterns \citep{heaps1978information}), and 2) the long-tailed frequency distribution of factual occurrences, analogous to Zipf's law in natural language \citep{zipf2013psycho,zipf2016human}. 

Motivated by the above idea, we leverage the nonparametric {\em Pitman–Yor process (PYP)} \citep{pitman1997two} for modeling the knowledge model $\knwModel$. A PYP is characterized by two real parameters, the {\em discount parameter} \(0 \leq \alpha < 1\) and the {\em concentration parameter} \(\beta > -\alpha\), and a base probability measure $\basemeasure$. 
We denote it as $\PYP(\alpha, \beta, \basemeasure)$.
A sample from the Pitman–Yor process \( \PYP(\alpha, \beta, \basemeasure) \) is a  random probability measure \( \knwModel = \sum_{i=1}^{\infty} p_i \delta_{\phi_i} \),  
which is a discrete distribution with countably infinite atoms,
where \( p = (p_1, p_2, \ldots) \) are the weights generated by the Pitman–Yor Chinese Restaurant Process ($\PYCRP$) (described below);
each atom \( \phi_i \sim \basemeasure \) is the  $i$-th cluster parameter independently drawn from the base measure \( \basemeasure \).

The weights $p = (p_1, p_2, \ldots)$ 
are generated by the Pitman–Yor Chinese Restaurant Process ($\PYCRP$), denote as \( p = (p_1, p_2, \ldots) \sim \PYCRP(\alpha, \beta) \).
$\PYCRP$ adopts a preferential attachment mechanism that 
naturally captures both the {\em sublinear scaling} of new factual discoveries and the {\em power-law distributed frequencies} of knowledge pieces (see Lemma~\ref{lem:PYCRP fre} in Appendix~\ref{subsec:PitmanYor}). 
$\PYCRP$ works as follows:
Imagining a restaurant where customers arrive sequentially, each choosing either to join an existing lively table or start their own.
The first customer sits at a new table.
Consider the \( n \)-th customer who just come to the restaurant.
Suppose \( N_k \) is the number of customers already at table \( k \), and \( K \) is the current number of occupied tables.
The \( n \)-th customer either joins an existing table \(k\) or starts a new table with the following probabilities: 
\begin{arxiv}
\[
\text{For the \(n\)-th customer:} \quad
\begin{cases}
\displaystyle \frac{N_k - \alpha}{n - 1 + \beta}, & \text{if joining an existing table \(k\),}\\[10pt]
\displaystyle \frac{\beta + \alpha K}{n - 1 + \beta}, & \text{if starting a new table.}
\end{cases}
\]
\end{arxiv}
\begin{nips}
The \(n\)-th customer joins an existing table \(k\) with probability
 $(N_k - \alpha)/(n - 1 + \beta)$,
 or starts a new table with probability
 $(\beta + \alpha K)/(n - 1 + \beta)$.
\end{nips}
%Suppose \( N_k \) is the number of customers already at table \( k \), and \( K \) is the current number of occupied tables.
%The \( n \)-th customer either joins an existing table \(k\) 
%with probability $(N_k - \alpha)/(n - 1 + \beta)$ or start a new table with probability $(\beta + \alpha K)/(n - 1 + \beta)$.
The weight 
$
p=(p_1, p_2, \dots)=\lim\limits_{n\to \infty}(N_1/n,N_2/n,\cdots)
$
are defined as the relative sizes of the tables as the number of customers \( n \to \infty \).  
More details of $\PYMM$ can be found in Appendix~\ref{sec:PYMM}.

\noindent
{\bf Hierarchical Data Model:}
In the Pitman–Yor Mixture Model, the Pitman–Yor Process serves as a prior defined over the set of mixture distributions. 
Recall that a sample from the Pitman–Yor Process is a random probability measure of the form \( \knwModel = \sum_{i=1}^{\infty} p_i \delta_{\phi_i} \), 
where \( \{p_i\} \) are the mixture weights and \( \{\phi_i\} \) are the atoms. This can be viewed as a mixture distribution in which the \( i \)-th cluster is chosen with probability \( p_i \), 
and the corresponding model is parameterized by \( \phi_i \), referred to as the \( i \)-th knowledge cluster (table).

The parameters \( \phi_i \) are drawn from a base measure \( \basemeasure\), which acts as a prior over the cluster parameters. 
We assume \( \basemeasure \) is supported on a bounded parameter space \( \knwfamily = \{ \phi \in \mathbb{R}^{\knwdimension} : \|\phi\|_2 \le 1 \} \), 
ensuring that each \( \phi_i \in \knwfamily \) for all \( i \in \mathbb{N}^+ \).

% \jian{base measure H? }

Now, it is ready to describe the Syntax-Knowledge model, which generates a sentence \( X \) according to the following hierarchical framework:
\begin{enumerate}
    \item We first independently sample the latent parameters of the knowledge and syntax models:
    \begin{arxiv}
    \[
    \knwModel \sim \PYP(\alpha, \beta, \basemeasure),\quad 
    \synModel=\{\synModel^{(1)},\synModel^{(2)},\dots,
    \synModel^{(\numsyntax)}\},\quad 
    \synModel^{(i)} \overset{\text{i.i.d.}}{\sim} \synprior(\synModel),
    \]
    \end{arxiv}
    \begin{nips}
    $
    \knwModel \sim \PYP(\alpha, \beta, \basemeasure),\quad 
    \synModel=\{\synModel^{(1)},\synModel^{(2)},\dots,
    \synModel^{(\numsyntax)}\},\quad 
    \synModel^{(i)} \overset{\text{i.i.d.}}{\sim} \synprior(\synModel),
    $
    \end{nips}
    where $\synprior$ is the prior distribution of $\synModel^{(i)}$ and supported on a bounded parameter space \( \synfamily = \{ \phi \in \mathbb{R}^{\syndimension} : \|\phi\|_2 \le 1 \} \).\, ensuring that each \( \synModel^{(i)} \in \synfamily \) for all \( 1\le i\le \numsyntax  \).
    The value \( \numsyntax \) denotes the number of distinct syntax parameter vectors, indicating that different types of knowledge should be expressed through different syntactic patterns.
    \item We generate the sentence from both 
    $\synModel$ and $\knwModel$. In fact, we first sample
    an (abstract) knowledge element $\Knwelem$
    (corresponding to a customer) from a knowledge cluster (corresponding to a table) in $\knwModel$,
    \begin{arxiv}
    \footnote{
       In our paper, the abstract knowledge element \(\Knwelem\) is not given any concrete form, since it is only relevant for the data generation process. We simply assume that a sentence can be generated from \(\Knwelem\) combined with an appropriate syntax element. There is no need to specify a concrete architecture for the syntax or the knowledge models neither, as long as they satisfy certain assumptions (e.g., Assumption~\ref{ass: mutual information}), which suffice for our theoretical results.
    }
    \end{arxiv}
    use $\Knwelem$ to determine which syntax $\synModel^{(i)}$ (e.g., which template or format) to use, and then use the syntax encoder to generate the corresponding sentence. 
    \begin{arxiv}
         See Figure~\ref{fig:kol}(b) for an illustration. More details and discussions can be found in Appendix~\ref{app:datageneration}.
    \end{arxiv}
    \begin{nips}
        See Appendix~\ref{app:datageneration} for details and \Cref{fig:data_model} for the data model schematic.
    \end{nips}
    
\end{enumerate}




\section{Explaining Scaling Laws}
\setlength{\textfloatsep}{8pt}
\label{sec:mainscalinglaws}

\subsection{A Bayesian Coding Framework}
\label{sec:Bayesian Framework}
In this section, we 
explain LLM scaling laws by
adopting the Bayesian sequential prediction framework
(also called online Bayesian coding game). This Bayesian setting has been studied extensively in information theory (especially in the context of universal coding
or universal compression), statistics and machine learning literature (see e.g., \citep{clarke1990information,clarke1994jeffreys}
and more recent expositions \citep{duchi2024lecture,jeon2024information}).
More details can be found in \Cref{sec:baye game}.

% \jian{the following part needs revision. also check the appendix
% i wrote. make the notation consistent.}

Given the data-generating distribution \( P_{\dataModel} \), the redundancy of the model \( \Model \) (recall the definition from \eqref{eq:reddef1})
with respect to samples \( X_{1:N} \) i.i.d. drawn from \( P_{\dataModel} \) is:
\begin{align}
\label{redundancy}
    \redundency_N(\Model, \dataModel) 
&= \E\limits_{X_{1:N}\sim P_{\dataModel}} \Bigl[ -\log P_{\Model}(X_{1:N}) \Bigr] 
- \E\limits_{X_{1:N}\sim P_{\dataModel}} \Bigl[ -\log P_{\dataModel}(X_{1:N}) \Bigr] \\
&= \DKL(P_{\dataModel}^N \,\Vert\, P_{\Model}), \notag
\end{align}
where the first term represents the cumulative cross-entropy of the model \( \Model \), and the second term corresponds to the irreducible entropy of the ground-truth data distribution \( \dataModel \).

% It is important to note that, although the data is assumed to be generated in an i.i.d. manner, the model itself is a sequential prediction model.

% It is important to note that, although the data is assumed to be generated in an i.i.d. manner, the model \( \Model \) is formulated as a sequential prediction model. This modeling choice is crucial; otherwise, the redundancy would reduce to a trivial lower bound:
% $\DKL\left(P_{\dataModel}^N \,\Vert\, P_{\Model}^N\right) 
% = N \DKL\left(P_{\dataModel} \,\Vert\, P_{\Model}\right) = \Omega(N)$,
% \jian{Omega?? the above sentence is a bit strange...}
% which scales linearly with the number of samples and fails to capture any sequential structure in the modeling process.






In Bayesian generative models, the data-generating parameter \( \dataModel \) is treated as a random variable drawn from a prior distribution \( \pi \). 
Let \( \datafamily \) denote the parameter space.
The \emph{Bayesian redundancy} of a model \( \Model \) is defined as the expected redundancy under the prior  \( \pi \):
\[
\redundency_N(\Model, \datafamily) := \int_{\datafamily} \pi(\dataModel)\, \redundency_N(\Model, \dataModel)\, \dd\dataModel.
\]
According to \Cref{lem:baye red} in Appendix~\ref{sec:baye game}, the {\em optimal} Bayesian redundancy (over all possible predictive models) 
is given by:
\[
\inf_{\Model} \int_{\datafamily} \pi(\dataModel) \redundency_N(\Model, \dataModel) \dd\dataModel = \int_{\datafamily} \pi(\dataModel) \DKL(P_{\dataModel}^N \Vert Q_{\pi}) \dd\dataModel = \mathbb{I}(X_{1:N};\dataModel),
\]
where \( Q_{\pi} = \int P_{\dataModel}^N \, \pi(\dataModel) \, \dd\dataModel \) denotes the marginal distribution over \( X_{1:N} \) obtained by first sampling \( \dataModel \sim \pi \), and then drawing \( X_{1:N} \overset{\text{i.i.d.}}{\sim} P_{\dataModel} \) conditioned on \( \dataModel \).





\subsection{Data Scaling Law (under the Bayesian framework)} 
\label{sec:baye data law}


Under the Bayesian sequential prediction framework, we derive an upper bound on the optimal Bayesian redundancy (which is equal to the mutual information $\mathbb{I}(X_{1:N};\dataModel)$ by Lemma~\ref{lem:baye red}) of our hierarchical data model $\dataModel$. 
%This bound is consistent with the empirical data scaling laws observed during the pre-training of LLMs.

We make the following natural assumption of the knowledge model:
if the parameters of two knowledge clusters are close,
the KL divergence between two induced distributions is small.

\begin{assumption}
\label{ass: mutual information}
% We assume that the log-likelihood induced by parameters in the parameter spaces \( \knwfamily \) and \( \synfamily \) are Lipschitz continuous. More specifically, there exist positive constants \( \knwLip \) and \( \synLip \) such that for all \( \knwModel^{(1)}, \knwModel^{(2)} \in \knwfamily \) and \( \synModel^{(1)}, \synModel^{(2)} \in \synfamily \), we have:
% \begin{align*}
%     \left|\log P_{\knwModel^{(1)}}(\Knwelem) - \log P_{\knwModel^{(2)}}(\Knwelem)\right| &\leq \knwLip \|\knwModel^{(1)}-\knwModel^{(2)}\|, \quad\text{for any knowledge element} \ \Knwelem,\\[6pt]
%     \left|\log P_{\synModel^{(1)}}(\Synelem) - \log P_{\synModel^{(2)}}(\Synelem)\right| &\leq \synLip \|\synModel^{(1)}-\synModel^{(2)}\|, \quad\text{for any syntax element} \  \Synelem.
% \end{align*}
The probability families $\synfamily$ and $\knwfamily$ satisfy the following: there exist positive constants \( \knwLip \) and \( \synLip \) such that for all \( \knwModel^{(1)}, \knwModel^{(2)} \in \knwfamily \) and \( \synModel^{(1)}, \synModel^{(2)} \in \synfamily \), we have:
\begin{align*}
    \DKL\left(P_{\knwModel^{(1)}}||P_{\knwModel^{(2)}}\right) \leq \knwLip \|\knwModel^{(1)}-\knwModel^{(2)}\|^2, \quad
    \DKL\left(P_{\synModel^{(1)}}||P_{\synModel^{(2)}}\right) \leq \synLip \|\synModel^{(1)}-\synModel^{(2)}\|^2.
\end{align*}
\end{assumption}


The constants $\knwLip$ and $\synLip$ may depend on the concrete form of the parametrization of $P_{\knwModel}$ and $P_{\synModel}$, and are typically related to the Fisher information (See \Cref{lem:KL and MSE} and \Cref{justify KL and MSE}). 
However, the particular form of the parametrization is not important for our later development.
The constants $\knwLip$ and $\synLip$ appear in logarithmic order in the upcoming \Cref{thm:data_scaling_Mutual_Information}. For the sake of clarity, we omit logarithmic terms in the statement of \Cref{thm:data_scaling_Mutual_Information}. 
% \jian{check, cite and justify.}
Under the Bayesian setting and the above assumption, we can derive the following upper bound of the optimal Bayesian redundancy.

\begin{theorem}
\label{thm:data_scaling_Mutual_Information}
Under the Bayesian sequential prediction framework and \Cref{ass: mutual information}, the averaged optimal Bayesian redundancy (per sentence) of the hierarchical data model $\dataModel$  satisfies:
    \begin{align*}
    % \mathcal{L}_{N,S}:&=\dfrac{\mathbb{I}(T_{1:N}; \synModel)}{N} = \Tilde{O}(d_SN^{-1}),\\
    \inf_{\Model}\dfrac{1}{N}\redundency_N(\Model,\datafamily)=\dfrac{1}{N}\mathbb{I}(X_{1:N}; \dataModel) = 
    \widetilde{O}\left(\frac{\knwdimension}{N^{1-\alpha}}+\frac{ n_s\syndimension}{N}\right).
    \end{align*}
where $\knwdimension$ and $\syndimension$ are the parameter dimensions of the knowledge and syntax clusters, respectively, and $n_s$ is the number of distinct syntax clusters.
    % \jian{what tilde hides? Zhixuan}
\end{theorem}

Using \eqref{eq:reddef1}, we can obtain the following decomposition
of the optimal Bayesian cross-entropy loss.

\begin{corollary}
\label{cor:datascalinglaw}
Suppose $\pi$ is the prior of $\dataModel$.
Under the same setting as Theorem~\ref{thm:data_scaling_Mutual_Information}, the averaged optimal Bayesian loss (per sentence) 
can be bounded as:
\begin{align*}
    % \mathcal{L}_{N,S}:&=\dfrac{\mathbb{I}(T_{1:N}; \synModel)}{N} = \Tilde{O}(d_SN^{-1}),\\
    %\inf_{\Model}\dfrac{1}{N}\mathbb{E}_{\dataModel\sim \pi}\mathbb{E}_{X_{1:N}\sim P_{\dataModel}} \left[ -\log P_{\Model}(X_{1:N}) \right] 
    \inf_{\Model}\frac{1}{N}
    \E_{\dataModel\sim\pi}\,
    \E_{X_{1:N}\sim P_{\dataModel}}
    \bigl[-\log P_{\Model}(X_{1:N})\bigr]
    =
    \widetilde{O}\left(\frac{\knwdimension}{N^{1-\alpha}}+\frac{ n_s\syndimension}{N}\right)+\dfrac{1}{N}H(X_{1:N}|\datafamily).
\end{align*}
where $H(X_{1:N}|\datafamily)=
\mathbb{E}_{\dataModel\sim \pi}[H(X_{1:N}\mid \dataModel)]$
is the irreducible part of the loss.
\end{corollary}
% \jian{Zhixuan: state a corollary about the test loss, reflecting the 
% irreducible part.}




We now interpret \cref{thm:data_scaling_Mutual_Information} and \cref{cor:datascalinglaw}.
The left-hand side of \cref{cor:datascalinglaw} corresponds to the cross-entropy loss during LLM training, while the bound itself can be viewed as   the upper bound of the data scaling law (for the optimal model). On the right-hand side, the entropy term represents the irreducible constant observed in empirical scaling laws, and the first two terms correspond to the Bayesian redundancy, which diminishes as the data size increases.

Moreover, the two terms correspond respectively to the syntax and knowledge models. These two models exhibit significantly different learning behaviors: the redundancy for the syntax model decreases rapidly at a rate \(\Tilde{O}(N^{-1})\), whereas the redundancy for the knowledge model decreases more slowly at a rate \(\Tilde{O}(N^{\alpha-1})\). 
These differences highlight two distinct training phases: an early stage dominated by syntax redundancy reduction, and a later stage dominated by knowledge redundancy reduction.
%(see below for the experimental validation)

Such a behavior is quite intuitive: Initially, the model primarily learns syntactic structures, since frequently recurring syntactic patterns yield substantial and immediate reductions in redundancy (and thus test cross-entropy loss). Moreover, as the syntax model can be captured by a parametric model, its redundancy decreases rapidly at the parametric learning rate \(\widetilde{O}(N^{-1})\). This is also consistent with standard Bayesian results for parametric models \citep{rissanen1986stochastic,clarke1994jeffreys}. As training progresses, the model gradually incorporates factual knowledge; however, knowledge elements with lower frequencies receive fewer training examples, resulting in a slower reduction in redundancy.

Finally, we note that a few recent studies have also 
derived the data scaling law from the power‑law data distribution 
in various stylized settings
\cite{hutter2021learning, michaud2023quantization, brill2024neural}. See \Cref{sec:relatedwork}
for more discussions.
\begin{arxiv}
Furthermore, the learning dynamic described above is similar in spirit with the phenomenon of {\em simplicity bias} in machine learning, where the model tends to learn simpler (or lower-complexity) structures first
(e.g., \citep{shah2020pitfalls,lyu2021gradient,xu2024overview,li2025feature}).    
\end{arxiv}

\begin{figure}[t]
    \centering
    \vspace{-2ex}
    \begin{minipage}[b]{0.48\textwidth}
        \centering
        \includegraphics[width=\textwidth,scale=0.8]{dataregression.png}
%        \caption*{(a)}
    \end{minipage}
    \hfill
    \begin{minipage}[b]{0.48\textwidth}
        \centering
        \includegraphics[width=\textwidth,scale=0.8]{datalaw.jpg}
%        \caption*{(b)}
    \end{minipage}
   \caption{
    (a) Validation loss as a function of training data size.  Models trained on data sampled from pretrained knowledge under various power-law distributions (i.e., \( p(i) \sim (x + b)^{\text{power}} \)) show clear power-law scaling of loss with data size, while uniform sampling does not. A more skewed data distribution leads to faster loss decay.  
    (b) Loss decomposition by data frequency class: high-frequency data is learned earlier, while lower-frequency data is acquired later during training.
    }
    \label{fig:data_scaling}
\end{figure}

\noindent
{\bf Experimental Validations:}  
We validate Theorem~\ref{thm:data_scaling_Mutual_Information} through experiments using datasets sampled from pretrained knowledge following various power-law distributions as well as a uniform distribution
(See the experimental details in \Cref{app:experiment}).  
In Figure~\ref{fig:data_scaling}(a), we observe that models trained with power-law sampled data exhibit a strong power-law relationship between validation loss and data size, accurately captured by the regression form \( \mathrm{loss} = (x/x_0)^{-\alpha} + \epsilon \). The fitted exponent $\alpha$ increases in magnitude as the data distribution becomes more skewed, indicating faster loss decay, which is consistent with our conclusion in \Cref{thm:data_scaling_Mutual_Information}. In contrast, models trained with uniformly sampled data deviate significantly from this power-law behavior.  We note our theory does not cover such uniform distribution, as
it cannot arise from our data modeling based on the PYCRP.


Figure~\ref{fig:data_scaling}(b) further decomposes validation loss by data frequency class over training steps. 
We find that during the initial training steps (from the topmost curve 
to the 2nd topmost curve), the loss reduces for data of all frequencies 
(although the loss reduction for higher frequency data is slightly larger).
This indicates the syntax learning phase.
Later, we observe that high-frequency data achieves more rapid loss reduction early in training, while medium- and low-frequency data improve later. This confirms a frequency-dependent learning dynamic: the model first captures more common patterns and then gradually incorporates rarer knowledge as training progresses, which is consistent
from the insight gained from Kolmogorov structure function
(see \Cref{sec:kolmogorov}).

\vspace{-0.15cm}



\subsection{Model Scaling Law}
\label{sec:model_scaling_law}



In this section, we derive a power-law characterization for the model scaling law. Our theoretical results here are derived under slightly simplified assumptions 
(but retain essential insights).
Specifically, we focus exclusively on the knowledge model and omit the syntax model, motivated by our earlier findings in Theorem~\ref{thm:data_scaling_Mutual_Information}, where we showed (and empirical observations confirm) that the syntax model is learned at a significantly faster rate. 

We continue modeling the knowledge model \( \knwModel = \sum_{i=1}^{\infty} p_i \delta_{\phi_i} \) as an infinite mixture model but now make the following assumption on the mixing probabilities \( p_i \). 

\begin{assumption}
\label{ass:powerlaw}
For the mixture model \( \knwModel = \sum_{i=1}^{\infty} p_i \delta_{\phi_i} \), the mixing probabilities \( p_i \) 
follow a power-law distribution:
$
p_i = \zeta(1/\alpha)^{-1} \, i^{-1/\alpha},
$
\begin{arxiv}
where \( \zeta(1/\alpha) = \sum_{i=1}^{\infty} i^{-1/\alpha} \) is the Riemann zeta function.
\end{arxiv}
\begin{nips}
where \( \zeta(1/\alpha) = \sum_{i=1}^{\infty} i^{-1/\alpha} \).
\end{nips}
\footnote{
This choice of exponent \(1/\alpha\) aligns with the asymptotic behavior of mixing weights in the Pitman–Yor Process \( \mathrm{PYP}(\alpha, \beta, \basemeasure) \) (see Lemma~\ref{lem:PYCRP fre}). 
}
\end{assumption}
In this section, we consider the optimal redundancy achievable by an omniscient model \( \Model^*_\capacity \) under the constraint $\capacity$. That is, we may construct the model \( \Model^*_\capacity \) from the true data distribution \( \dataModel \). The only constraint of 
\( \Model^*_\capacity \) is that the model capacity is at most $\capacity$ bits.
Due to the finite capacity, the model 
\( \Model^*_\capacity \) 
cannot memorize all the data nor the true distribution
\( P_{\dataModel} \), hence
must apply lossy compression to the true model.
In particular, We denote by 
$\ratedistortion_i(m_i)$
the redundancy incurred by answering questions from knowledge cluster $\phi_i$, and \( m_i \) denotes the constraint of the mutual information between
the model and $\phi_i$ (one may think it as the memory allocated to compress \( \phi_i \)).  
The minimal achievable redundancy under a given memory constraint can be characterized by a distortion-rate function, and we make the following assumption of  
$\ratedistortion_i(m_i)$.
\begin{arxiv}
(it is decrease as an exponential function of $m_i$).
\end{arxiv}
For the formal definition of 
$\ratedistortion_i(m_i)$
and the justification of 
the assumption, see Appendix~\ref{appsec:model scaling}.

\begin{assumption}
\label{ass:model scaling}
We assume that the distortion-rate function \( \ratedistortion_i(R) \) satisfies
\begin{arxiv}
    \begin{enumerate}
    \item there exists positive constant $c_3,c_4$ such that for all $i\in \mathcal{N}^+,R\le c_3$, the distortion-rate function $\ratedistortion_i(R)\ge c_4$,
    \item there exists positive constant $c_{\max},b_{\max}$ such that for all $i\in \mathcal{N}^+$, the distortion rate function $\ratedistortion_i(R) \le c_{\max} b_{\max}^{-R}.$
\end{enumerate}
\end{arxiv}
\begin{nips}
\\
    1. There exists $c_3,c_4$ such that for all $i\in \mathcal{N}^+,R\le c_3$, the distortion-rate function $\ratedistortion_i(R)\ge c_4$.
    
    2. There exists $c_{\max},b_{\max}$ such that for all $i\in \mathcal{N}^+$, the distortion rate function $\ratedistortion_i(R) \le c_{\max} b_{\max}^{-R}.$
\end{nips}

    % \item $D_i(R)$ is non-negative and monotonically non-increasing with respect to $R$,
    % \item $D_i(R)$ is convex,
% where \( c_{\min}, c_{\max}>0, b_{\min}, b_{\max}>1 \) are some constants.
\end{assumption}

The redundancy minimization problem can be formulated as the following optimization problem:
\begin{align}
\label{eq:constrainedoptimization}
    \text{minimize} \quad & \mathbb{E}_i[\ratedistortion_i(m_i)] = \sum_{i=1}^{\infty} p_i \ratedistortion_i(m_i),   \\
    \text{subject to} \quad & \mathbb{I}(\Phi_0; \Model^*_\capacity ) \le \capacity, \quad m_i = \mathbb{I}(\phi_i; \Model^*_\capacity) \ge 0 \ \text{for all } i \in \mathbb{N}^+, \nonumber
\end{align}
where \( \Phi_0 = (\phi_1, \phi_2, \ldots) \).
If we further assume that \( \phi_i \perp \Phi_{-i} \mid \Model^*_\capacity \) where \( \Phi_{-i} = (\phi_{1},\ldots, \phi_{i-1}, \phi_{i+1}, \ldots) \), i.e., \( \phi_i \) is conditionally independent of the remaining cluster parameters given the model \( \Model^*_\capacity \).
In other words, conditional on the model \( \Model^*_\capacity \), knowing \(\phi_i\) does not provide any additional information about \(\phi_{j}\) for \(j \neq i\). 
Under this natural assumption, the model capacity constraint can be simplified as:
$
    \mathbb{I}(\Phi_0;  \Model^*_\capacity ) 
    = \sum_{i=1}^{\infty} \mathbb{I}
    (\phi_i;  \Model^*_\capacity ) \leq \capacity. 
$
See Appendix~\ref{appsec:model scaling} for the derivation.
Let \( \testmodelerror(\capacity) \) denote the optimal value of the optimization problem in \eqref{eq:constrainedoptimization}, representing the minimal achievable redundancy under the model capacity constraint.


\begin{theorem}
\label{thm:model scaling} 
 Under \Cref{ass:model scaling} and \Cref{ass:powerlaw}, the optimal value of the optimization problem under the given model size constraint satisfies:
\[\testmodelerror(\capacity) = \Theta(\capacity^{-1/\alpha+1}).\] 
% \jian{say something for each class?}
% \zhixuan{we can not We cannot derive a rigorous theoretical conclusion for each class unless we assume that the rate-distortion function is exactly an exponential function}
%\jian{Zhixuan:state the loss ,including the irreducible part}
Moreover, if we further assume that $D_k(R)=a_kb_k^{-R}$ for some constant $b_{\min}\le b_k\le b_{\max},a_{\min}\le a_k\le a_{\max}$,  
the contribution of the $k$'s cluster is
\begin{arxiv}
\[
p_k\ratedistortion_k(m_k^*)=\Theta(\min\{k^{-1/\alpha},C^{-1/\alpha}\}),
\]
\end{arxiv}
\begin{nips}
$
p_k\ratedistortion_k(m_k^*)=\Theta(\min\{k^{-1/\alpha},C^{-1/\alpha}\}),
$
\end{nips}
where $m_k^*$ is the solution of the optimization problem~\eqref{eq:constrainedoptimization}.
\end{theorem}


\noindent
{\bf Experimental Validation:} 
\begin{figure}[t]
    \centering
    \begin{minipage}[b]{0.48\textwidth}
        \centering
        \includegraphics[width=\textwidth,scale=0.95]{modellaw_regression.jpg}
%        \caption*{(a)}
    \end{minipage}
    \hfill
    \begin{minipage}[b]{0.48\textwidth}
        \centering
        \includegraphics[width=\textwidth,scale=0.95]{modellaw.jpg}
%        \caption*{(b)}
    \end{minipage}
    \caption{
    (a) Validation loss as a function of model size (excluding token embedding head).  Regression analysis of model scaling: using the fitting form $\mathrm{loss} = (x/x_0)^{-\alpha} + \epsilon$, we find that when the data is generated from a power-law distribution, the loss decreases with model size following a power-law trend. In contrast, the loss does not exhibit power-law scaling with model size when the data is generated from a uniform distribution.
    (b) Validation loss for different frequency classes: larger models are able to better fit lower-frequency (rarer) data, while for a fixed model size, more frequent data is learned more effectively.}
    \label{fig:model_scaling}
    \vspace{-0.2cm}
\end{figure}
In Figure~\ref{fig:loss_modelscaling_decomposition}(a), we show the empirical decomposition of validation loss by knowledge token frequency class as model size increases. Figure~\ref{fig:loss_modelscaling_decomposition}(b) presents the corresponding theoretical prediction, derived from the optimal solution to the constrained optimization problem~\eqref{eq:constrainedoptimization}. The empirical results on power-law distributed data closely align with the theory: in both panels, high-frequency classes exhibit loss reduction with smaller models, while low-frequency classes only begin to improve as model capacity increases.

In Figure~\ref{fig:model_scaling}(a), we present the model scaling law across different data distributions. Consistent with our theoretical predictions, we observe that the more skewed the data distribution (i.e., the heavier the power-law tail), the larger the exponent in the fitted scaling law, indicating faster loss decay as model size increases. In contrast, 
{\em for data generated from a uniform distribution, the loss curve significantly deviates from a power law}: the loss remains nearly flat until the model reaches a critical capacity threshold, after which it drops sharply, and then plateaus again. 
This pattern arises for the following reason: initially, none of the individual items stand out and are fully learned, resulting in a flat loss. As the model grows, it begins to capture a subset of properties for most individuals, causing a rapid drop. Once all such properties have been learned, the loss flattens out again. A more detailed analysis of the learning dynamics for each property is presented in Appendix~\ref{para:data_hetero}. 
We also note that \citet{allen2024physics} investigated factual knowledge acquisition using the bioS dataset consisting of uniformly distributed individuals, and one of their main findings is that the amount of acquired knowledge (in bits) scales linearly with the model size. However, they did not examine the scaling behavior in terms of cross-entropy loss, which follows a very different scaling than power law (Figure~\ref{fig:model_scaling}(a)).

Comparing the learning curves of uniformly distributed data with those of power-law-distributed data (across both data scaling and model scaling) suggests that {\em it can be advantageous to have power-law-distributed data}, because the model can gradually learn knowledge in the order of frequency, which is more effective than the uniform case, where no one element stands out and the model lacks guidance on what to prioritize. 
Exploring how adjusting the frequency or mixing ratio of data can enhance (or accelerate) learning performance is an intriguing direction for future research (see \citep{gu2025data} for a recent study in this direction).

Our theorem also implies that for a fixed model capacity $\capacity$, if a knowledge element appears with a frequency below a certain threshold, the model will choose not to learn it, despite that the model may have seen it many times during training. This aligns with our experimental findings: hallucination tends to occur when the total amount of knowledge exceeds the model's capacity, specifically, when higher-frequency knowledge already saturates the model's capacity, lower-frequency elements are ignored. As shown in Figure~\ref{fig:accuracy_freq_modelsize}, for a 7.2M-sized model, knowledge that occurs fewer than 508 times is consistently hallucinated, regardless of the number of pretraining epochs.

\begin{arxiv}
    \subsection{Real-world Data Validation}
\label{sec:real_world_validation}

To validate the applicability of our theoretical framework to real-world scenarios, we conducted a series of experiments. We employed the same model architecture used in the synthetic data experiments and pre-trained our models on the large-scale Fineweb-edu dataset. We then analyzed and fitted the data scaling law and model scaling law observed during training. Furthermore, we evaluated the models' grasp of factual knowledge using the PopQA dataset, a popular question-answering benchmark focusing on long-tail knowledge. In our PopQA experiments, we used the object popularity (o\_pop) of a knowledge item in the training set as a proxy for its frequency.

By comparing the empirical exponents from the data and model scaling laws on Fineweb-edu against those from our synthetic experiments (which were generated using knowledge distributions with varying underlying power-law exponents, see discussions surrounding Figure~\ref{fig:data_scaling} and Figure~\ref{fig:model_scaling}), we infer that the knowledge distribution in Fineweb-edu is effectively characterized by a power-law exponent between $1.35$ and $1.5$. This observation provides a potential estimate for the exponent range of real-world knowledge distributions and is consistent with the priors in our theoretical model (e.g., Assumption~\ref{ass:powerlaw}). When fitting the model scaling law, we observed that the inherent irreducible loss for real human language (English, in this case) is considerably larger than that observed for synthetic language. This is intuitive, as the complexity and inherent stochasticity of natural language far exceed the synthetic setting. Note that the irreducible loss values derived from the model scaling law and the data scaling law differ due to our finite model and dataset sizes. In practice, the observed loss can be expressed as $\mathcal{L} = (D/D_0)^{-\alpha} + (M/M_0)^{-\beta} + \varepsilon$, so when fitting the model scaling law (by fixing data size), the estimated irreducible term effectively includes the data-dependent component, i.e., $(D/D_0)^{-\alpha} + \varepsilon$; similarly, when fitting the data scaling law, the irreducible loss includes the model-dependent term, $(M/M_0)^{-\beta} + \varepsilon$. In other words, the discrepancy in irreducible loss arises from the limited model size or data size in each respective setting.

Tests on the PopQA dataset further illuminated the characteristics of knowledge acquisition in real models. Consistent with findings from our synthetic data experiments (see Figure~\ref{fig:model_scaling}(b) and discussion around Figure~\ref{fig:accuracy_freq_modelsize}), we observed a ``knowledge cutoff'' phenomenon in real models as well. Specifically, even with sufficient training, smaller models struggle to learn knowledge that appears with low frequency in the training data. As model size increases, they become capable of learning knowledge at progressively lower frequencies.

\begin{figure}[t] % Or [t] or [!htb] depending on desired placement
    \centering
    \begin{minipage}[b]{0.32\textwidth}
        \centering
        \includegraphics[width=\textwidth,scale=0.95]{fineweb_datalaw.jpg}
        % \caption*{(a)} % If subcaptions are desired and handled by a package
    \end{minipage}
    \hfill
    \begin{minipage}[b]{0.32\textwidth}
        \centering
        \includegraphics[width=\textwidth,scale=0.95]{fineweb_modellaw.jpg}
        % \caption*{(b)}
    \end{minipage}
    \hfill
    \begin{minipage}[b]{0.32\textwidth}
        \centering
        \includegraphics[width=\textwidth,scale=0.95]{popqa.jpg}
        % \caption*{(c)}
    \end{minipage}
    \caption{
    Experimental validation on real-world data.
    (a) Data scaling law on Fineweb-edu: Validation loss as a function of training data size. The fit indicates a power-law relationship, consistent with a power-law distribution of knowledge in the dataset. The estimated exponent for this power-law distribution of knowledge is between $1.35$ and $1.5$.
    (b) Model scaling law on Fineweb-edu: Validation loss as a function of model size. The plot shows a power-law decrease in loss with increasing model parameters, and also suggests a substantially larger irreducible loss component for natural language compared to synthetic data.
    (c) Knowledge cutoff on PopQA: Model performance (e.g., accuracy) on PopQA questions, categorized by the frequency of the underlying knowledge in the Fineweb-edu training set. Smaller models exhibit a clear inability to answer questions about low-frequency knowledge, while larger models demonstrate improved performance on rarer knowledge, effectively lowering the frequency cutoff for knowledge acquisition.
    }
    \label{fig:real_world_validation}
    \vspace{-0.4cm}
\end{figure}
\end{arxiv}

\section{Discussions and Related Work}
\label{sec:relatedwork}

\noindent
{\bf Prediction and Compression:}
The link between prediction and compression is fundamental in both Shannon's probabilistic information theory and Kolmogorov's algorithmic information theory, forming the basis for efficient encoding and decoding of data \citep{cover2006elements,mackay2003information,li2008introduction}. In particular, the better one can predict the distribution of next symbol, the better one can compress the data (via  arithmatic code) and vise versa \citep{rissanen1976generalized,rissanen1983universal,deletanglanguage,li2025lossless}.  
The connection of Kolmogorov's theory to LLMs and artificial intelligence was outlined in 
Hutter's theory of universal intelligence \citep{hutter2005universal}
and Ilya Sutskever's talk at Simons institute \citep{Ilyatalk2023}.
\citep{deletanglanguage} advocate viewing language prediction as a compression problem and show that modern LLMs serve as powerful general-purpose compressors, outperforming traditional text compression tools such as gzip.


\vspace{0.1cm}
\noindent
{\bf Heap's, Zipf's Laws and Generative Models of 
Power Laws:}
Heap’s Law \citep{heaps1978information} 
an empirical relationship stating that the  vocabulary size grows sublinearly ($\propto N^\beta$, for some $\beta$ between 0.4 and 0.7) with the size of a corpus $N$, revealing that new words appear at a diminishing rate. Zipf’s Law 
\citep{zipf2016human}, meanwhile, states that 
\(f(r)\), the frequency of the \(r\)th frequent word, $\propto \frac{1}{r^{\alpha}}$, leading to a heavily skewed distribution dominated by leading terms. These complementary observations expose the
disproportionate frequency of the words, especially “long tail” of rare words, and these long-tail effects are observed not only in linguistics but also in phenomena like city populations, animal populations, website traffic, etc.

One foundational theoretical model that produces
power-law patterns is the classic Simon’s model \citep{simon1955class}, which posits that each new element is more likely to replicate already-popular elements, thus creating a “rich-get-richer” effect. 
Bayesian nonparametric models, such as the Chinese Restaurant Process (CRP) and its generalization, the Pitman–Yor CRP (PYCRP), can also yield power-law distributions while benefiting from the property of exchangeability (the probability of any particular clustering remains unchanged by the order in which data points are observed).

 

\vspace{0.1cm}
\noindent
{\bf Syntax-Knowledge Modeling of Language:}
There is a body of literature that
explicitly separates or conceptually distinguishes syntax and knowledge models within language models. Here we only mention a few representive ones. \citep{dyer-etal-2016-recurrent} proposed a RNN-based model that learn syntactic structures (in the form of parse trees) alongside the generation of words. They did not explicitly incorporate a separate knowledge model. 
\citep{kusner2017grammar} proposed
Grammar Variational Autoencoder combining variational autoencoders (VAEs) with formal grammars to generate syntactically valid structured data, explicitly separates the syntactic and semantic elements.
\citep{konstas2017neural}
proposed the neural Abstract Meaning Representation model that, in some pipelines, first generates a syntactic skeleton, then integrates semantic content from AMR.
Recently, \citep{gong2025disentangling} proposed a stylized data model that comprises disentangled two-type features, to capture respectively syntax and semantics in language, and analyzed the Transformer's learning process, demonstrating that a two-stage dynamic emerges where the syntax component is learned first, which aligns with our \cref{thm:data_scaling_Mutual_Information}.




\vspace{0.1cm}
\noindent
{\bf Scaling Laws:}
The study of neural scaling laws began with the observation that the population loss of trained deep neural networks follows a power-law relationship with respect to dataset size and model parameters. Early work by \cite{rosen2020feldconstructive} introduced a joint error function that captured these dependencies, laying the groundwork for empirical analyses. Henighan et al.\citep{henighan2020scaling} subsequently expanded scaling laws to a broader range of architectures and tasks, while Kaplan et al.\citep{kaplan2020scaling} demonstrated their robustness at vastly larger scales, showing that the loss scales as $L \propto (N/N_{\text{min}})^{-\alpha_N}(D/D_{\text{min}})^{-\alpha_D}$ where \(N\) is the number of parameters, \(D\) is the dataset size, and $\alpha_N$, $\alpha_D$ are scaling exponents. The Chinchilla study \citep{hoffmann2022training} later refined this framework by fitting their scaling law, and then identifying the compute-optimal frontier, showing that many prior models were undertrained, and demonstrating that scaling both parameters and data yields superior performance under fixed compute budgets. 
On the theoretical front, Bahri et al.\citep{bahri2024explaining} distinguished between variance-limited and resolution-limited regimes, identifying four distinct scaling behaviors. Sharma et al.\citep{sharma2020neural} then linked scaling exponents to the intrinsic dimension of data manifolds, highlighting the role of data geometry in performance. 
More recently, Havrilla et al.\citep{havrilla2024understanding} employed statistical and approximation theory to explain transformer scaling laws for low-dimensional data, and Daliri et al.\citep{daliri2024unlocking} established convergence guarantees for 1-bit quantized networks, extending scaling principles to extreme weight-precision settings. 
Recent work has begun to incorporate the learning rate schedule into scaling laws. Empirically, \citep{luo2025multi} proposed a multi-power law to predict entire LLM loss curves under unseen schedules, while \citep{li2025functional} theoretically introduced a {\em functional scaling law} in the kernel regression setting to capture the full loss dynamics under arbitrary learning rate schedules.

Several works have also aimed to explain the origin of power-law behavior observed in scaling laws. Hutter \citep{hutter2021learning} analyzed data scaling law under a stylized setting and showed that when the data distribution follows a power law, the resulting loss curve follows a power-law decay, rather than the common $1/N$ rate. Subsequent works \citep{michaud2023quantization, brill2024neural} extended the setting in \citep{hutter2021learning} and include aspects of model scaling, but still assumed a form for the loss on each data category, without considering the specifics of the learning algorithm. Maloney et al. \citep{maloney2022solvable} derived a power-law form for data scaling of the loss using random matrix theory, linking it to the power-law structure in the spectral properties of the data distribution. 
As in earlier studies, we also derive that a power‑law data distribution is the primary driver of the scaling law in our theoretical model. However, our work differs in the following crucial aspects. 
First, we explicitly link LLMs to compression and the scaling laws to the Kolmogorov structure function, a powerful view that can incorporate more complicated forms of language (rather than just facts). Second, we enrich the data‑generation process with a syntax component, separating from underlying knowledge and yielding a more realistic model of language generation. Third, leveraging the coding/compression view, we rigorously quantify the redundancy for both Bayes‑optimal models and any learning algorithm under mild regularity assumptions. 
%Unlike these prior works, we adopt an information-theoretic and Kolmogorov complexity perspective. We analyze the redundancy of the models that are either learnable by optimal Bayesian algorithms or learning algorithms satisfying certain assumptions, and derive a power-law characterization of their loss.

\vspace{0.1cm}
\noindent
{\bf Bayesian Mixture Code:}
A common approach to constructing a universal code is to employ the \emph{Bayesian mixture code}.
It is well known that this Bayesian mixture code minimizes the Bayes risk/ redundancy 
\citep{aitchison1975goodness} and the Bayes or minimax risk/redundancy can be characterized by the mutual information \(I(\theta; X)\), which is closely tied to channel capacity \citep{gallager1979source,cover2006elements}
(see also \citep{duchi2024lecture}); in the fixed-dimensional case, the worst (capacity-achieving) prior is the Jeffreys prior \citep{clarke1994jeffreys}. Recently, Jeon et al. \citep{jeoninformation} derived a Bayes risk/redundancy upper bound for the family of deep transformers and proposed a Bayesian meta-learning model to explain in-context learning.
\footnote{Although their work did not mention universal coding explicitly, 
their main result can be interpreted in terms of universal coding and redundancy.
}
Our data-generation model draws inspiration from their meta-learning model but differs in two key aspects: we explicitly separate the syntax model 
from the knowledge model, 
and we model the growing and power-law nature of knowledge. 
Jeon and Roy \citep{jeon2024information}
proposed to employ an infinitely wide neural network as a nonparametric data-generation model to explain scaling laws, which is similar to our modeling in spirit. However, their theory predicts that the loss scales as 
\(\widetilde{O}(1/M)\) and \(\widetilde{O}(1/N)\), where \(M\) denotes the model size and \(N\) the data size, which does not capture the power-law scaling behavior observed in LLM practice (with exponents less than 1).


\vspace{0.1cm}
\noindent
{\bf Knowledge Acquisition:}
Researchers have explored multiple mechanisms for external knowledge acquisition in LLMs, including pretraining on massive datasets \citep{chang2024large}, which further reveals a power-law relationship between training steps and the forgetting of memorization and generalization of factual knowledge, and shows that LLMs trained with duplicated training data are more prone to forgetting. In addition, recent studies find that continued pretraining contributes little to factual recall unless overfitting occurs \citep{hoffbauer2024knowledge}, whereas targeted knowledge augmentation—via paraphrasing and diverse retrieval contexts—can more reliably inject domain-specific knowledge into LLMs for RAG applications \citep{bhushan2025systematic}. \citet{mallen2023not} also suggests that LLMs struggle to memorize less popular, long-tail facts even at large scales, and that retrieval augmentation can outperform much larger unassisted models while also improving efficiency. Recently, \citet{gu2025data} demonstrated that knowledge acquisition can exhibit phase transitions with respect to data mixing and model size. In particular, their findings indicate that data with a mixing ratio below a certain threshold (depending on the model size) cannot be effectively learned, a result that is consistent in spirit with our model scaling law (see the proof of Theorem~\ref{thm:model scaling}, in which there is a similar threshold).
Furthermore, \citet{ou2025llms} examined the internal representation of knowledge by introducing the concept of “knowledge circuits,” while addressing challenges related to ensuring factual accuracy and mitigating biases in scaled models. In addition, \citet{allen2024physics} provide a quantitative perspective by estimating the number of knowledge bits a model can store, showing that LLMs are capable of storing up to 2 bits of factual knowledge per parameter and analyzing how architectural choices, training regimes, and data properties influence this capacity. \citet{lu2024scaling} also discussed the relationship between fact knowledge capacity and model size and training epochs, finding they exhibit a linear and negative exponential law relationship.

\vspace{0.1cm}
\noindent
{\bf Cause of Hallucination:} While we focus on hallucinations arising from limited model capacity, prior work has identified a variety of causes. These include erroneous, outdated, or domain-incomplete data in the pretraining corpus \citep{zhang2023vibe}; biased data distributions \citep{ladhak2023pre}; instruction fine-tuning on unfamiliar or underrepresented data \citep{kang2025unfamiliar}; and knowledge shadowing \citep{zhang-etal-2025-law}, where dominant knowledge within the model suppresses less prominent information, leading to the generation of fabricated or inaccurate details.
\cite{kalai2024calibrated} propose a distributional model (on facts) and 
prove that hallucinations must occur at a certain rate if the model satisfies a statistical calibration condition. 

% \jian{their nonparametric model}
\vspace{0.1cm}
\noindent
{\bf Solomonoff's Universal Predictor:}
We would like to mention that Solomonoff introduced a universal Bayesian mixture over all \emph{Turing-computable} predictors \citep{Solomonoff1964a,Solomonoff1964b}, known as the \emph{Solomonoff predictor},
which is a major inspiration of the design of our knowledge model.
It is known that Solomonoff's predictor
achieves universally optimal prediction and compression rates in expectation for any computable sequence-generation process \citep{RathmannerHutter2011}. 
The connection to LLMs and meta-learning is also alluded in 
\citep{deletanglanguage, graulearning}.
However, this predictor remains a purely theoretical model and is not computable in practice. Exploring how modern LLMs might approximate the Solomonoff predictor (or a constrained version of it) is an intriguing direction for future research.




% \subsection{Memorization}
% In this section, we discuss the issue of memorization: is the LLM learning the underlying distribution of the language, or is it memorizing the training data \(X\)?

% Let's revisit the inequality we discussed earlier in inequality \ref{Inequality},
% \begin{align*}
%     K(X) &\le K(P(\cdot|LLM)) - \log P(X|LLM) \\
%     &\le M - \log P(X|LLM).
% \end{align*}
% As we mentioned earlier, for this inequality to hold, at least \(M\) should be no smaller than the minimal sufficient statistic. When \(M\) is much larger than the minimal sufficient statistic, if the equality holds, some memorization is inevitably happening.

% It is also possible that when \(M\) is much larger than the minimal sufficient statistic, the equality of the inequality may not hold. This corresponds to the following problem:
% \[
% \hat{\lambda}_x(\alpha) = \min_S \left\{\Lambda(S) : S \ni x, K(S) = \alpha \right\},
% \]
% where \(\hat{\lambda}_x(\alpha)\) is defined this way. The question is whether \(\hat{\lambda}_x(\alpha)\) is a monotonic function of \(\alpha\).

% 但是如果memorization确实是按我们所说的这种形式发生，它可能并不会像一般所想的那样损害LLM generalization的能力。
% \[K(LLM)-K(P(\cdot|LLM))\le -\log P(X|LLM)-K(X|LLM^*).\]
% $P(\cdot|LLM)$是我们希望LLM学会的部分，实际得到的是LLM，所以左式可以看成是衡量LLM memorization的项。
% \subsection{In-Context Learning}
% The pre-trained model has already learned the modules and the module controller, while the in-context prompt provides conditions to the controller.

% Here, the difference between the in-context prompt and the training data is reflected: training mainly involves learning the modules from the data, while the in-context prompt tells the LLM's controller which pre-trained modules to use. Training on long texts can be thought of as having two parts: one is training the modules, which is continuously done during training; the other is training the controller, which should decide which modules to use based on a small portion of the initial context in the long text. During training, the controller may initially make incorrect decisions about which modules to use, and then improve the controller. With this training process in mind, if a high-quality controller is trained, it should be able to use a small in-context prompt to call the correct modules, thus demonstrating in-context learning.

% Does task vector superposition correspond to module sharing?

% \subsection{Reasoning}
% From the perspective of Reinforcement Learning (RL), intelligence should simultaneously learn to explore and exploit. However, from the perspective of next-token prediction task design, a pre-trained LLM may have only learned the ability to exploit, while the ability to explore is missing. Compression alone does not bring exploration-based intelligence. 
% Search vs. Compress.

% \subsection{Hallucination}

% \textcolor{blue}{Hallucination: learn format before learn knowledge (learning before memorization)}

% From the perspective of compression, format or features should correspond to modules like grammar and semantics that are used almost every time during generation. Since these parts are used every time, they must be compressed as efficiently as possible and learned quickly. On the other hand, knowledge refers to memory, which should only be compressed or remembered well once it appears frequently in the data.

% For example, training data \(X\) may contain many examples of sentences like "The capital of XX country is XX city." When these examples appear infrequently, the model learns only the sentence structure, such as "The capital of XX country is XX city." If this sentence structure is not well compressed, significant compression redundancy will occur. As training progresses, sentences about common countries and capitals, such as "The capital of China is Beijing" or "The capital of the United States is Washington," appear more frequently. These frequently occurring pieces of knowledge must also be remembered; otherwise, compression redundancy will persist.

% In this case, the model size should be partially used to learn the distribution of the language and partially used to remember knowledge. So what kind of knowledge should be remembered first? Apart from the frequently occurring knowledge, it should be those pieces that are poorly compressed, i.e., the training data with higher PPL loss, which should be remembered. A potential validation method could be a knowledge extraction experiment, similar to the "physics of language modeling" experiment, to observe which knowledge is learned in the absence of knowledge augmentation.



\section{Concluding Remarks}
\label{sec:Concluding Remarks}

\begin{arxiv}
We propose a nonparametric generative framework for language, the Syntax-Knowledge model. It has a language syntax encoder,
captured by a parametric model, and a factual knowledge model, captured by the nonparametric 
Pitman-Yor mixture model. Our theoretical analysis under this model demonstrates that the optimal compressors (corresponding to LLMs), first learn to compress frequently occurring syntactic patterns, then progressively acquire knowledge elements in order of frequency. This mechanism explains both data and model scaling laws, as well as the effects of fine-tuning (for either instruction following or knowledge injection). Our findings are consistent with experimental results and several phenomena reported in the literature.

Importantly, our approach is {\em data-centric}: we construct a generative model of language and derive scaling laws under the optimal predictive model (as suggested by Kolmogorov structure function). This perspective sidesteps the mathematical intractability of directly analyzing complex architectures such as multi-layer transformers and the optimization dynamics, while still capturing several important statistical behaviors of LLMs.
\end{arxiv}

There are several exciting directions to extend this work:
\begin{enumerate}
\item
While our analysis has focused on factual knowledge, real-world datasets contain far more diverse forms of information. An important future direction is to incorporate richer knowledge structures—such as relational, procedural, or commonsense knowledge—along with compositional reasoning and inference mechanisms into our theoretical framework.

\item
Our data-centric model could also shed light on how data filtering, domain mixing, and synthetic data generation influence model scaling behaviors. For instance, in data mixing, samples from multiple domains are typically combined in each batch under specific ratios \citep{liuregmix,yedata,chenaioli},. Adjusting these ratios effectively changes the underlying data distribution, which can be modeled as altering the generative probabilities of different clusters in our syntax–knowledge model. 

\item
As shown in our experiments (in \cref{fig:property_capacity_frequency}), power-law data distributions lead to smoother learning curves compared to uniform ones. It would be intriguing to further investigate which data distributions are most suitable for efficient learning, a question that could also inform the design of synthetic training datasets.

\item
Another interesting challenge is understanding how LLMs can approximate universal predictors (e.g., the Solomonoff predictor, which is also based on Kolmogorov complexity) within practical computational constraints (see, e.g., \citep{graulearning,lee2022relaxing}). Bridging these theoretical frameworks with real-world LLMs could deepen our understanding of their behaviors and pave the way for developing models that are more controllable and more reliable.
\end{enumerate}


\begin{arxiv}

\section{Acknowledgment}

We would like to thank Wenjie Huang for providing constructive feedback for an earlier version of the paper. The work is sponsored by Doubao Fund.

    
\end{arxiv}



\bibliography{ICML2025}
% COLM
\bibliographystyle{colm2025_conference}

% APPENDIX

\newpage
\appendix
\onecolumn

\section{Prediction and Compression}
\label{app:prediction_compression}


\subsection{LLMs as Compressors}
\label{app:llmcompressor}


For completeness, we provide a brief introduction to 
{\em arithmetic coding} and explain how Large Language Models (LLMs) can serve as lossless data compressors using this method. More details can be found in, e.g., \citep{rissanen1976generalized, cover2006elements, deletanglanguage,li2025lossless}.

Suppose we have an autoregressive LLM $\llm$ that predicts the next-token probability \(P_{\llm}(x_n \mid x_{1:n-1})\). Arithmetic coding is a popular method for lossless data compression, encoding each symbol (or token) based on its predicted probability. Specifically, arithmetic coding represents a data sequence as an interval within the real line between 0 and 1, sequentially refining this interval based on conditional probabilities. The process is as follows:

\begin{enumerate}
\item \textbf{Initialization:} Start with the interval \([0, 1)\).
\item \textbf{Subdivision:} Divide the current interval into subintervals, each proportional to the probabilities assigned to symbols in the alphabet.
\item \textbf{Encoding tokens:} For each symbol in the input sequence, refine the current interval to the corresponding subinterval associated with that symbol.
\item \textbf{Output:} After processing the entire sequence, select the shortest binary number (base 2) that lies within the final interval as the encoded output.
\end{enumerate}

\begin{example}
Consider an alphabet with symbols \(A\), \(B\), and \(C\), having probabilities \(0.5\), \(0.3\), and \(0.2\), respectively. The initial interval \([0, 1)\) is divided as follows:
\[
A: [0, 0.5),\quad B: [0.5, 0.8),\quad C: [0.8, 1).
\]

To encode the message ``\(AB\)'', start with the interval \([0, 1)\). First, narrow it to \(A\)'s range \([0, 0.5)\). When processing the second symbol \(B\), subdivide the interval \([0, 0.5)\) according to the conditional probabilities:
\[
A: [0, 0.25),\quad B: [0.25, 0.4),\quad C: [0.4, 0.5).
\]

Thus, the interval for ``\(AB\)'' is \([0.25, 0.4)\). We select the number \(0.25\), whose binary representation is \(0.01\).
\end{example}

The efficiency of arithmetic coding directly relates to the predictive accuracy of the underlying LLM. Formally, we have the following proposition
\citep{rissanen1976generalized}
(see also \cite[Ch. 13]{cover2006elements}).

\begin{proposition}
\label{prop:arithmeticcode}
Let \(P_{\llm}\) be the probability distribution predicted by an LLM $\llm$ for a data sequence \(x_{1:n}\). Using arithmetic coding, the code length 
$\codelength(x_{1:n})$ required to encode the sequence \(x_{1:n}\) satisfies:
$$
\codelength(x_{1:n}) \leq -\log P_{\llm}(x_{1:n}) + O(1) = -\sum_{i=1}^n \log P_{\llm}(x_i \mid x_{1:i-1}) + O(1).
$$
\end{proposition}

Consequently, viewing an LLM $\llm$ as a compressor, the total description length of the data \(x\) includes two parts: the complexity \(K(\llm)\) required to describe the model itself (architecture and parameters), and the data encoding length \(\codelength(x)\), bounded as in
Proposition~\ref{prop:arithmeticcode}.

\citep{deletanglanguage,li2025lossless} showed that modern LLMs 
(such as Chinchilla 70B) can serve as powerful general-purpose compressors, significantly outperforming traditional text compression tools such as gzip and LZMA2.

\subsection{Universal Coding and The Coding Game}
\label{subsec:universalcoding_codinggame}
In this appendix, we briefly introduce the concepts 
of universal coding and the coding game and 
how these concepts are connected to perplexity 
minimization in LLMs.


\subsubsection{Universal Coding}

The celebrate Shannon’s source coding theorem (\cite{Shannon1948}) establishes a fundamental limit on achievable code rates by stating that for any code, the average code length $L$
is at least
$H(P)$, 
where 
\(\displaystyle 
H(\phi) =
H(P)=\mathbb{E}_{x \sim P}\bigl[-\log_2 P(x)\bigr]
\)
is the Shannon entropy of the source distribution $P$. Moreover, if we know the source distribution $P$,
we can encode the source message with average code length approaching $H(P)$.


However, real-world data sources often cannot be fully
characterized by a single, known distribution.
Designing coding schemes that adapt effectively to unknown distributions is known as \emph{universal coding}. 
Universal coding
seeks a single coding scheme \(c\) that achieves near-optimal performance
in terms of average code length over every distribution in a given family $\Theta$ (see e.g. \cite[Ch. 13]{cover2006elements}).
The additional cost compared to 
the entropy $H(P)$
is called the \emph{redundancy} of the code $c$.
suppose $\ell(c(x))$ is the length of the code $c(x)$
and $Q_c(x)=2^{-\ell(c(x))}$
is defined to be the predictive probability corresponding to code $c$ ($Q_c$ is a valid probability distribution by Kraft inequality).
The redundancy of code $c$ (or distribution $Q_c$) 
is formally defined as follows:
\begin{align}
\label{eq:reddef}
    \redundency(Q_c, P) & =
    \mathbb{E}_{x \sim P}
    \bigl[\ell\bigl(c(x)\bigr)\bigr]
    -H(P)
    =\mathbb{E}_{x \sim P}
    \bigl[-\log Q_c(x)\bigr]
    -H(P)
     \\
    & = \crossentropy(P\,\|\,Q_c)-H(P) 
    = \DKL(P\,\|\,Q_c). \notag
\end{align}

Hence, finding an efficient universal code that minimizes the redundancy is equivalent to finding a predictive probability that minimizes the cross-entropy.
A central goal in the study of universal coding
is to design efficient universal coding schemes
that can achieve sublinear redundancy
($\redundency = o(n)$ for sequence of length $n$).
Efficient universal coding schemes and tight redundancy bounds have been studied extensively in the information theory literature for a wide variety of distribution families (see, for example, \cite{cover2006elements,rissanen1984universal,clarke1994jeffreys,ziv1978compression,davisson1981efficient,atteson1999asymptotic,feder1996hierarchical}).


\subsubsection{A Coding Game and the Bayesian (Mixture) Strategy}
\label{sec:baye game}
In this subsection, we introduce a coding game
that has many connections to online learning, information theory and Bayesian statistics. 
We mostly follow the exposition in \cite{duchi2024lecture}.

Consider the following online Bayesian coding game.
Suppose $P_\theta$ is a distribution indexed
by $\theta$. Here the set of indices is $\Theta$ and the prior distribution of $\theta$ is $\pi(\theta)$
defined over $\Theta$.
The player's goal is model the distribution $P_\theta$
as well as possible using a distribution $Q$.
The nature chooses the distribution $P_\theta$
(according to the prior)
and sample $n$ random variables $x_i \sim P_\theta$
sequentially.
At step $i$, the player observes the history 
$x_{1:i-1}$,
chooses the distribution $Q(x_i\mid x_{1:i-1})$,
and suffers the log loss
$-\log Q(x_i\mid x_{1:i-1})$.
The overall objective is to minimize the Bayesian log-loss as follow:
$$
\inf_Q \int_\Theta \pi(\theta) \E_{x_{1:n}\sim P_\theta}
\left[\log \frac{1}{Q(x_{1:n})}\right] 
\dd \theta
=
\inf_Q \int_\Theta \pi(\theta) \sum_{i=1}^n\E_{P_\theta}
\left[\log \frac{1}{Q(x_i\mid x_{1:i-1})}\right] 
\dd \theta.
$$
In view of Proposition~\ref{prop:arithmeticcode},
the above objective can also be thought as 
minimizing the Bayesian code length (hence the name
of the game).

By Shannon's source coding theorem,
we can see that the average 
code length cannot be smaller than 
$$
\int_\Theta \pi(\theta) \E_{x_{1:n}\sim P^n_\theta}
\left[\log \frac{1}{P^n_{\theta}(x_{1:n})}\right] 
\dd \theta
=
\int_\Theta \pi(\theta) H(P^n_\theta) 
\dd \theta.
$$
Here the superscript $n$
in $P^n_\theta$ indicates the distribution over 
a sequence of $n$ random variables.
Subtracting this lower bound, the objective becomes
minimizing 
the {\em Bayesian redundancy}:
\begin{align}
& \,\,\, \inf_Q \int_\Theta \pi(\theta) \E_{x_{1:n}\sim P_\theta}
\left[\log \frac{1}{Q^n(x_{1:n})}
-\log \frac{1}{P^n_{\theta}(x_{1:n})}
\right] 
\dd \theta \notag\\
=\,\,\, &
\inf_Q\int_\Theta \pi(\theta) \DKL(P^n_\theta\| Q^n) \dd \theta
\,\,\,=\,\,\,
\inf_Q\int_\Theta \pi(\theta) \redundency(Q^n, P^n_\theta) \dd \theta
\,\,\,\overset{\Delta}{=}\,\,\,
\inf_Q\redundency_n(Q, \Theta)
\end{align}

\noindent
{\bf Bayesian Strategy:}
For a probability family \(\{P_\theta\}_{\theta \in \Theta}\), 
a common approach to constructing a universal code is to employ the \emph{Bayesian strategy} (also called {\em Bayesian mixture code}).
We place some prior distribution \(\pi\) defined over \(\Theta\),
and consider the  Bayesian mixture measure defined as 
$$
Q^n_{\pi}(x_{1:n}) \;=\; \int_\Theta P^n_\theta(x_{1:n})\,\pi(\theta)\,\dd\theta.
$$
At time step $i$, the Bayesian strategy
which uses the Bayes (posterior) estimator 
$$
Q_{\pi}(x_k = x \mid x_{1:k-1}) \;=\; \frac{Q^{k}_{\pi}(x_{1:k})}
{Q^{k-1}_{\pi}(x_{1:k-1})}
$$
as the next-token predictor. 
It is a classic result that this Bayesian mixture code minimizes the Bayesian redundancy 
$
\redundency_n(Q, \Theta) \;=\; \int_\Theta \DKL\bigl(P_\theta \,\big\|\, Q\bigr)\,\pi(\theta)\,\dd\theta
$
\citep{aitchison1975goodness}. In fact, the Bayes redundancy can be characterized by the mutual information \(I(\theta; X)\), which is closely tied to channel capacity \citep{gallager1979source,clarke1994jeffreys,cover2006elements}. 
See also \citep[Ch. 19]{duchi2024lecture} and 
\cite{jeoninformation}. For reader's convenience, we summarize
the results as the following lemma.

\begin{lemma}[\citep{clarke1994jeffreys,duchi2024lecture,jeon2024information}]
\label{lem:baye red}
The minimum Bayesian redundancy is attained by the Bayesian mixture code $Q_{\pi}$, and is equal to the mutual information between random variable $\theta$ (from the prior $\pi$ over $\Theta$) and the data $x_{1:n}$.
\begin{align*}
\inf_Q\redundency_n(Q, \Theta)=
\inf_Q\int_\Theta \pi(\theta) \DKL(P^n_\theta\| Q^n) \dd \theta
=\int_\Theta \pi(\theta) \DKL(P^n_\theta\| Q_{\pi}^n) \dd \theta
= I(x_{1:n}; \theta).
\end{align*}
Here, $\theta\in \Theta$ is sampled from the prior $\pi$,
and $x_{1:n}$ are sampled from $P^n_{\theta}$.
\end{lemma}

\begin{proof}
We provide a simple proof for completeness.
The proof is given in the single‐observation case.  
The exact same argument applies in the \(n\)-observation case, by replacing \(x\) by \(x_{1:n}\).
We first show that minimum Bayesian redundancy is attained by the Bayesian mixture code $Q_{\pi}$.
For any distribution $Q$, we can see that  
\begin{align*}
\redundency_n(Q_{\pi}, \Theta) - \redundency_n(Q, \Theta) &\;=\;
    \int_{\Theta}\!\pi(\theta)\,
      \Bigl[\DKL\!\bigl(P_{\theta}\,\|\,Q_{\pi}\bigr)
            -
            \DKL\!\bigl(P_{\theta}\,\|\,Q\bigr)\Bigr]
   \,d\theta \\
   & \;=\;
   \int_{\Theta}\!\pi(\theta)\,
   \sum_{x} P_{\theta}(x)\,\ln\!\frac{Q(x)}{Q_{\pi}(x)}
   \,d\theta 
   \;=\; 
  \sum_{x} \Bigl[\ln\!\frac{Q(x)}{Q_{\pi}(x)}\Bigr]
            \int_{\Theta}\!\pi(\theta)\,P_{\theta}(x)\,d\theta \\
   & \;=\; 
   \sum_{x}\,Q_{\pi}(x)\,\ln\!\frac{Q(x)}{Q_{\pi}(x)} 
   \;=\; 
   -\,\DKL\!\bigl(Q_{\pi}\,\|\,Q\bigr)
   \;\;\le\;\; 0.
\end{align*}

The second part is simply rewriting the the definition of mutual information, as follows:
\[
\begin{aligned}
  I(X; \theta)
  &=  \mathbb{E}_{(\theta,X)}\!\Bigl[\ln \tfrac{P(X\mid \theta)}{P(X)}\Bigr]
  = \mathbb{E}_{\theta\sim \pi,X\sim P_{\theta}}\Bigl[
       \ln \frac{P_\theta(X)}{\int_{\Theta} \pi(\theta')\,P_{\theta'}(X)\,d\theta'}
     \Bigr] \\
  &= \int_{\Theta} \pi(\theta)\!\int_{X} 
        P_{\theta}(X)
        \,\ln \frac{P_\theta(X)}{Q_{\pi}(X)}
     \,dX\,d\theta \\
  &= \int_{\Theta} \pi(\theta)\,\DKL\!\Bigl(P_{\theta} \,\Big\|\, Q_{\pi}\Bigr)\,d\theta \;=\; \redundency(Q_{\pi},\Theta) .
\end{aligned}
\]
This proves the lemma.
\end{proof}




\section{More on Kolmogorov Structure Functions}
\label{app:kolstructurefunc}
In this appendix, we 
briefly review Kolmogorov structure function
and related concepts.
See \cite[Ch 5.5]{li2008introduction} or \citep{vereshchagin2004kolmogorov,rissanen2005kolmogorov}
for more details.

We consider the case where the data and the code
are expressed in binary bits.
One can describe every binary string $X$ by a two-part description:
the model description in the form of a finite set $S$ that 
contains $X$, and the data-to-model code describing $X$
given $S$ (using $K(X\mid S)$ bits).

In Kolmogorov complexity, we say that a binary string $X$ 
of length $n$ is {\em random} if its Kolmogorov complexity $K(X)=n\pm O(1)$
(i.e., it is impossible to compress $X$ significantly).
Extending this idea, one can define the
{\em randomness deficiency} of $X$ with respect to a set $S$
(such that $X\in S$) as
$$
\delta(X\mid S)\dot{=}\log |S| - K(X\mid S).
$$

\begin{definition}[\textbf{``Best Fit" function}]
It is also called the {\em minimal randomness deficiency} function,
which is defined as:
$$
\beta_X(\alpha) = \min\nolimits_S
\{
\delta(X\mid S) : S \ni X; K(S)\leq \alpha.
\}
$$
\end{definition}
We say $X$ is a {\em typical} or {\em random} element
of a finite set $S$, or $S$ is a fitting model for $X$,
if $X\in S$ and the randomness deficiency
$\delta(X\mid S)=O(1)$.
Basically, it says that to describe $X$ in $S$,
one essentially needs $\log |S|\pm O(1)$ bits
(meaning $X$ is not very special in $S$).
This definition parallels the concept of typical set 
in information theory \cite[Ch. 3]{cover2006elements}.
If $\delta(X\mid S)$ is small enough,
$X$ satisfies {\em all properties} of low complexity
with high probability for the elements
of $S$ (see \citep{vereshchagin2004kolmogorov} for the details).



\begin{definition}[\textbf{Kolmogorov structure function}]
Formally, it is defined as:
$$
h_X(\alpha) = \min\nolimits_S
\{
\log |S| : S \ni X; K(S)\leq \alpha,
\}
$$
for any $\alpha>0$.
The set \(S\) can be viewed as a candidate “typical set” for the data \(X\). One may understand $S$ as the {\em model} of $X$
(the complexity of $S$, 
should be at most $\alpha$)
and the term \(\log\lvert S\rvert\) measures how many bits are needed to single out \(X\) within \(S\). 
\end{definition}

% The set \(S\) can be viewed as a candidate “typical set” for the data \(X\). One may understand $S$ as the {\em model} of $X$
% (the complexity of $S$, 
% should be at most $\alpha$)
% and the term \(\log\lvert S\rvert\) measures how many bits are needed to single out \(X\) within \(S\). 
From the two-part description, it is easy to see that
$
K(X)\leq K(S)+\log |S|+O(1).
$
Hence, we have that 
$$
K(X)\leq \alpha+h_X(\alpha) +O(1).
$$
Hence, $h_X(\alpha)$ cannot be below the {\em sufficiency line} $L\,: \, L(\alpha)+\alpha = K(X)$ (by more than an additive constant).
For those $\alpha$'s such that 
$K(X)\leq \alpha+h_X(\alpha) +O(1)$,
we say the corresponding model {\em sufficient statistics},
and the smallest such $\alpha$ the {\em minimum sufficient statistics}
\citep{gacs2001algorithmic}.
The notion of sufficient statistics here parallels the same notion in the probabilistic setting: a sufficient statistic for the data captures all essential and relevant information and suffices for answering any downstream questions. Any additional information beyond the sufficient statistic is treated as “random noise” with respect to these tasks. For instance, the exact phrasing of a fact may be irrelevant for any downstream tasks unless it is so widely quoted that the particular wording appears frequently in the data, in which case it belongs to the  the sufficient statistic, rather than the noise.

% \begin{figure}[t]
%     \centering
%     \begin{minipage}[b]{0.48\textwidth}
%         \centering
%         \includegraphics[width=\textwidth]{kol_Figure.pdf}
%     \end{minipage}
%     \hfill
%     \begin{minipage}[b]{0.48\textwidth}
%         \centering
%         \includegraphics[width=\textwidth]{test_loss_Figure.pdf}
%     \end{minipage}
%     % \caption{(a) The red dashed curve represents the minimal sufficient statistic, while the black solid curve denotes the Kolmogorov complexity. The minimal sufficient statistic grows sublinearly with data size, exhibiting a power-law trend. This scaling behavior of the minimal sufficient statistic is a key factor driving the power-law decay in loss. The Kolmogorov complexity roughly equals the minimal sufficient statistic plus a linear term with slope less than one.
%     % % \jian{need to draw a test loss.
%     % % Zhixuan pls handles this.} \zhixuan{111111}
%     % (b) This figure presents a 3D illustration of the relationship among data size, model size, and loss. For details on the model size perspective, see \Cref{fig:kol}. }
%     \caption{The red dashed curve in Data Size-Model Size place represents the minimal sufficient statistic, the blue dashed curve in Model Size-Loss place represents the test loss. 
%     % For details on the model size perspective, see \Cref{fig:kol}.
%     (a) An overview of the relationship among data size, model size, and loss.
%     (b) An additional illustration of test loss curve in Model Size-Loss place. We replace the cumulative loss with the average loss on the loss axis to better align with the setting of scaling laws. The test loss of the model trained with $X_{1:n}$ is approximately given by $\frac{h_{\alpha}(X_{1:n'})}{n'},$ where $\alpha$ denotes the minimal sufficient statistics of $X_{1:n}$ and $n' \gg n$ so that the average loss on $X_{1:n'}$ serves as a good approximation of the test loss.
% }
%     \label{fig:kol_3d}
% \end{figure}

In this paper, we use the probabilistic extension of the above definition
(see \cite{grunwald2003kolmogorov,vereshchagin2004kolmogorov}).
As Kolmogorov himself pointed out, the finite set model class is equivalent (up to small additive terms) to the model class of probability density functions
\citep{vereshchagin2004kolmogorov,gacs2001algorithmic}
If \(X\) is genuinely “typical” under the model distribution \(P_{\Model}\) (parametrized by $\Model$), then \(\log\lvert S\rvert\) parallels the 
negative log-likelihood, \(-\log P_{\Model}(X)\) (which is the code length
of $X$ under model $\Model$, ignoring the integer constraint).
Hence, we define 
$$
h_X(\alpha) = \min\nolimits_{\Model}
\{
-\log P_{\Model}(X) : K(P_{\Model})\leq \alpha.
\}
$$
% \textcolor{blue}{\subsection*{Detailed Explanation for Figure 7}
% As introduced in Section 3 (and illustrated in Figure~\ref{fig:kol}), the Kolmogorov structure function for a given data sequence X provides insights into model scaling. Here in Figure~\ref{fig:kol_3d}, we extend this concept to incorporate the data scale. We posit that the ``knowledge'' encapsulated within a natural language dataset expands as the dataset grows. This phenomenon is illustrated in the left panel of Figure~\ref{fig:kol_3d}, where the minimal sufficient statistics increase with the data size.\\
% The right panel of Figure~\ref{fig:kol_3d} illustrates the relationship between the Kolmogorov structure function and the test loss of an LLM. Crucially, the test loss represents the expected loss of a trained model on unseen data, effectively an average loss over an infinitely large test set. In the right panel, this is approximated using a test set of size $n'$ much larger than the training set size $n$ (i.e., $n' \gg n$). To maintain consistency with this definition, the vertical axis in our diagrams represents \textit{average loss} rather than cumulative loss.\\
% Intuitively, guided by the Kolmogorov structure function, the behavior of the test loss as a function of model size can be segmented into two phases around the minimal sufficient statistics. The connection is as follows: Before the model capacity reaches the minimal sufficient statistics for the training data $X_{1:n}$, the optimal test loss corresponds to the average structure function evaluated on a much larger dataset $X_{1:n'}$. As model size increases in this region, the test loss decreases because the model is capturing more of the data's true, generalizable structure. However, once the model capacity surpasses the minimal sufficient statistics, it begins to overfit by memorizing ``pure noise'' specific to the training set $X_{1:n}$. At this point, further increases in model size no longer improve generalization, causing the optimal test loss to plateau.
% \\
% }




The shape of the structure function $h_X(\alpha)$ has been studied by several authors
\citep{vereshchagin2004kolmogorov,rissanen2005kolmogorov,li2008introduction}.
A major results in this line of research is that $h_X(\alpha)$
can essentially assume all possible shapes as along as it satisfies monotonicity, $h_X(0)=|X|$, $h_X(K(X))=0$
and above the sufficiency line
(up to logarithmic additive terms).


\subsection{Kolmogorov Structure Functions with Power-Law Shape}

An instructive question to consider is:
{\em What structure in the data \(X\) causes the structure function \(h_X(\alpha)\) to follow a power-law shape?} Indeed, as pointed out by \cite{hutter2021learning}, different data structures can lead to different learning curves, which in our context correspond to the curve of the Kolmogorov structure function. To build intuition, first consider the example below that naturally yields a piecewise structure. 

% \begin{example}
% Consider a data sequence \(X = (S_1, S_2, \dots, S_m)\), where each element \(S_j\) is a 0/1 string with length $n$. These segments are generated from a mixture of two distinct random sources $P_{\text{mix}}$, defined by probability distributions \(P_1\) and \(P_2\). 
% \begin{itemize}
%     \item \textbf{Source 1 (Major):} \(\text{Source}(P_1)\). 99\% of the segments in \(X\) are drawn from this source.
%     \item \textbf{Source 2 (Minor):} \(\text{Source}(P_2)\). The remaining 1\% of the segments are drawn from this source, which is significantly different from \(\text{Source}(P_1)\).
% \end{itemize}
% We assume that for both $P_1$ and $P_2$, any model $M$ whose capacity is below their own complexity is a poor approximation for either source.
% For instance, we can think of a keyed noisy XOR process $P_K$:
% $
% X_t = \left( \bigoplus_{i=1}^l (X_{t-i} \land K_i) \right) \oplus Z_t
% $
% where $K = (K_1, \dots, K_l)$ is a incompressible secret key string, $\land$ is the logical AND operator, and \(Z_t\) is a Bernoulli noise variable with \(P(Z_t=1) = \epsilon\) for some very small \(\epsilon > 0\). The true entropy rate of this process is almost zero, but it appears as pure random noise to any model $M$ whose capacity is less than the complexity required to describe the secret key $K$. For these low-capacity models, the cross-entropy rate \(\entropyrate(P_K \| M) \approx 1\).

% For simplicity, we assume that the model capacity required to describe the distribution \(P_1\) is approximately the same as that required to describe \(P_2\). We denote this capacity for describing a single source as \(C_S\).

% We now analyze the behavior of the structure function \(h_X(C)\) for the entire sequence \(X\) as the allowed model capacity \(C\) increases.
% We make a standard assumption that \(X\) is a typical sequence drawn from this mixture distribution. Based on the Asymptotic Equipartition Property, we can use the cross-entropy rate $\entropyrate(P\| M)$ as the approximation for the average codelength required to compress a segment from the true source $P$ using a given model $M$.

% At low capacity \(C < C_S\), the set of allowed models $\{M: K(M) \le C\}$ does not include any model with capacity sufficient to approximate $P_1$ or $P_2$. It is to say that the best available model $M_{\text{low}}$ is a poor approximation for both sources, meaning the cross-entropy rates $\entropyrate(P_1\|M_{\text{low}})$ and $\entropyrate(P_2\|M_{\text{low}})$ are both significantly higher than the true entropy rates $H(P_1)$ and $H(P_2)$, respectively. Consequently, the value of $h_X(C)$ remains significantly higher than the true entropy rate of the data.

% As \(C\) increases, it reaches the first threshold $C_S + O(1)$. The set of allowed models $\{M: K(M) \le C\}$ now includes the model for the major source $M_1 \approx P_1$. Consequently, the 99\% of segments originating from $P_1$ can be efficiently compressed, achieving an average codelength of $\entropyrate(P_1\| M_1) \approx H(P_1)$. However, the 1\% of segments originating from $P_2$ are also compressed using $M_1$, resulting in an inefficient average codelength of $\entropyrate(P_2\| M_1)$. The structure function value is thus $h_X(C) \approx nm(0.99 H(P_1) + 0.01\entropyrate(P_2\| M_1))$. This value of $h_X(C_S + O(1))$ is now lower than $h_X(C)$ where $C<C_S$, but still higher than the true entropy rate of the data.

% As \(C\) increases further, it reaches the second threshold \(2C_S + O(1)\). The set of allowed models $\{M: K(M) \le C\}$ now includes the full two-source mixture model $M_{\text{mix}}\approx P_{\text{mix}}$. Consequently, all of the data can be efficiently compressed, achieving an average codelength of $\entropyrate(P_{\text{mix}}\|M_{\text{mix}}) \approx H(P_{\text{mix}})$, which is the true entropy rate of the data source.

% As \(C\) beyond \(2C_S+O(1)\), no new structural compression is possible. While the value of $h_X(C)$ will continue to decrease as more complex models can merely serve to memorize the pure
% randomness within the specific data $X$.
% \end{example}


\begin{example}
Consider a data sequence \(X = (S_1, S_2, \dots, S_m)\), where each element \(S_j\) is a 0/1 string with length $n$. These segments are generated from a mixture of three distinct random sources $P_{\text{mix}}$, defined by probability distributions \(P_1\), \(P_2\), and \(P_3\).
\begin{itemize}
    \item \textbf{Source 1:} \(\text{Source}(P_1)\). 60\% of the segments are drawn from this source,
    \item \textbf{Source 2:} \(\text{Source}(P_2)\). 30\% of the segments are drawn from this source,
    \item \textbf{Source 3:} \(\text{Source}(P_3)\). 10\% of the segments are drawn from this source.
\end{itemize}
We begin by defining the entropy rate. For a random process $P$, the entropy rate $\entropyrate(P)$ is defined as:
$ \entropyrate(P) := \lim\limits_{n \to \infty} \frac{1}{n} \mathbb{E}[-\log P(X_1, \dots, X_n)] $
where $(X_1, \dots, X_n)$ is a sequence drawn from $P$. This limit is known to exist for any stationary and ergodic process \citep{cover2006elements}. 
Throughout this section,
We assume that the entropy rate, as well as the analogously defined cross-entropy rate $\entropyrate(P\|Q)$ and KL-divergence rate $\KLrate(P\|Q)$, exist.
%we assume that all sources under consideration are stationary and ergodic processes. 

For simplicity of discussion, we further assume that the entropy rate of these three sources and the model capacity required to describe each source are approximately the same. We denote this capacity for describing a single source as \(C_S\). 
Moreover, we assume that all three sources cannot be approximated by any model $M$ whose capacity is below their own complexity.
\footnote{
The precise form of the models is not important 
in our discussion. But such a model exists:
For instance, a keyed noisy XOR process $P_K$:
$
X_t = \left( \bigoplus_{i=1}^l (X_{t-i} \land K_i) \right) \oplus Z_t
$
where $K = (K_1, \dots, K_l)$ is an incompressible secret key string, $\land$ is the logical AND operator, and \(Z_t\) is a Bernoulli noise variable with \(P(Z_t=1) = \epsilon\) for some very small \(\epsilon > 0\). 
}

We consider the asymptotic case as $n\rightarrow \infty$. Based on the Asymptotic Equipartition Property \citep{cover2006elements}, we can use the entropy rate $\entropyrate(P)$ as the approximation for the average codelength required to compress a segment from source $P$. Similarly, we use the cross-entropy rate $\entropyrate(P\| M)$ as the approximation for the average codelength required to compress a segment from source $P$ using model $M$. 
As $n\rightarrow \infty$, the entropy rate of the mixture of the three sources can be decomposed as\citep{beirami2014universal}
: $\entropyrate(P_{\text{mix}})= 0.6\entropyrate(P_1)+0.3\entropyrate(P_2)+0.1\entropyrate(P_3)$.
\footnote{Corollary 3 in \citep{beirami2014universal} provides the asymptotic decomposition for the entropy of a mixture of parametric sources when the number of mixture sources is sufficiently small relative to the sequence length. The total mixture entropy $H_n(P_{\text{mix}})$ for a sequence of length $n$ is given by:$ H_n(P_{\text{mix}}) = \sum_{i=1}^K \pi_i H_n(P_i) + H(\pi) + o(1) $
where $H_n(P_{\text{mix}})$ is the total entropy of the mixture source, $H(\pi)$ is the entropy of the mixture weights and $\pi = (\pi_1, \dots, \pi_K)$ are mixture weights. When we consider the entropy rate, the terms on the right-hand side that are independent of $n$, namely $H(\pi) + o(1)$, vanish. We thus obtain the decomposition of the entropy rate $ \entropyrate(P_{\text{mix}}) = \sum_{i=1}^K \pi_i \entropyrate(P_i) $.}



We now analyze the behavior of the structure function \(h_X(C)\) for the entire sequence \(X\) as the allowed model capacity \(C\) increases.


At low capacity \(C < C_S\), the set of allowed models $\{M: K(M) \le C\}$ does not include any model with capacity sufficient to describe a single source. By our assumption, the available model $M_{\text{low}}$ is a poor approximation for all three sources, meaning the cross-entropy rates $\entropyrate(P_i\|M_{\text{low}})$ for $i=1,2,3$ are all significantly higher than the true entropy rates $H(P_i)$. Consequently,  $h_X(C)$ remains significantly higher than the true entropy rate of the mixture.

As \(C\) increases, it reaches the first threshold $C_S + O(1)$. The set of allowed models $\{M: K(M) \le C\}$ now includes models for the individual sources (e.g., $M_i \approx P_i$ for $i=1,2,3$). Clearly, the best model for the whole data $X$ will be $M_1$, as it captures the largest component, resulting the largest redundancy reduction. Consequently, 60\% of segments from $P_1$ are compressed efficiently, achieving average codelength of $\entropyrate(P_1\|M_1)\approx H(P_1)$. However, the segments from $P_2$ and $P_3$ are compressed inefficiently using $M_1$. The average structure function is thus:
\begin{align*}
\dfrac{1}{nm}h_X(C_S+O(1)) & \approx 0.6 \entropyrate(P_1) + 0.3 \entropyrate(P_2\| M_1) + 0.1 \entropyrate(P_3\| M_1)) \\ & = 0.6 \entropyrate(P_1) + 0.3\entropyrate(P_2)+ 0.1\entropyrate(P_3) + 0.3 \KLrate(P_2\| M_1) + 0.1\KLrate(P_3\|M_1)
\\ & \approx \entropyrate(P_{\text{mix}})+0.3 \KLrate(P_2\| P_1) + 0.1\KLrate(P_3\|P_1).   
\end{align*}
This value is now lower, but still higher than the true entropy rate because 40\% of the data originating from $P_2, P_3$ cannot be compressed efficiently.

As \(C\) increases further, it reaches the second threshold \(2C_S + O(1)\). The set of allowed models $\{M: K(M) \le C\}$ now includes models that are mixtures of two sources. The best available model will be $M_{1,2}\approx P_{1,2}$, where $P_{1,2}=\frac{0.6P_1+0.3P_2}{0.9}$ covers the two largest components. This model efficiently compresses 90\% of the data, but still compresses the 10\% from $P_3$ inefficiently. The average structure function is thus: 
\begin{align*}
    \dfrac{1}{nm}h_X(2C_S+O(1))  \approx 0.9\entropyrate(P_{1,2})+0.1\entropyrate(P_3\|M_{1,2})
     \approx \entropyrate(P_{\text{mix}})+ 0.1\KLrate(P_3\|P_{1,2}).
\end{align*}
This value is now lower, but still higher than the true entropy rate because 10\% of the data originating from $P_3$ can not be compressed efficiently.

As \(C\) increases further, it reaches the third threshold \(3C_S + O(1)\). The set of allowed models $\{M: K(M) \le C\}$ now includes the full three-source mixture model $M_{\text{mix}}\approx P_{\text{mix}}$. Consequently, all of the data can be efficiently compressed, achieving an average structure function: 
$$
\dfrac{1}{nm}h_X(3C_S+O(1))\approx\entropyrate(P_{\text{mix}}\|M_{\text{mix}}) \approx \entropyrate(P_{\text{mix}}),
$$ 
which is the true entropy rate of the data source.

As \(C\) beyond \(3C_S+O(1)\) (the minimal sufficient statistics), no new structural compression is possible. While $h_X(C)$ will continue to decrease as more complex models can merely serve to memorize the pure
randomness within the specific data $X$
($h_X(C)$ coincides with the sufficiency line).
\qed
\end{example}

\begin{figure}
    \centering
    \includegraphics[width=0.7\linewidth]{Kol_model_figure.pdf}
    \caption{Illustration of the structure function $h_X(C)$ for a finite mixture of sources.  The gentle decrease (blue line) represents memorization rather than structural compression. The sharp vertical drops (black line) represent structural compression, which occurs when the model capacity $C$ becomes sufficient to capture a new source component (at $C_S$, $2C_S$, etc.). The red dashed line represents the structure function curve if the mixture components could be approximated by low-complexity models.}
    \label{fig:placeholder}
\end{figure}

The preceding example illustrates how a finite mixture of sources creates a piecewise structure. We can now build upon this intuition to understand the power-law case by generalizing from a finite to an infinite mixture model.

Consider that the data \(X\) consists of segments generated by an infinite mixture of components, where the \(i\)-th component distribution $P_i$ is associated with a mixture probability \(\pi_i\), satisfying $\sum_{i=1}^\infty \pi_i = 1$. We assume these that components are ordered by their probability, such that \(\pi_1 \ge \pi_2 \ge \dots\). We define the mixture of the $j$-th to $k$-th components as $P_{j:k} = \frac{\sum_{i=j}^{k} \pi_i P_i}{\sum_{i=j}^{k} \pi_i}$. We make the same assumption that
the entropy rate $\entropyrate(P_{j:k})$ can be approximated as: $\entropyrate(P_{j:k})\approx \sum\limits_{i=j}^k \pi_i\entropyrate(P_i).$



Following the spirit of the previous example, a model with a limited capacity \(C\) can only capture a finite number, say \(M(C)\), of these components. The best available model at this capacity level will naturally capture the mixture of the \(M(C)\) most probable components, as this yields the largest compression gain. The model then compresses the data originating from the remaining tail components inefficiently.
% The structure function $h_X(C)$ will thus be higher than the irreducible entropy of the data. The redundancy $R(C)$ is then defined as the extra description length needed for the uncaptured tails. Assuming the model code length required to describe the parameters of each of the components and the redundancy for the uncaptured components are approximately constants. Then the redundancy can be approximated as
The average structure function is thus:
\begin{align*}
    \dfrac{1}{|X|}h_X(C) & \approx
 \left(\sum\limits_{i=1}^{M(C)}\pi_i \right) \entropyrate(P_{1:M(C)}) + \left( \sum\limits_{i=M(C)+1}^{\infty}\pi_i\right)  \entropyrate\left(P_{M(C)+1:\infty}\|P_{1:M(C)}\right)
 \\ & = \left(\sum\limits_{i=1}^{M(C)}\pi_i \right) \entropyrate(P_{1:M(C)}) +  \sum\limits_{i=M(C)+1}^{\infty}\pi_i \left(\entropyrate\left(P_{i}\right)+\KLrate(P_i\|P_{1:M(C)})\right)
 \\ & \approx \entropyrate(P_{1:\infty}) + \sum\limits_{i=M(C)+1}^{\infty}\pi_i\KLrate(P_i\|P_{1:M(C)})
\end{align*}




% Furthermore, we introduce the following simplifying assumptions:
% \begin{itemize}
%     \item[1.] The description length of each component is a constant, i.e., $M(C)=\Theta(C)$.
%     \item[2.] The redundancy associated with uncaptured components is a constant, i.e., \\
%     $\KLrate\left(P_{M(C)+1:\infty}\|P_{1:M(C)}\right)=\Theta(1)$.
% \end{itemize}

Furthermore, we introduce two simplifying assumptions. First, the description length of each component is assumed to be a constant, i.e., $M(C)=\Theta(C)$. Second, the redundancy associated with uncaptured components is assumed to be a constant, i.e., $\KLrate\left(P_i\|P_{1:M(C)}\right)=\Theta(1)$ for all $i$.

Then the redundancy $R(C)$ is thus:
\begin{align*}
    R(C):=\dfrac{1}{|X|}h_X(C)-H(P_{1:\infty})\approx \left(\sum\limits_{i=M(C)+1}^{\infty}\pi_i\right)\KLrate\left(P_{M(C)+1:\infty}\|P_{1:M(C)}\right) = \Theta\left(\sum\limits_{i=M(C)+1}^{\infty}\pi_i\right).
\end{align*}




The crucial insight is as follows: If the mixture probabilities \(\pi_i\) themselves follow a power-law distribution, for instance:
\[
\pi_i \propto i^{-\beta} \quad (\text{for } \beta > 1).
\]
Then the redundancy also follows a power law:
\begin{align}
\label{eq:kol_model_scaling}
    R(C)=\Theta\left(\sum_{i=M(C)+1}^{\infty} i^{-\beta} \right) = \Theta\left(\int_{M(C)}^\infty x^{-\beta} dx \right)= \Theta\left(M(C)^{1-\beta}\right)=\Theta\left(C^{1-\beta}\right).
\end{align}
Conversely, if we empirically observe that the structure function's redundancy follows a power law, we can infer that the underlying data structure also exhibits a power law. This bidirectional relationship can serve as a theoretical bridge connecting empirically observed phenomena in language, such as Zipf's Law and Heaps' Law, with the model scaling observed in large language models. 

We also note that the Kolmogorov structure function is closely related to the rate-distortion function in Shannon's information theory \cite{grunwald2003kolmogorov}. The rate-distortion function can be viewed as the probabilistic counterpart to the Kolmogorov structure function, which is defined in terms of algorithmic complexity. Following this conceptual link, our study of model scaling laws in \Cref{appsec:model scaling} leverages the rate-distortion framework to rigorously prove the conclusions corresponding to the power-law intuition developed above. More specifically, in \Cref{thm:model scaling}, we considered the case of $\pi_i \propto i^{-1/\alpha}$ and used rate-distortion theory to derive an optimal model redundancy of $\Theta(C^{-1/\alpha + 1})$ for a model size constraint $C$, which aligns with our preceding result \eqref{eq:kol_model_scaling}. 





% Inspired by the example, a power-law shape requires the model to successively capture finer-grained aspects. Concretely, text whose event frequencies follow a power-law distribution represents such a natural and realistic candidate: at each step up in complexity, the model encodes additional subsets of minorr events, resulting in a gradual decrease in \(\log|S|\).





% We also note that Kolmogorov structure function
% is closely related to the rate-distortion function in Shannon's
% information theory, see e.g., \cite{grunwald2003kolmogorov}.
% In our study of model scaling law in \Cref{appsec:model scaling}, we leverage the concept of
% rate-distortion function as well.






% \textcolor{blue}{Then we intuitively answer the question {\em What structure in the data \(X\) causes the structure function \(h_X(\alpha)\) to follow a power-law shape?}. We have illustrated the Kolmogorov structure function in a classical, finite setting in \cref{sec:kolmogorov} as shown in the left panel of \Cref{fig:kol}. In this setting, we consider a finite-length sequence $X$ and the model code length for $X$ in its two-part code description is also finite. However, natural language exhibits a property where knowledge scales with data, compelling us to consider the case when model code length grows to infinity as the data size grows to infinity. To build intuition, we consider a simplified scenario that satisfies the above condition: data generated by an infinite-dimensional mixture model. We also assume the components of this model are encoded independently (i.e., the code length for a set of components is approximately the sum of their individual code lengths) for simplicity. We explore the intuition for model scaling in the asymptotic regime where data size grows infinitely large relative to the model size. Given a model size sufficient to encode $M$ components, the redundancy induced by this model size constraint(See in \Cref{fig:kol}) corresponds to the sum of the tail mixture probabilities (i.e., the sum of all mixture probabilities excluding the $M$ largest). If the mixture probabilities of this model with infinite dimensions follow a power law, then the resulting redundancy (the sum of a tail that also follows a power law) also follows a power law. Conversely, if the sum of the tail probabilities follows a power law, we can also deduce that the mixture probabilities themselves follow a power law.
% This bidirectional relationship provides the intuition for our model scaling theory in \Cref{appsec:model scaling}, and it can also serve as a theoretical bridge connecting empirically observed phenomena in language, such as Zipf's Law and Heaps' Law, with the model scaling observed in large language models. 
% \\
% We also note that Kolmogorov structure function
% is closely related to the rate-distortion function in Shannon's
% information theory. Specifically, the expected Kolmogorov structure function can be shown to be approximately equal to the rate-distortion function, see e.g., \cite{grunwald2003kolmogorov}.
% In our study of model scaling law in \Cref{appsec:model scaling}, we leverage the concept of
% rate-distortion function as well.
% \\
% As illustrated above, the Kolmogorov structure function provides insights into model scaling. In Figure~\ref{fig:kol_3d}, we extend this concept to incorporate the data scaling. Following the above intuition, we consider that the data is generated from an infinite-dimensional power-law mixture model. In this case, the minimal sufficient statistics with respect to the model size also exhibit a power law, as shown in the left panel of \Cref{fig:kol_3d}.
% \\
% The right panel of Figure~\ref{fig:kol_3d} illustrates the relationship between the Kolmogorov structure function and the test loss of an LLM training with finite data. Crucially, the test loss represents the expected loss of a trained model on unseen data, effectively an average loss over an infinitely large test set. In the right panel, this is approximated using a test set of size $n'$ much larger than the training set size $n$ (i.e., $n' \gg n$). To maintain consistency with this approximation, the y-axis in the right panel of \Cref{fig:kol_3d} represents average loss($h_{\alpha}(X_{1:n})/n$) rather than cumulative loss($h_{\alpha}(X_{1:n})$).
% \\
% Intuitively, guided by the two-part code description and Kolmogorov structure function, the behavior of the test loss as a function of model size can be segmented into two phases around the minimal sufficient statistics. The connection is as follows: Before the model capacity reaches the minimal sufficient statistics for the training data $X_{1:n}$, the optimal test loss corresponds to the average structure function evaluated on a much larger dataset $X_{1:n'}$. As model size increases in this region, the test loss decreases because the model is capturing more of the data's true, generalizable structure. However, once the model capacity surpasses the minimal sufficient statistics, it begins to overfit by memorizing pure noise specific to the training set $X_{1:n}$. At this point, further increases in model size no longer improve generalization, causing the optimal test loss to plateau.
% \\
% Therefore, the optimal redundancy achievable given the training data $X_{1:n}$ is equivalent to the redundancy obtained with infinite data, but under a model size constraint determined by the minimal sufficient statistics of $X_{1:n}$. We have previously established that for infinite data, the redundancy exhibits a power-law relationship with the model size constraint. Furthermore, as we show in the left panel of \Cref{fig:kol_3d} and demonstrated before, the size of the data and the size of its minimal sufficient statistics also follow a power-law relationship. Combining these two factors implies that the optimal redundancy for training data $X$ also follows a power-law relationship with respect to the size of training data.}










\eat{
Sometimes, it is more convenient to the following closely related "MDL" function, used in 
the Minimum Description Length (MDL) theory
(see e.g., \cite{li2008introduction,rissanen1978modeling,
barron1998minimum,grunwald2007minimum}):

\begin{definition}[\textbf{"MDL" function}]
    The "MDL" function $\lambda_{X}(\alpha)$ denotes the length of the minimal two-part description length for \(X\) consisting of the model code length \(K(S)\) and the compressed length of \(X\) given model \(S\) with the constrain that the Kolmogorov complexity of the model $S$ is not larger than $\alpha$. Formally,
$
\lambda_X(\alpha) = \min_S \{\Lambda(S) :  K(S) \leq \alpha\},
$
where \(\Lambda(S) = K(S) -\log P_S(X)  \geq K(X) - O(1)\) is the total length of two-part code of \(X\) with help of model \(M\). The smallest 
$\alpha$ that satisfies $\lambda_X(\alpha)+O(1)=K(X)$ is called \textbf{Minimal Sufficient Statistics} of $X$ (\cite{li2008introduction}).
\jian{notation inconsistent with previous discussions}
\zhixuan{We need to use M to replace S?}
\end{definition} 
}
\section{Details of the Syntax-Knowledge Model}
\label{app:datageneration}

%We assume that the data consists of a collection of independent sentences, each represented as a sequence of tokens of finite length. Sentences are generated by a two-stage stochastic process: first, a piece of knowledge is generated; second, a complete sentence is generated based on the generated knowledge. The second stage which generates sentences based on knowledge represents linguistic components such as grammatical structures. It does not grow indefinitely as the volume of generated data increases, so we assume that it can be characterized by a finite-dimensional parameter $\boldsymbol{\theta}$. On the other hand, the first stage which generates knowledge is different, as knowledge continuously increases as data increases. This stochastic process is considered a non-parametric or infinitely-dimensional stochastic process. 

In this section, we propose a hierarchical data generation model, called the {\em syntax-knowledge} model. In this model, where each sentence in the training set is generated by a {\em syntax encoder} that encode a (factual) knowledge element, sampled from the {\em knowledge model}. 
In our model, the syntax model (e.g., a probabilistic CFG or English grammar) does not grow indefinitely with the size of the dataset. Therefore, we assume that the distribution of the syntax model can be parametrized using a finite-dimensional parameter $\synModel$. On the other hand, the knowledge model employs a {\em non-parametric} stochastic process to account for two empirically observed phenomena: 1) the unbounded growth of factual information as datasets grow (mirroring Heap's Law in lexical growth patterns \citep{heaps1978information}), and 2) the long-tailed frequency distribution of factual occurrences, analogous to Zipf's law in natural language \citep{zipf2013psycho,zipf2016human}. 

To operationalize this above idea, we leverage the {\em Pitman–Yor Mixture Model(PYMM)} \citep{pitman1997two} for modeling the knowledge model. The preferential attachment mechanism 
in Pitman-Yor Process naturally captures both the {\em sublinear scaling} of new factual discoveries and 
the power-law distributed frequencies of knowledge pieces.
% We discuss the limitation and possible extension of 
% our data generation model in Appendix~\ref{sec:limit}.

\vspace{0.1cm}
\noindent
{\bf Notations:} We denote the conditional probability \( p(x \mid \phi) \) as \( p_\phi(x) \), where \( \phi \) is the latent variable that determines this probability distribution.
% Let \(\Voc\) be the set of all tokens. For simplicity, 
% we assume that $\Voc$ can be divided into two categories: 
% syntax-related tokens, denoted \(\Vocsyn\), and knowledge-related tokens, denoted \(\Vocknow\).
The syntax model is parameterized by $\synModel$, the knowledge model is denoted as $\knwModel$, and the entire data model is denoted as 
$\dataModel=\{\synModel,\knwModel\}$. 

% \jian{please use my notations...}

The syntactic elements generated by the syntax model are denoted by 
\( \Synelem_{1:N} \), and the knowledge elements generated by the knowledge model are denoted by \( \Knwelem_{1:N} \). 
The corpus consists of sentences \( X_{1:N} \), where each sentence \( X_i \) is generated by the syntax model from a syntax–knowledge pair \( (\Synelem_i, \Knwelem_i) \). 
Now, we provide the details below.




% \jian{please look at the notation I wrote in previous sections....}
% Our knowledge modeling  proceeds as follows: each piece of factual knowledge is represented as a pair \((x, y)\), and the model is trained to map \(x\) to \(y\). One may interpret \(x\) as a “question” and \(y\) as the corresponding “answer,” or more generally view \(x\) as the identifier of a knowledge piece and \(y\) as its content.
% Define \(\Vocknow^*= \bigcup_{i=1}^{\infty} (\Vocknow)^i\) as the set of all nonempty sequences formed by tokens from \(\Vocknow\) . Then the set of all possible of $x,y$ belong to $\Vocknow^*$, denoted as $X,Y$ respectively.


% \jian{I don't understand these notations at all. 
% what is the domain of x and y.
% The description like "item ID" is horrible. 
% the reviewer will stop reading right here.
% }

\subsection{The Non-Parametric Knowledge Model}
\label{sec:PYMM}

We model the knowledge model as a Pitman-Yor Mixture Model ($\PYMM$), which is a {\em nonparametric Bayesian} model. 
We first introduce the Pitman-Yor Process, from which the $\PYMM$ is constructed.


\noindent
{\bf Pitman-Yor Process:}
The {\em Pitman-Yor Process (PYP)}, also known as the {\em Pitman-Yor two-parameter Poisson–Dirichlet process}, extends the Dirichlet Process (DP) and the Chinese Restaurant Process \citep{pitman1997two}. It has been applied
to model a growing number of topics in the literature of 
topic modeling \citep{teh2006hierarchical, goldwater2011producing}. A Pitman–Yor Process (PYP) is characterized by two real-valued parameters: the \emph{discount parameter} \( 0 \leq \alpha < 1 \) and the \emph{concentration parameter} \( \beta > -\alpha \), along with a base probability measure \( \basemeasure \). We denote the process as \( \PYP(\alpha, \beta, \basemeasure) \).


A sample from the Pitman–Yor process \( \PYP(\alpha, \beta, \basemeasure) \) is a  
random probability measure \( \knwModel = \sum_{i=1}^{\infty} p_i \delta_{\phi_i} \),  
which is a discrete distribution with countably infinite atoms, where:
\begin{itemize}
    \item Each atom \( \phi_i \sim \basemeasure \) is the $i$-th cluster parameter,  independently drawn from the base measure \( \basemeasure \);
    \item \( \delta_{\phi_i} \) denotes the Dirac delta measure centered at \( \phi_i \);
    \item \( p = (p_1, p_2, \ldots) \) are the weights generated by the Pitman–Yor Chinese Restaurant Process ($\PYCRP$) (described below).
\end{itemize}
% We denote this sampling procedure as \( p = (p_1, p_2, \ldots) \sim \PYCRP(\alpha, \beta) \).

Imagining a restaurant where customers arrive sequentially, each choosing either to join an existing lively table or start their own, resulting in a naturally evolving clustering structure.

\begin{itemize}
    \item The first customer sits at a new table.
    \item 
Suppose \( N_k \) is the number of customers already at table \( k \), and \( K \) is the current number of occupied tables.
The \( n \)-th customer either joins an existing table \(k\) or starts a new table with the following probabilities:
\[
\text{For the \(n\)-th customer:} \quad
\begin{cases}
\displaystyle \frac{N_k - \alpha}{n - 1 + \beta}, & \text{if joining an existing table \(k\),}\\[10pt]
\displaystyle \frac{\beta + \alpha K}{n - 1 + \beta}, & \text{if starting a new table,}
\end{cases}
\]
% \jian{we may want a figure in appendix. customs and tables.
% it looks nice.}
\end{itemize}
As the number of customers \( n \to \infty \), the relative sizes of the tables (i.e., the proportion of customers at each table) converge in distribution to a random probability vector(\cite{pitman2006combinatorial},\Cref{thm:PYP rank fre}), denoted as:
\[
p=(p_1, p_2, \dots)=\lim\limits_{n\to \infty}(N_1/n,N_2/n,\cdots) \sim \PYCRP(\alpha, \beta)
\]
This probability vector $p=(p_1, p_2, \dots)$ asymptotically exhibits a power-law distribution(see \Cref{lem:PYCRP fre}).

% \jian{mention the power law lemma. zhixuan}
\noindent
{\bf Pitman–Yor Mixture Model (PYMM):}
In the Pitman–Yor Mixture Model, the Pitman–Yor Process  
\( \PYP(\alpha, \beta, \basemeasure) \) 
serves as a prior over the space of mixture distributions. 
In particular, a random probability measure 
\( \knwModel = \sum_{i=1}^{\infty} p_i \delta_{\phi_i} \), 
sampled from \( \PYP(\alpha, \beta, \basemeasure) \), 
can be interpreted as a discrete mixture model: to generate a knowledge element from \( \knwModel \), 
one first selects index \( i \in \mathbb{N} \) with probability \( p_i \), 
then draws a sample from the conditional distribution \( P_{\phi_i} \), 
where \( \phi_i \sim \basemeasure \) is the $i$-th cluster parameter.

The knowledge elements generated by the knowledge model are abstract representations that capture the factual component underlying a sentence. 
We assume that the support of the knowledge element lies in a discrete set \( \knwset \) (i.e., each knowledge element
$\Knwelem\in \knwset$). 
Although the cardinality \( |\knwset| \) may be potentially very large, we assume that \( \log |\knwset| \) 
is quite small (since \( \log |\knwset| \) is roughly the number of bits or tokens required to encode the knowledge element), and particularly much smaller compared to \( N^\alpha \) for any $\alpha>0$.
%\jian{we need to be careful here.}
% \zhixuan{do we need this sentence here?}

% \jian{stop here.}

We assume that the base measure \( \basemeasure \), from which the knowledge cluster parameters \( \phi_i \) are drawn, is supported on the bounded parameter space 
\( \knwfamily = \{ \knwModel \in \mathbb{R}^{\knwdimension} : \|\knwModel\|_{2} \le 1 \} \).
 Given the parameter spaces \( \knwfamily \), we define corresponding parametric families of probability distributions:
$
\knwprobability = \left\{ P(\cdot\mid \knwModel) : \knwModel \in \knwfamily \right\},
$
where \( P(\cdot\mid\knwModel) \) denotes the conditional distribution of the knowledge model given parameter \( \knwModel \).



\subsection{Parametric Syntax Model}
% The syntax model generates a syntax model conditioned on a piece of knowledge. We assume that when conditioned on all possible pieces of knowledge, there are fixed number \( \numsyntax < \infty \) different distributions of generated syntax, each determined by a corresponding latent variable \( \phi_s^{(i)} \)(sampled from prior $\synPrior$). Thus, the entire syntax model can be parameterized by 
% \(
% \synModel=\{\phi_s^{(1)},\phi_s^{(2)},\cdots,\phi_s^{(\numsyntax)}\}
% \). 
% Specifically, we partition all knowledge into $\numsyntax$ distinct sets, denoted by $\knwset_1, \ldots, \knwset_{\numsyntax}$.
The syntax model generates syntactic elements conditioned on knowledge elements. 
Each knowledge element \( \Knwelem \in \knwset \) deterministically selects an index \( i \in \{1, \dots, \numsyntax\} \), corresponding to a latent variable \( \synModel^{(i)} \) that parameterizes the conditional syntax model. 

Conditioning the generation of syntax on a knowledge element \( \Knwelem \) is thus equivalent to conditioning on the corresponding latent variable:
\[
P(\cdot \mid \synModel, \Knwelem) = P(\cdot \mid \synModel^{(i)}).
\]
Accordingly, the overall syntax model is parameterized as:
\[
\synModel = \{ \synModel^{(1)}, \synModel^{(2)}, \dots, \synModel^{(\numsyntax)} \}.
\]


The syntactic elements are abstract representations capturing surface-level variation such as word choice (e.g., synonyms), phrase structure (e.g., active vs. passive voice), and other randomness that does not affect the underlying semantics.


We assume that the prior distribution \( \synprior \), from which the conditional syntax model parameters \( \synModel^{(i)} \) are drawn, is supported on the bounded parameter space 
\( \synfamily = \{ \synModel \in \mathbb{R}^{\syndimension} : \|\synModel\|_{2} \le 1 \} \) 
for all \( 1 \le i \le \numsyntax \).
 Given the parameter spaces \( \synfamily \), we define corresponding parametric families of probability distributions:
\[
\synprobability = \left\{ P(\cdot\mid \synModel) : \synModel \in \synfamily \right\},
\]
where \( P(\cdot\mid \synModel) \) denotes the conditional distribution of the syntax model given parameter \( \synModel \).

% \jian{syntax component means the entire
% syntax components. we need different name. Zhixuan}


% Each sample drawn from the syntax model generates a sentence template that does not contain knowledge. Specifically, the syntax model generates a sentence template with a special token $\knowledge$. This token indicates the presence of a knowledge-related token instead of directly providing the specific token.

% For example, for the sentence "Alice was born on 2000," the corresponding sentence template generated by the syntax model would be "$\knowledge$ was born on $\knowledge$." The knowledge-related token $\knowledge$ are then generated and filled in by the knowledge model.


% Since the contextual relationship is primarily determined by the knowledge components rather than the way the sentence is expressed, we assume that the sentence templates $T_{1:n}$ in the training data are independently sampled from $P_{\synModel}$.


\subsection{Hierarchical Syntax-Knowledge Data Model}
We propose that the Syntax–Knowledge model generates a sentence \( X \) according to the following hierarchical Bayesian framework
(see Figure~\ref{fig:kol}):
\begin{enumerate}
\item Sample the latent variables for the knowledge and syntax models:
    \[
    \knwModel \sim \PYP(\alpha,\beta,\basemeasure), \quad 
    \synModel = \{ \synModel^{(1)}, \synModel^{(2)}, \dots, \synModel^{(\numsyntax)} \}, \quad 
    \synModel^{(i)} \overset{\text{i.i.d.}}{\sim} \synprior(\synModel).
    \]
    
    \item Sample a knowledge element $\Knwelem$:
    \[
    \Knwelem \sim P(\Knwelem \mid \knwModel).
    \]
    
    \item Sample a syntactic element $\Synelem$ conditioned on the sampled knowledge element $\Knwelem$:
    \[
    \Synelem \sim P(\Synelem \mid \synModel, \Knwelem).
    \]
    
    \item The syntax encoder generates the sentence \( X \) from the syntax–knowledge pair \( (\Synelem, \Knwelem) \). 
    We assume that the syntactic element $\Synelem$ has already captured
    the randomness in the syntax, and there is a one-to-one mapping between
    the sentence \( X \) and the syntax–knowledge pair \( (\Synelem, \Knwelem) \).
    Note that we do not assume the one-to-one mapping to be explicitly known or learnable. For the purpose of our theoretical analysis, it is sufficient to assume the existence of such a mapping, which ensures identifiability of the observed sentence with respect to its underlying syntax and knowledge elements.
\end{enumerate}
     

By independently repeating steps 2, 3, 4 for \( N \) times, we obtain a corpus \( X_{1:N} \) consisting of $N$ i.i.d. sentences generated from the Syntax–Knowledge model.

We now illustrate how our data model generates data through a simple 
bioS dataset example. Note that our data model is not limited to such examples.
\begin{example}
A concrete example is the generation of the bioS dataset designed in
\cite{allen2024physics}. For example, a knowledge element $\Knwelem$  encodes the factual knowledge about an individual: \textit{[Person: Alice], [Event: Birth], [Date: January 1, 2000]}. 
%This knowledge element contains the factual knowledge without specifying how it should be expressed in natural language.
A syntax element $\Synelem$ corresponding to a sentence template.
Conditioned on the knowledge element $\Knwelem$, the syntax model chooses a syntax element $\Synelem$, such as ``[Subject] was born on [Date]''.  
By composing $(\Synelem, \Knwelem)$, the resulting sentence is: ``Alice was born on January 1, 2000.''
Alternatively, the syntax model may select a different syntax element, ``[Subject] entered the world on [Date]'', yielding the sentence: ``Alice entered the world on January 1, 2000.''
Note that the choice of syntax element is dependent on the knowledge element; for instance, if the knowledge element pertains to an animal, a different syntax element would be chosen.
\end{example}
\begin{nips}
    \begin{figure}
        \centering
        \includegraphics[width=0.5\linewidth]{data_model.pdf}
        \caption{An illustration of the hierarchical Syntax-Knowledge data model}
        \label{fig:data_model}
    \end{figure}
\end{nips}




% \begin{figure}[htbp]
%     \centering
%         \includegraphics[width=0.9\linewidth]{Arxiv/data model.pdf}
%         \caption{\zhixuan{some illustration}}
%         \label{fig:data model}
% \end{figure}



% The generative process of the set of data points $\{x_i\}_{i=1}^N$ is:
% \begin{enumerate}
%     \item Draw the base measure $\knwModel$ from the PYP:
%     \[
%     \knwModel \sim \PYP(\beta,\alpha , P_0)
%     \]
%     \item For each data point $x_i$:
%     \begin{enumerate}
%         \item Draw the cluster assignment:
%         \[
%         z_i \mid \knwModel \sim \knwModel
%         \]
%         \item Draw the observation:
%         \[
%         x_i \mid z_i \sim F(\phi_{z_i})
%         \]
%     \end{enumerate}
% \end{enumerate}
% where $F$ is the likelihood function and $\phi_{z_i}$ are the parameters associated with cluster $z_i$.

% Note that given the base measure \( \knwModel \), the data points \( x_i \) are conditionally i.i.d. according to the  mixture distribution \( x_i \mid \knwModel \sim \sum_{k=1}^{\infty} \pi_k \delta_{\phi_k} \).

% A key property of the Pitman-Yor Process is that the ranked cluster proportions follow a Poisson-Dirichlet distribution: \[(\pi_1, \pi_2, \ldots) \sim\text{PD}(\theta, d).\]

%  When $d = 0$, the Pitman-Yor process reduces to the Dirichlet process, and the Poisson-Dirichlet distribution becomes $\text{PD}(\theta, 0)$, which corresponds to the \knwModelEM (\knwModelriffiths-Engen-McCloskey) distribution:
% \[
% \pi_k = V_k \prod_{j=1}^{k-1} (1 - V_j), \quad V_k \sim \text{Beta}(1, \theta)
% \]
% When $d > 0$, the weights $\pi_k$ exhibit a power-law behavior, leading to a more flexible cluster size distribution.

% \jian{a strange sentence here..}

% \subsection{The Data Generation Process}


\section{Explaining Scaling Laws}
\label{sec: scaling law}
% This section is organized into three main parts.
% In \Cref{sec: estimation Bayesian Red}, we introduce additional assumptions under which the results in \Cref{sec:baye data law} can be further extended by providing an estimate of the Bayesian redundancy, accurate up to an $O(1)$ term. We note that the derivation of the estimation of the Bayesian redundancy is not fully rigorous, but rather consists of heuristic calculations intended to provide intuition.
% Then we address the distinction between the cumulative loss considered in \Cref{sec:baye data law} and the practically more relevant test loss. Under slightly modified assumptions, we derive corresponding results for the test loss scenario in \Cref{sec: test loss scaling}. Finally, we also provide the omitted details from the model scaling law part (\Cref{sec:model_scaling_law}) in the main text.
In \Cref{sec:baye data law}, we established a data scaling law under the Bayesian setting for our Syntax-Knowledge model, showing that the optimal Bayesian redundancy decreases according to a power law with respect to data size, with an exponent larger than \(-1\).  Given our earlier results from Section~\ref{sec:baye data law}, we know the syntax model is learned at a significantly faster rate (also observed empirically), and thus the  scaling law is primarily driven by the knowledge model. For simplicity, we ignore the syntax model in this section and focus exclusively on the knowledge model. In this section, we derive a similar power-law behavior under a slightly different set of assumptions. The key differentiating assumption is that knowledge is represented in a question-answering format, which ensures the identifiability of the mixture knowledge model. These assumptions greatly simplify the theoretical proof and allows us to derive both upper and lower bounds on redundancy, as long as the model satisfies Assumption~\ref{ass:data law} (without requiring it to be an optimal Bayesian predictor). The central insight of this section is that the empirically observed scaling laws primarily stem from the power-law-distributed structure in the data. As a result, the assumptions made here do not affect the validity of our main conclusion.
We also provide the omitted details from the model scaling law part (\Cref{sec:model_scaling_law}) in the main text.

% \subsection{Heuristic Estimation of Bayesian Redundancy}
% \label{sec: estimation Bayesian Red}
% Recall that the optimal Bayesian redundancy introduced in \Cref{sec:Bayesian Framework} is
% \begin{align*}
%     \inf_{\Model} \redundency_N(\Model,\datafamily) 
%    =\inf_{P_{\Model}}\int_{\datafamily} \mathcal{P}(\dataModel)\DKL(P^N_{\dataModel}(\Knwelem,\Synelem)||P_{\Model}(\Knwelem_{1:N},\Synelem_{1:N})) \dd\dataModel
% \end{align*}
% We first apply the chain rule of KL divergence to decompose the Bayesian redundancy into knowledge-related and syntax-related terms.
% \begin{align}
%     &\inf_{P_{\Model}}\int_{\datafamily} \mathcal{P}(\dataModel)\DKL(P^N_{\dataModel}(\Knwelem,\Synelem)||P_{\Model}(\Knwelem_{1:N},\Synelem_{1:N})) \dd\dataModel \nonumber
%    \\=&\inf_{P_{\Model}}\int_{\datafamily} \mathcal{P}(\dataModel)\left(\DKL(P^N_{\knwModel}(\Knwelem)||P_{\Model}(\Knwelem_{1:N}))+\DKL(P^N_{\synModel}(\Synelem|\Knwelem)||P_{\Model}(\Synelem_{1:N}|\Knwelem_{1:N}))\right) \dd\dataModel \nonumber
%    \\=&\underbrace{\inf_{P_{\Model}(\Knwelem_{1:N})}\int_{\knwfamily} \mathcal{P}(\knwModel)\DKL\left(P^N_{\knwModel}(\Knwelem)||P_{\Model}(\Knwelem_{1:N})\right)\dd\knwModel}_{\text{knowledge-related term}} \nonumber
%    \\+&\underbrace{\inf_{P_{\Model}(\Synelem_{1:N}|\Knwelem_{1:N})}\int_{\synfamily}\mathcal{P}(\synModel)\DKL\left(P^N_{\synModel}(\Synelem|\Knwelem)||P_{\Model}(\Synelem_{1:N}|\Knwelem_{1:N})\right) \dd\synModel}_{\text{syntax-related term}} 
%    \label{eq:knw syn decompose}
% \end{align}
% In our data model, the generation of syntax elements conditioned on the knowledge elements corresponds to a family of distributions over a space of dimension $\numsyntax \cdot \syndimension$. Then, according to Theorem 2.1 in \citep{clarke1990information}, under certain continuity and regularity conditions, the syntax-related term in \Cref{eq:knw syn decompose} satisfies:
% \begin{align}
%     \inf_{P_{\Model}(\Synelem_{1:N}|\Knwelem_{1:N})}\int_{\synfamily}\mathcal{P}(\synModel)\DKL\left(P^N_{\synModel}(\Synelem|\Knwelem)||P_{\Model}(\Synelem_{1:N}|\Knwelem_{1:N})\right) \dd\synModel \approx \dfrac{\numsyntax \syndimension}{2}\log \dfrac{N}{2\pi}.
%     \label{eq:approx syn}
% \end{align}
% where the approximation symbol $\approx$ indicates that additive terms independent of $N$ have been omitted.

% However, for the knowledge-related term, since our knowledge model is a mixture model, the identifiability condition (i.e., $\theta_1 \ne \theta_2 \rightarrow
% P_{\theta_1} \ne P_{\theta_2}$) typically does not hold. More specifically, the non-identifiability of mixture model parameters can be viewed as a form of redundancy in the parameter space—some dimensions are not strictly necessary for uniquely specifying the distribution. Consequently, we can only provide an upper bound on the redundancy, while deriving a meaningful lower bound is generally difficult.
% % Due to the joint convexity of the KL divergence (\Cref{lem:mixture KL} used in the proof of \Cref{thm:data_scaling_Mutual_Information} is a generalization of the joint convexity), it is generally difficult to obtain an upper bound for the Bayesian redundancy of a mixture knowledge model. 
% In this section, we consider a simplified mixture knowledge model in which the knowledge elements $\Knwelem$ generated by different knowledge clusters $\phi_i$ are all distinct. We denote the support of $P_{\phi_i}$ as $\knwset_i$. Under this assumption, the mixture model is identifiability. Thus the knowledge-related term in \Cref{eq:knw syn decompose} can be decomposed as follows:
% \begin{align*}
%     &\inf_{P_{\Model}(\Knwelem_{1:N})}\int_{\knwfamily} \mathcal{P}(\knwModel)\DKL\left(P^N_{\knwModel}(\Knwelem)||P_{\Model}(\Knwelem_{1:N})\right)\dd\knwModel
%     \\=& \inf_{P_{\Model}(\Knwelem_{1:N})}\int_{\knwfamily} \mathcal{P}(\knwModel)\sum\limits_{i=1}^N\DKL\left(P_{\knwModel}(\Knwelem)||P_{\Model}(\Knwelem_{i})|\Knwelem_{1:i-1}\right)\dd\knwModel
%     \\=& \inf_{P_{\Model}(\Knwelem_{1:N})}\int_{\knwfamily} \mathcal{P}(\knwModel)\sum\limits_{i=1}^N\sum\limits_{j=1}^{\infty}p_i\DKL\left(P_{\phi_i}(\Knwelem)||P_{\Model}(\Knwelem_{i})|\Knwelem_{1:i-1},\Knwelem_i\in \knwset_i\right)\dd\knwModel
%     \\=& \underbrace{\inf_{P_{\Model}(\Knwelem_{1:N})}\int_{\knwfamily} \mathcal{P}(\knwModel)\sum\limits_{i=1}^N\sum\limits_{j=1}^{m}p_i\DKL\left(P_{\phi_i}(\Knwelem)||P_{\Model}(\Knwelem_{i})|\Knwelem_{1:i-1},\Knwelem_i\in \knwset_i\right)\dd\knwModel}_{\text{frequently selected term}}
%     \\+&\underbrace{\inf_{P_{\Model}(\Knwelem_{1:N})}\int_{\knwfamily} \mathcal{P}(\knwModel)\sum\limits_{i=1}^N\sum\limits_{j=m+1}^{\infty}p_i\DKL\left(P_{\phi_i}(\Knwelem)||P_{\Model}(\Knwelem_{i})|\Knwelem_{1:i-1},\Knwelem_i\in \knwset_i\right)\dd\knwModel}_{\text{rarely selected term}}
% \end{align*}
% where $m = N^{\alpha}$ is the truncation point. The first term corresponds to the redundancy of the clusters $\phi_j$ with $j \leq m$, while the second term accounts for the redundancy from clusters with $j > m$. Note that the choice of $m$ is motivated by the conclusion of \Cref{lem:PYCRP fre}, which allows us to divide the clusters into two groups: frequently selected and rarely selected. 

% Since the clusters in the second term are rarely selected, the KL divergence in each integral is likely to be of constant order(approximately equal to $c$) with high probability. Then according to \Cref{lem:L_j}, we estimate the rarely selected term as follows:
% \begin{align}
%     &\inf_{P_{\Model}(\Knwelem_{1:N})}\int_{\knwfamily} \mathcal{P}(\knwModel)\sum\limits_{i=1}^N\sum\limits_{j=m+1}^{\infty}p_i\DKL\left(P_{\phi_i}(\Knwelem)||P_{\Model}(\Knwelem_{i})|\Knwelem_{1:i-1},\Knwelem_i\in \knwset_i\right)\dd\knwModel \nonumber
%     \\ \approx & \inf_{P_{\Model}(\Knwelem_{1:N})}\int_{\knwfamily} \mathcal{P}(\knwModel)\sum\limits_{i=1}^N\sum\limits_{j=m+1}^{\infty}p_i \ c\  \dd\knwModel = \Theta(N^{\alpha})
%     \label{eq:approx rarely}
% \end{align}
% On the other hand, because the clusters in the second group are rarely selected, the first term can be approximated as the redundancy of an identifiable distribution over an $m \cdot \knwdimension$-dimensional space. According to Theorem 2.1 in \citep{clarke1990information}, the corresponding estimation of the frequently selected term is given as follows:
% \begin{align}
%     \inf_{P_{\Model}(\Knwelem_{1:N})}\int_{\knwfamily} \mathcal{P}(\knwModel)\sum\limits_{i=1}^N\sum\limits_{j=1}^{m}p_i\DKL\left(P_{\phi_i}(\Knwelem)||P_{\Model}(\Knwelem_{i})|\Knwelem_{1:i-1},\Knwelem_i\in \knwset_i\right)\dd\knwModel \approx \dfrac{\knwdimension N^{\alpha}}{2}\log \dfrac{N}{2\pi}
%     \label{eq:approx fre}
% \end{align}
% Combining \eqref{eq:approx rarely} and \eqref{eq:approx fre}, we know that
% \begin{equation}
%     \inf_{P_{\Model}(\Knwelem_{1:N})}\int_{\knwfamily} \mathcal{P}(\knwModel)\DKL\left(P^N_{\knwModel}(\Knwelem)||P_{\Model}(\Knwelem_{1:N})\right)\dd\knwModel\approx \dfrac{\knwdimension N^{\alpha}}{2}\log \dfrac{N}{2\pi} 
%     \label{eq:approx knw}
% \end{equation}
    

% where the approximation symbol $\approx$ indicates that we omit terms of order smaller than $N^{\alpha} \log N$ as well as additive constants independent of $N$.

% The heuristic estimates \eqref{eq:approx syn} and \eqref{eq:approx knw} in this section match the order of the upper bound provided in the main text by \Cref{thm:data_scaling_Mutual_Information}. This suggests that, although only an upper bound is formally established in the main text, it is likely to be tight in terms of order. Under the heuristic assumptions made in this section, we obtain the same order estimation, indicating that the scaling behavior of the Bayesian redundancy is well captured by our theoretical upper bound in \Cref{thm:data_scaling_Mutual_Information}.





% \subsection{More on Scaling Laws}
% \label{sec: test loss scaling}
% In \Cref{sec:baye data law} and \Cref{sec: estimation Bayesian Red}, we established a data scaling law under the Bayesian setting for our Syntax-Knowledge model, showing that the optimal Bayesian redundancy decreases according to a power law with respect to data size, with an exponent larger than \(-1\).  Given our earlier results from Section~\ref{sec:baye data law} and \Cref{sec: estimation Bayesian Red}, we know the syntax model is learned at a significantly faster rate (also observed empirically), and thus the scaling law is primarily driven by the knowledge model. For simplicity, we ignore the syntax model in this section and focus exclusively on the knowledge model. In this section, we derive a similar power-law behavior under a slightly different set of assumptions. These assumptions greatly simplify the theoretical proof and allows us to derive both upper and lower bounds on redundancy, as long as the model satisfies Assumption~\ref{ass:data law} (without requiring it to be an optimal Bayesian predictor). The central insight of this section is that the empirically observed scaling laws primarily stem from the power-law-distributed structure in the data. As a result, the assumptions made here do not affect the validity of our main conclusion.






% The knowledge model \( \knwModel = \sum_{i=1}^{\infty} p_i \delta_{\phi_i} \) remains an infinite mixture model, but with fixed mixing probabilities \( p_i \). In particular, we assume that the mixing probabilities follow a power-law distribution of the form
% \[
% p_i = \frac{1}{\zeta(1/\alpha) \, i^{1/\alpha}},
% \]
% where \( \zeta(1/\alpha) = \sum_{i=1}^{\infty} \frac{1}{i^{1/\alpha}} \) is the Riemann zeta function. The power-law exponent \( 1/\alpha \) is chosen to match the asymptotic behavior of the mixing weights in the Pitman–Yor Process \( \mathrm{PYP}(\alpha, \beta, \basemeasure) \).
% For cluster parameter $\phi_i$, we still assume that $\phi_i \overset{\text{i.i.d.}}{\sim} \basemeasure $.

We continue modeling the knowledge model \( \knwModel = \sum_{i=1}^{\infty} p_i \delta_{\phi_i} \) as an infinite mixture model but now make the following assumption on the mixing probabilities \( p_i \). 

\begin{assumption}
\label{appass:powerlaw}
For the mixture model \( \knwModel = \sum_{i=1}^{\infty} p_i \delta_{\phi_i} \), the mixing probabilities \( p_i \) 
follow a power-law distribution:
$
p_i = \zeta(1/\alpha)^{-1} \, i^{-1/\alpha},
$
where \( \zeta(1/\alpha) = \sum_{i=1}^{\infty} i^{-1/\alpha} \) is the Riemann zeta function.

Note that the choice of exponent \(1/\alpha\) aligns with the asymptotic behavior of mixing weights in the Pitman–Yor Process \( \mathrm{PYP}(\alpha, \beta, \basemeasure) \) (see Lemma~\ref{lem:PYCRP fre}). 

\end{assumption}



In this section, we represent knowledge as question-answer pairs. We assume knowledge is given by pairs \( (\question, \answer) \), where \( \question \) denotes the question description, and \( \answer \) is the answer of the question. 
% Each knowledge cluster corresponds to a question description set \( \questionset_i \) and is associated with a conditional distribution \( P_{\phi_i}(\answer \mid \question_i) \) over the generated answers \( \answer \). 
Each knowledge cluster corresponds to a set of question descriptions \( \questionset_i \), and generates knowledge pairs \( (\question, \answer) \sim P_{\phi_i} \) where \( \question \in \questionset_i \).  
We assume that the sets of question descriptions corresponding to different knowledge clusters are pairwise disjoint, i.e., \( \questionset_i \cap \questionset_j = \emptyset \) for \( i \ne j \). 
% \zhixuan{Do you need to write this as begin\{assumption\}?}
% \jian{No.}



Here, we are only concerned with whether the model can provide the corresponding answer given a specific question. Consequently, the redundancy of the model with respect to the \( i \)-th knowledge cluster is defined as follows:


\begin{definition}
The redundancy of the model \( \Model \) associated with the \( i \)-th knowledge cluster and conditioned on knowledge being drawn from this cluster, is defined as
\[
\ferror^{(i)} := \mathbb{E}_{q \sim P_{\phi_i}} \left[ \DKL\left( P_{\phi_i}(\answer \mid \question=q) \,\Vert\, P_{\Model}(\answer \mid \question=q) \right) \right],
\]
where \( P_{\phi_i} \) denotes the data distribution induced by the cluster-specific parameter \( \phi_i \), and \( P_{\Model} \) is the predictive distribution defined by the model \( \Model \).



The total expected redundancy of the model \( \Model \) under the prior over knowledge clusters is then given by:
% \begin{align*}
% \testerror 
% &:= \mathbb{E}_{\knwModel \sim \mathrm{PYP}(\alpha,\beta,\basemeasure)}\left[\sum_{i=1}^{\infty} p_i \ferror^{(i)}(\Model)\right] \\
% &= \mathbb{E}_{p \sim \mathrm{PYCRP}(\alpha,\beta)}\left[\sum_{i=1}^{\infty} p_i \mathbb{E}_{\phi_i \sim \basemeasure}\left[\ferror^{(i)}(\Model)\right] \right],
% \end{align*}
% where \( p = (p_1, p_2, \ldots) \) denotes the mixing weights of $\knwModel$.
\begin{align*}
\testerror 
&:= \sum_{i=1}^{\infty} p_i \mathbb{E}_{\phi_i \sim \basemeasure}\left[\ferror^{(i)}\right].
\end{align*}








% In particular, we derive a model scaling law by solving a constrained optimization problem that minimizes the Bayesian redundancy under model-size constraints.

% Under certain assumptions, we formulated the optimal loss constrained by model size as a constrained optimization problem. Then we proved that the solution to this optimization problem follows a power-law relationship with model size. 


% We no longer assume a specific form of knowledge
% % \jian{what form?? confusing} 
% but focus on its impact on error through two key aspects: 
% \begin{itemize}
%     \item The probability of occurrence of a specific cluster of knowledge;
%     \item The error when the knowledge occurs during testing is belongs to a specific cluster.
% \end{itemize}
% \jian{no introduction of cluster yet...}

 % For the second aspect, we use the following definition

% Since the PYCRP process we discussed in \Cref{sec:PYCRP} is exchangeable, according to de Finetti’s theorem, we know that there exists a latent variable $\phi$ such that conditioned on $\phi$, knowledge pieces are i.i.d. Furthermore, in the PYCRP process, the number of customers at each table follows a power-law distribution.

% Inspired by Zipf's Law \citep{zipf2016human}, we assume that knowledge is independently and identically sampled from a power-law distribution.
% \begin{assumption}
%     The probability of knowledge occurrence follows a power law:
%     $
%     p_k \propto \dfrac{1}{k^\alpha}, \quad p_k = \dfrac{1}{\zeta(\alpha)k^{\alpha}},$
%     where $\alpha>1$ is a constant, $\zeta(\alpha) = \sum\limits_{n=1}^{\infty} \frac{1}{n^{\alpha}}$ is the Riemann zeta function, and $p_k$ denotes the occurrence probability of the $k$-th knowledge.
% \end{assumption}


% First, the index $k$ of the knowledge: Different pieces of knowledge have different levels of learning difficulty. This includes the inherent difficulty of the knowledge itself as well as the difficulty arising from potential forgetting due to differences in knowledge occurrence frequency. Second, the occurrence count $t$ of the knowledge in the training data. We believe that data size influences error through this factor.
% Third, the amount of storage $m$ allocated for the knowledge in the model. We believe that model size influences error through this factor.

% Let \( \ferror^{(i)} \) denote the redundancy associated with the \( i \)-th knowledge cluster. Since this section focuses on the data scaling law, the size of the training data primarily influences \( \ferror^{(i)} \) through the number of samples \( t_i \) drawn from the \( i_{th} \) knowledge cluster. Hence, we express the redundancy \( \ferror^{(i)} \) as a function of the sample count \( t_i \). The formal definition is:



% Then the expected redundancy $\testdataerror$ is:
% \begin{align*}
%     \testdataerror:=\mathbb{E}_{\knwModel,i,t_i}\left[\dataerror^{(i)}(t_i) \right].
%     % = \sum\limits_{k=1}^{\infty} p_k \mathbb{E}_{t_k}[\ferror(k,t_k,m_k)].
% \end{align*}


\end{definition}

\subsection{Data Scaling Laws}
\label{sec:data scaling}
% In this section, we analyze the data scaling law. We consider the case where the model size is sufficiently large. 

% We continue to adopt the previously mentioned $\PYMM$ as our knowledge model, i.e. the knowledge model $\knwModel=\sum_{i=1}^{\infty}p_i\delta_{\phi_i}$. 

% Denote $\ferror^{(i)}$ as the redundancy of $i_{th}$ cluster of knowledge. Since this section focuses on the data scaling law, the size of the training data primarily affects redundancy $\ferror^{(i)}$ through the number of samples drawn from $i_{th}$ cluster knowledge. Therefore, we express redundancy $\ferror^{(i)}$ as a function of the sampling account $t_i$.

% Let \( \ferror^{(i)} \) denote the redundancy associated with the \( i \)-th knowledge cluster. Since this section focuses on the data scaling law, the size of the training data primarily influences \( \ferror^{(i)} \) through the number of samples \( t_i \) drawn from the \( i_{th} \) knowledge cluster. Hence, we express the redundancy \( \ferror^{(i)} \) as a function of the sample count \( t_i \). The formal definition is:

% We now discuss the factors that may influence $\ferror^{(i)}$.

% \begin{itemize}
%     \item The count of occurrences of this cluster of knowledge in the training data.  The size of the training data mainly affects the test redundancy through this factor.
%     \item The amount of storage $m$ allocated for this cluster of knowledge. The size of model mainly affects the test redundancy through this factor.
% \end{itemize}

% Accordingly, we use $\ferror^{(i)}(t)$ to represent the redundancy of $i_{th}$ cluster of knowledge during testing. 
% The formal definition is:
% \begin{definition}
%      The redundancy associated with the \( i \)-th knowledge cluster, given that knowledge occurs from this cluster, is defined as follows:
% \[
% \dataerror^{(i)}(t_i) :=  \DKL\left( P_{\phi_i} \,\Vert\, P_{M} \right),
% \]
% where \( P_{\phi_i} \) is the data distribution corresponding to parameter \( \phi_i \), and \( P_{M} \) is the predictive distribution induced by the model \( M \).

% Then the expected redundancy $\testdataerror$ is:
% \begin{align*}
%     \testdataerror:=\mathbb{E}_{\knwModel,i,t_i}\left[\dataerror^{(i)}(t_i) \right].
%     % = \sum\limits_{k=1}^{\infty} p_k \mathbb{E}_{t_k}[\ferror(k,t_k,m_k)].
% \end{align*}


% \end{definition}






% Accordingly, we use $\ferror^{(i)}(t,m)$ to represent the redundancy of $i_{th}$ cluster of knowledge during testing. The formal definition is:
% \begin{definition}
%      The redundancy when the knowledge occurs is belongs to $i_{th}$ cluster is defined as following:
%     \[\ferror^{(i)}(t_i,m_i):=\mathbb{E}_{\phi_i}\left[\DKL (P_{\phi_i}\|P_{M})\right],\] where $\M$ is the model which is derived based on $t_i$ times of observations from $\phi_i$ and stored using $m_i$ bits of memory. 

%     Then the expected test redundancy $\testerror$ is:
% \begin{align*}
%     \testerror:=\mathbb{E}_{\knwModel,i,t_i}\left[\ferror^{(i)}(t_i,m_i) \right].
%     % = \sum\limits_{k=1}^{\infty} p_k \mathbb{E}_{t_k}[\ferror(k,t_k,m_k)].
% \end{align*}
% \end{definition}



% Here, we highlight the key differences in the setup between this section and the \Cref{sec:Bayesian}:
% \begin{itemize}
%     % \item In \Cref{sec:Bayesian}, we focus on the redundancy of a learning algorithm. Here we consider the redundancy of a given model ????. 
%     \item In \Cref{sec:Bayesian}, we focus on the optimal redundancy. Here we consider the redundancy of a model which may not be optimal but satisfies certain assumptions.
% \end{itemize}
In this section, we derive the data scaling law under the assumption that there is no model capacity constraint. There are two key differences between this section and \Cref{sec:baye data law}. First, we analyze the redundancy of a potentially suboptimal model under certain assumptions, whereas \Cref{sec:baye data law} focuses on the optimal case. Second, we study the test redundancy, in contrast to the Bayesian redundancy considered in \Cref{sec:baye data law}, which corresponds to the averaged cumulative redundancy.
Moreover, unlike \Cref{sec:baye data law}, which provides only an upper bound, this section establishes both upper and lower bounds.
The effect of dataset size on \( \ferror^{(i)} \) primarily arises from the number of occurrences of the question description \( q\in \questionset_i \) in the training data. Accordingly, we express \( \ferror^{(i)} \) as $\dataerror(t_i)$, where $t_i$ is the number of such occurrences
in the training data. 
The expected redundancy under data size constraint is then given by
\[ \testdataerror(N) = \mathbb{E}_{t_i}\left[\sum_{i=1}^{\infty} p_i \mathbb{E}_{\phi_i \sim \basemeasure}\left[\dataerror^{(i)}(t_i)\right] \right],\]
where \( t_i \) denotes a random variable representing the number of occurrences of the \( i \)-th knowledge cluster in a training dataset of size \( N \).






Next, we present the assumptions on \( \dataerror^{(i)} \) and derive the upper and lower bound for \( \testdataerror \) under these assumptions.

\begin{assumption}
\label{ass:data law}
(\uppercase\expandafter{\romannumeral 1})
The expected redundancy of unseen cluster of knowledge (i.e., $t_i=0$) exceeds a certain constant $c_1$, i.e., $\mathbb{E}_{\phi_i\sim \basemeasure}\left[\dataerror^{(i)}(0)\right]>c_1$ for all $k\in \mathcal{N}^+$.

(\uppercase\expandafter{\romannumeral 2})
For all cluster of knowledge, there exists a constant $c_2$ such that 
 $\mathbb{E}_{\phi_i\sim \basemeasure}\left[\dataerror^{(i)}(t)\right]\le \dfrac{c_2}{t+1}$ for all $i,t\in \mathcal{N}^+$.
\end{assumption}

Assumption (\uppercase\expandafter{\romannumeral 1}) posits that the model exhibits constant redundancy when encountering unseen knowledge, which we consider reasonable. We now turn to the validity of Assumption (\uppercase\expandafter{\romannumeral 2}). 
During language model training, the objective is to minimize the negative log-likelihood, which is equivalent to maximizing the likelihood function. When the model is viewed as a maximum likelihood estimator (MLE), the asymptotic normality of the MLE implies that the redundancy, in the asymptotic regime, is of the same order as specified in Assumption (\uppercase\expandafter{\romannumeral 2}). 
\begin{lemma}[Asymptotic Normality of MLE]
Let \( X_1, X_2, \dots, X_n \) be $n$ i.i.d. samples drawn from a distribution \( P_{\theta^*} \), and let \( \hat{\theta}_n = \arg \max_{\theta} \ell(\theta) \) denote the maximum likelihood estimator. Then, as \( n \to \infty \), the MLE satisfies
\[
\sqrt{n}(\hat{\theta}_n - \theta^*) \xrightarrow{d} \mathcal{N}(0, J(\theta^*)^{-1}),
\]
where \( J(\theta^*) \) is the Fisher information matrix evaluated at \( \theta^* \). 
\end{lemma}
Combining this conclusion with the result in \Cref{lem:KL and MSE}, we conclude that, under the asymptotic regime, the redundancy of the MLE matches the order assumed in Assumption (\uppercase\expandafter{\romannumeral 2}). We assume that this order holds in all cases. Note that the assumption is stated with $1/(t+1)$ instead of $1/t$ to include the case where $t = 0$.





% See details in \Cref{}.

% \zhixuan{Justify the assumptions}





\begin{theorem}
\label{thm: data scaling lower bound}
Under (\uppercase\expandafter{\romannumeral 1}) in \Cref{ass:data law}, the total expected redundancy $\testdataerror$ satisfies
    \begin{align*}
    \testdataerror(N) = \Omega(N^{-1+\alpha}).
\end{align*}
Under (\uppercase\expandafter{\romannumeral 2}) in \Cref{ass:data law}, the total expected redundancy $\testdataerror$ satisfies
    \begin{align*}
    \testdataerror(N) = O(N^{-1+\alpha}).
\end{align*}
\end{theorem}
\begin{proof}[Proof of \Cref{thm: data scaling lower bound}]
% By \Cref{thm:PYP rank fre}, we know that there exists $N_0$ such that for all $n\ge N_0$, 
% $\dfrac{S_{\alpha,\beta}}{2n^{1/\alpha}}\le p_n\le \dfrac{2S_{\alpha,\beta}}{n^{1/\alpha}}$ almost surely. Then we have
Denote $f_i(t_i)=\mathbb{E}_{\phi_i\sim \basemeasure}[\dataerror^{(i)}(t_i)]$.

We begin by proving the first part of the theorem.
    \begin{align*}
        \testdataerror(N)&=\sum\limits_{k=1}^{\infty} p_k \mathbb{E}_{t_k}[f_k(t_k)]
        \\&\ge \sum\limits_{k> N^{\alpha}} p_k\mathbb{E}_{t_k}[f_k(t_k)] 
        \\& \overset{(1)}{=} \sum\limits_{k> N^{\alpha}}\sum\limits_{l=0}^N p_k^{l+1} (1-p_k)^{N-l}\binom{N}{l}f_k(l) \\
        &\overset{(2)}{\ge} \sum\limits_{k> N^{\alpha}} p_k (1-p_k)^{N} f_k(0)\\
        &\overset{(3)}{\ge} \dfrac{c_1}{\zeta(1/\alpha)}\sum\limits_{k> N^{\alpha}} \dfrac{1}{k^{1/\alpha}} \left(1-\dfrac{1}{\zeta(1/\alpha)N}\right)^{N} \\
        &\ge \dfrac{c_1}{2e\zeta(1/\alpha)}\sum\limits_{k> N^{\alpha}} \dfrac{1}{k^{1/\alpha}} \\
        &=  \Omega(N^{-1+\alpha}).
    \end{align*}
(1) is obtained by expanding the expectation; (2) corresponds to taking the term with $l = 0$; and (3) is derived by substituting the value of $p_k=\dfrac{1}{k^{1/\alpha}\zeta(1/\alpha)}$.

% \end{proof}


% \begin{theorem}[Restatement of \Cref{thm: data scaling upper bound}]
% Assume there exists a constant $C$ such that 
%  $(t+1)\dataerror(k,t)\le C$ for all $k,t\in \mathcal{N}^+$. Then the test error $\testdataerror$ satisfies
%     \begin{align*}
%     \testdataerror(N)\le O(N^{-1+\alpha}).
% \end{align*}
% \end{theorem}
% \begin{proof}[Proof of \Cref{thm: data scaling upper bound}]
% By \Cref{thm:PYP rank fre}, we know that there exists $N_0$ such that for all $n\ge N_0$, 
    
% $\dfrac{S_{\alpha,\beta}}{2n^{1/\alpha}}\le p_n\le \dfrac{2S_{\alpha,\beta}}{n^{1/\alpha}}$ almost surely.

    Next, we prove the second part. For all $k\in \mathcal{N}^+$, we know that 
    % \begin{align*}
    %     \sum\limits_{k\le (2S_{\alpha,\beta}N)^{\alpha}}\mathbb{E}_{\knwModel}\left[p_kE_{t_k}[\dataerror(k,t_k)]\right]  = O(N^{-1+\alpha}).
    % \end{align*}

    
    % Let the training data be $X$, and the model trained on $X$ has a test loss for knowledge $Y$ denoted by $g_X(Y)$.
    % \begin{align*}
    %     E_{X,Y}[g_X(Y)] &= E_{X}\left[\sum\limits_{k=1}^{\infty}p_kf(t_k)\right] \\
    %     &= \sum\limits_{k=1}^{\infty}p_kE_{X}[f(t_k)].
    % \end{align*}
    % Here, $t_k$ is the number of times the $k$-th knowledge appears in the training data.
    \begin{align*}
        p_k\mathbb{E}_{t_k}[f_k(t_k)] &\overset{(1)}{=} p_k\sum\limits_{l=0}^N p_k^l (1-p_k)^{N-l}\binom{N}{l}f_k(l) \\
        &= \sum\limits_{l=0}^N p_k^{l+1} (1-p_k)^{N-l}\dfrac{l+1}{N+1}\binom{N+1}{l+1}f_k(l) \\
        &\le \dfrac{\max\limits_{l} (l+1)f_k(l)}{N+1}\sum\limits_{l=0}^N p_k^{l+1} (1-p_k)^{N-l}\binom{N+1}{l+1} \\
        &\overset{(2)}{\le} \dfrac{\max\limits_{l} (l+1)f_k(l)}{N+1}\\ &\le \dfrac{c_2}{N+1}.
    \end{align*}
    (1) is obtained by expanding the expectation, and (2) is derived by utilizing 
    \\
    $\sum\limits_{l=0}^N p_k^{l+1} (1-p_k)^{N-l}\binom{N+1}{l+1}\le \sum\limits_{l=-1}^N p_k^{l+1} (1-p_k)^{N-l}\binom{N+1}{l+1}=1.$

    Therefore,
    \begin{align*}
        \sum\limits_{k\le N^{\alpha}}p_k\mathbb{E}_{t_k}[f_k(t_k)] \le  \dfrac{c_2}{N+1} N^{\alpha} =O(N^{-1+\alpha}).
    \end{align*}
    For $k>N^{\alpha}$, we have 
    \begin{align*}
        \sum\limits_{k> N^{\alpha}}p_k\mathbb{E}_{t_k}[f_k(t_k)] &= \sum\limits_{k> N^{\alpha}}\sum\limits_{l=0}^N p_k^{l+1} (1-p_k)^{N-l}\binom{N}{l}f_k(l) \\
        &\overset{(1)}{\le}\ \sum\limits_{k> N^{\alpha}}\sum\limits_{l=0}^N \left(\dfrac{1}{\zeta(1/\alpha)k^{1/\alpha}}\right)^{l+1} \binom{N}{l}f_k(l) \\
        &\overset{(2)}{\le} \sum\limits_{k> N^{\alpha}}\sum\limits_{l=0}^N \left(\dfrac{1}{k^{1/\alpha}}\right)\left(\dfrac{1}{N}\right)^l \binom{N}{l}f_k(l) \\
        &\overset{(3)}{\le}\ \sum\limits_{k> N^{\alpha}}\sum\limits_{l=0}^N \left(\dfrac{1}{k^{1/\alpha}}\right)\dfrac{f_k(l)}{l!}  
        \\&= \sum\limits_{l=0}^N \dfrac{f_k(l)}{l!} \sum\limits_{k> N^{\alpha}}\dfrac{1}{k^{1/\alpha}}
        \\ &\le c_2\sum\limits_{l=0}^N \dfrac{1}{(l+1)!} \sum\limits_{k> N^{1/\alpha}}\dfrac{1}{k^{1/\alpha}}
        \\ &\le c_2e\sum\limits_{k> N^{\alpha}}\dfrac{1}{k^{1/\alpha}} 
        = O(N^{-1+\alpha}).
    \end{align*}
    (1) is obtained by substituting $p_k=\dfrac{1}{k^{1/\alpha}\zeta(1/\alpha)}$ and using the inequality $(1 - p_k)^{N-l} \le 1$; (2) follows from the fact that $\zeta(1/\alpha) > 1$ and $k > N^{\alpha}$; (3) uses the bound $\binom{N}{l} \le \dfrac{N^l}{l!}$.

    Therefore, we can see that
    \begin{align*}
   \testdataerror(N)=\mathbb{E}_{k,t_k}
   \left[\sum\limits_{k=1}^{\infty}p_kf_k(t_k) \right]\le O(N^{-1+\alpha}).
    \end{align*}
This completes the proof.   
\end{proof}


% \begin{theorem}
% \label{thm: data scaling upper bound}
%     Assume  Then the test error $\testdataerror$ satisfies
%     \begin{align*}
%     \testdataerror(N)\le O(N^{-1+1/\alpha}).
% \end{align*}
% \end{theorem}


% \begin{remark}
%    This assumption corresponds to the empirical Zipf's law for word frequencies. We might also be able to provide some theoretical justification, such as the idea in "Producing Power-Law Distributions and Damping Word Frequencies with Two-Stage Language Models." We intend to demonstrate that this power law is the primary cause for the power law in the loss.
% \end{remark}

% \begin{assumption}
%     We assume that a knowledge appearing $l$ times in the training data will result in a knowledge error of $\ferror(l)$ during testing, and we also assume that there exists a constant $C$ such that for any positive integer $n$,
%     \[
%     \max\limits_{l} (l+1)f(l) \le C.
%     \]
% \end{assumption}

% \begin{remark}
%     This assumption is quite strong.

%     It implicitly assumes that the learning difficulty for all knowledge is the same and that the model size is infinite. We might be able to improve upon this assumption.
% \end{remark}


% \begin{theorem}
%     Suppose our training data contains $n$ sentences, then the expected test loss is $\Theta(n^{-1+1/\alpha})$.
% \end{theorem}





% \subsubsection{Knowledge Combination}
% 如果我们考虑更复杂一些的知识组合，也会有类似的结论。考虑如下data model：

% \begin{itemize}
%     \item 每个sentence可能包含多个knowledge，但是用一个pair $(x,y)$来表示。
%     \item 每个sentence中出现的各个knowledge的概率是独立的，也就是说当一个sentence包含k个knowledge时，出现$i_1,i_2,\cdots,i_k$这些knowledge各$j_1,j_2,\cdots,j_k$次的概率是$p_{i_1}^{j_1}\cdots p_{i_k}^{j_k}$
%     \item sentence包含的knowledge数量是有上界的。
%     \item $p_k$仍然follow power law
%     \[p_k\propto \dfrac{1}{k^\alpha},p_k=\dfrac{1}{\zeta(\alpha)k^{\alpha}}.\]
%     \item sentence中仅包含一个knowledge的概率是$\Theta(1)$的
% \end{itemize}
% \begin{theorem}
%     假设我们的训练数据总共包含n个sentences，则从期望上测试loss是$\Theta(n^{-1+1/\alpha})$
% \end{theorem}
% \begin{proof}
%     这里我们只需要证明如果test sentence包含超过一个knowledge，loss从期望上仍然是不超过$n^{-1+1/\alpha}$的。

%     假设test sentence包含了$k$个knowledge

%     记训练数据为$X$,在X上训练的模型对于Knowledge Y的Test Loss为$g_X(Y)$.
%     \begin{align*}
%         E_{X,Y}[g_X(Y)]&=E_{X}[\sum\limits_{\alpha_1+\cdots+\alpha_j=k}^{\infty}p_{i_1}^{\alpha_1}\cdots p_{i_j}^{\alpha_j} f(t_{i_1},\cdots,t_{i_k})]
%         \\&=\sum\limits_{\alpha_1+\cdots+\alpha_j=k}^{\infty}p_{i_1}^{\alpha_1}\cdots p_{i_j}^{\alpha_j}E_{X}[\ferror(t_{i_1},\cdots,t_{i_k})].
%     \end{align*}
%     其中$t_k$第k条knowledge在训练数据中出现的次数。
%     \begin{align*}
%         \ &p_{i_1}^{\alpha_1}\cdots p_{i_j}^{\alpha_j}E_{X}[\ferror(t_{i_1},\cdots,t_{i_k})]
%         \\ &=p_{i_1}^{\alpha_1}\cdots p_{i_j}^{\alpha_j}\sum\limits_{l_1,\cdots,l_j}^n p_{i_1}^{l_1}\cdots p_{i_j}^{l_j} (1-p_{i_1}-\cdots-p_{i_k})^{n-l_1-\cdots-l_j}\dfrac{n!}{l_1!\cdots l_j!}f(l_1,\cdots,l_j)
%         \\&=\sum\limits_{l_1,\cdots,l_j} p_{i_1}^{\alpha_1+l_1}\cdots p_{i_j}^{\alpha_j+l_j} (1-p_{i_1}-\cdots-p_{i_j})^{n-l_1-\cdots-l_j}\dfrac{(n+k)!}{(l_1+\alpha_1)!\cdots(l_j+\alpha_j)!}\dfrac{(l_1+1)\cdots(l_1+\alpha_1)\cdots (l_j+1)\cdots(l_j+\alpha_j)}{n(n+1)\cdots(n+k)}f(l_1,\cdots,l_j)
%         \\&\le \dfrac{\max\limits (l_1+1)\cdots(l_1+\alpha_1)\cdots (l_j+1)\cdots(l_j+\alpha_j)f(l_1,\cdots,l_j)}{(n+1)\cdots(n+k)}\sum\limits_{l_1,\cdots,l_j}^n p_{i_1}^{\alpha_1+l_1}\cdots p_{i_j}^{\alpha_j+l_j} (1-p_{i_1}-\cdots-p_{i_j})^{n-l_1-\cdots-l_j}\dfrac{(n+k)!}{(l_1+\alpha_1)!\cdots(l_j+\alpha_j)!}\dfrac{(l_1+1)\cdots(l_1+\alpha_1)\cdots (l_j+1)\cdots(l_j+\alpha_j)}{n(n+1)\cdots(n+k)}
%         \\&\le \dfrac{\max\limits (l_1+1)\cdots(l_1+\alpha_1)\cdots (l_j+1)\cdots(l_j+\alpha_j)f(l_1,\cdots,l_j)}{(n+1)\cdots(n+k)}
%         % \\& = \dfrac{1}{n+1}\sum\limits_{l=0}^n (l+1)f(l)
%     \end{align*}
%     所以
%     \begin{align*}
%         &\sum\limits_{k\le n^{1/\alpha}}p_kE_{X}[\ferror(t_k)]
%         \\  \le& \dfrac{\max\limits_{l} (l+1)f(l)}{n+1} n^{1/\alpha}
%         \\ \le & Cn^{-1+1/\alpha}.
%     \end{align*}
%     又因为
%     \begin{align*}
%         &\sum\limits_{k> n^{1/\alpha}}p_kE_{X}[\ferror(t_k)]
%         \\  =& \sum\limits_{k> n^{1/\alpha}}\sum\limits_{l=0}^n p_k^{l+1} (1-p_k)^{n-l}\binom{l}{n}f(l)
%         \\ \le& \sum\limits_{k> n^{1/\alpha}}\sum\limits_{l=0}^n \left(\dfrac{1}{\zeta(\alpha)k^{\alpha}}\right)^{l+1} \binom{l}{n}f(l)
%         \\ \le& \sum\limits_{k> n^{1/\alpha}}\sum\limits_{l=0}^n \left(\dfrac{1}{\zeta(\alpha)k^{\alpha}}\right)\left(\dfrac{1}{n}\right)^l \binom{l}{n}f(l)
%         \\ \le & \sum\limits_{k> n^{1/\alpha}}\sum\limits_{l=0}^n \left(\dfrac{1}{\zeta(\alpha)k^{\alpha}}\right)\dfrac{f(l)}{l!}
%         \\ \le & C\sum\limits_{k> n^{1/\alpha}}\dfrac{1}{k^{\alpha}}
%         \\ \le& C_1\dfrac{1}{k^{\alpha-1}} = C_1n^{-1+1/\alpha}.
%     \end{align*}
%     所以\[\sum\limits_{k\le n^{1/\alpha}}p_kE_{X}[\ferror(t_k)]\le C_2n^{-1+1/\alpha}.\]
% \end{proof}
% \newpage

\subsection{Model Scaling Laws}
\label{appsec:model scaling}


In this section, we consider the optimal redundancy achievable by an omniscient model \( \Model^*_C \) under a model capacity constraint. (That is, the model \( \Model^*_C \) has access to the true data distribution \( \dataModel \).)  
Due to the finite capacity, the model \( \Model^*_C \) must apply lossy compression to each \( \phi_i \), resulting in a corresponding redundancy \( \modelerror^{(i)}(m_i) \), where  \( m_i \) denotes the constraint of the mutual information between
the model and $\phi_i$ (one may think it as the memory allocated to compress \( \phi_i \)).
The minimal achievable redundancy under a given memory constraint conditioned on a given cluster of knowledge is characterized by the  distortion-rate function:
\[
\ratedistortion_i(R) = \min_{\mathbb{I}(\phi_i; \Model^*_C) \le R} \mathbb{E}_{\phi_i\sim \basemeasure}\left[\mathbb{E}_{q \sim P_{\phi_i}} \left[ \DKL\left( P_{\phi_i}(\answer \mid \question=q) \,\Vert\, P_{\Model^*_C}(\answer \mid \question=q) \right) \right]\right],
\]
where \( \mathbb{I}(\phi_i; \Model^*_C) \) represents the information retained about \( \phi_i \) after compression.  
% Since \( \Model^*_C \) is omniscient and omnipotent, we assume it achieves the rate-distortion lower bound, i.e., the redundancy satisfies \( \modelerror^{(i)}(m_i) = \ratedistortion_i(m_i) \).

Furthermore, the omniscient model \( \Model^*_C \) allocates memory to each \( \phi_i \) in a manner that minimizes the expected redundancy. Formally, the memory allocation corresponds to the solution of the following optimization problem:
\begin{align}
    \text{minimize} \quad & \mathbb{E}_i[\ratedistortion_i(m_i)] = \sum_{i=1}^{\infty} p_i \ratedistortion_i(m_i),   \nonumber \\
    \text{subject to} \quad & \mathbb{I}(\Phi_0; \Model^*_C) \le C, \quad m_i = \mathbb{I}(\phi_i; \Model^*_C) \ge 0 \ \text{for all } i \in \mathbb{N}^+, \nonumber
\end{align}
where \( \Phi_0 = (\phi_1, \phi_2, \ldots) \) and $\mathbb{I}(\Phi_0; \Model^*_C) \le C$ represents the model capacity constraint.

% Let \( \Phi_{-i} = (\phi_{i+1}, \phi_{i+2}, \ldots) \).
% \jian{zhixuan. see the main text. do we really need this}
If we further assume that \( \phi_i \perp \Phi_{-i} \mid \Model^*_\capacity \), i.e., \( \phi_i \) is conditionally independent of the remaining cluster parameters given the model \( \Model^*_\capacity \),
where 
\( \Phi_{-i} = (\phi_{1},\ldots, \phi_{i-1}, \phi_{i+1}, \ldots) \).
In other words, conditional on the model \( \Model^*_\capacity \), knowing \(\phi_i\) does not provide any additional information about \(\phi_{j}\) for \(j \neq i\). 
% \jian{justification, look at what i wrote}
Under this assumption, the model capacity constraint can be decomposed as follows:
\begin{align*}
    C\ge\mathbb{I}(\Phi_0; \Model^*_C) 
    &= \mathbb{I}(\Phi_{-1}; \Model^*_C) + \mathbb{I}(\phi_1; \Model^*_C \mid \Phi_{-1}) \\
    &= \mathbb{I}(\Phi_{-1}; \Model^*_C) + \mathbb{I}(\phi_1; \Model^*_C) \\
    &= \mathbb{I}(\Phi_{-2}; \Model^*_C) + \mathbb{I}(\phi_2; \Model^*_C) + \mathbb{I}(\phi_1; \Model^*_C) \\
    &= \sum_{i=1}^{\infty} \mathbb{I}(\phi_i; \Model^*_C).
\end{align*}

Therefore, the optimization problem can be rewritten as:
\begin{align}
    \text{minimize} \quad & \sum_{i=1}^{\infty} p_i \ratedistortion_i(m_i), \label{app:model scaling} \\
    \text{subject to} \quad & \sum_{i=1}^{\infty} m_i \le C, \quad m_i \ge 0 \ \text{for all } i \in \mathbb{N}^+. \nonumber
\end{align}


% \jian{what is the definition of R? pi?}
Let \( \testmodelerror(C) \) denote the optimal value of the problem in \eqref{app:model scaling}, representing the minimal achievable redundancy under the model capacity constraint.


In general, the distortion-rate function is hard to express analytically. 
% Here, we consider a simplified scenario, in which the distortion $\ratedistortion_i(R)$ is on the same order as an exponential function with respect to $R$.
 % Here, we consider a simplified scenario, in which the distortion $\ratedistortion_i(R)$ satisfies the following assumptions.
 Here, we make the following assumptions.
 
\begin{assumption}
\label{appass:model scaling}
We assume that the distortion-rate function \( \ratedistortion_i(R) \) satisfies
\begin{itemize}
    % \item $D_i(R)$ is non-negative and monotonically non-increasing with respect to $R$,
    % \item $D_i(R)$ is convex,
    \item there exists positive constant $c_3,c_4$ such that for all $i\in \mathcal{N}^+,R\le c_3$, the distortion-rate function $\ratedistortion_i(R)\ge c_4$,
    \item there exists positive constant $c_{\max},b_{\max}$ such that for all $i\in \mathcal{N}^+$, the distortion rate function $\ratedistortion_i(R) \le c_{\max} b_{\max}^{-R}.$
\end{itemize}
% where \( c_{\min}, c_{\max}>0, b_{\min}, b_{\max}>1 \) are some constants.
\end{assumption}


% In this setting, we assume that the cluster parameters \( \phi_i \) in the data model \( \knwModel = \sum_{i=1}^{\infty} p_i \delta_{\phi_i} \) have been perfectly learned—that is, they are estimated without error in the infinite data limit. 






% We define a function similar to the rate-distortion function to characterize the achievable approximation quality of the target distribution under a given memory constraint.

% Denote
%     \begin{align*}
%         R_i(m) = \min_{P_{K}(\cdot|x))\in P_m} \{D_{\text{KL}}(P_{\theta_{x_i}}\| P_{K}(\cdot|x_i))\}
%     \end{align*}
% where $P_m$ denotes the set of distributions that can be stored within $m$ bits memory. 

% We further assume that different cluster of knowledge should be stored independently. Then, the optimal redundancy under the given model size constraint can be formulated as the optimal value of the following optimization problem.

% One possible approach is to treat the model size limitation as a constraint in an optimization problem. An optimal algorithm should find the best solution to this optimization problem, or we can use this optimal solution as the lower bound of the loss.

% Here, let's first consider an extreme case, where the data is infinite, and even every knowledge appears infinitely many times. Suppose it takes $k$ bits to remember a knowledge, resulting in a loss $g(k)$. In this case, the optimal loss under the model size limitation should be the solution to the following optimization problem.




% However, the optimization problem is challenging to solve when the exact form of the function $\ferror$ is unknown. To address this issue, we provide the following heuristic explanation from the perspective of floating-point precision in computing.



% First, we consider the case of storing a bounded real number within $m$ bits. For simplicity, we assume that this real number lies within $(0,1)$. Then storing it using $m$ bits means retaining $m$ decimal places in the binary representation. There are two types of rounding methods: one is deterministic but biased methods and the other one is stochastic and unbiased methods. The biased methods such as truncation introduce an absolute error of $\Theta(2^{-m})$. The unbiased methods such as stochastic rounding introduce an estimation variance of $\Theta(4^{-m})$. 

% Recall that our error term \(\ferror\) is the KL divergence between the target distribution and the predicted distribution. First, consider the case where the target distribution is determined by a single real-valued parameter. If we apply a deterministic but biased rounding method to this real number, the KL divergence will be roughly of the same order as the absolute error of this real number. If we apply a stochastic and unbiased rounding method, the KL divergence will be roughly of the same order as the variance of the rounding. In both cases, the error term is an exponential function of the storage space \(m\). For the case where the target distribution is determined by a \(d\)-dimensional parameter,  the error term remains an exponential function of the storage space \(m\) assuming we consider $d$ as a constant.



% The expected error should be $\Theta(2^{-m})$. In other words, the variance of the estimation is $\Theta(2^{-m})$. Then, using the derivation similar to \Cref{sec:Error of Syntax} which relates the variance of the estimation to the KL divergence of the induced distribution, we can similarly conclude that the KL divergence is also $\Theta(2^{-m})$. Thus we assume that $\ferror$ follows an exponential function to $m$.
% \begin{assumption}
% \label{assumption:model scaling}
%     For all $k,m$, we assume $c_{\min} b_{\min}^{-m}\le \ferror(k,m) \le c_{\max} b_{\max}^{-m}$ where $c_{\min},c_{\max},b_{\min},b_{\max}>0$ are some constant.
% \end{assumption}

% We assume the prior distribution of the cluster distribution $\phi_k$ is  $\text{Unif}(\{\phi \in \mathbb{R}^d,\|\phi\|_{\infty}\le 1\})$


% \begin{lemma} For some methods that lossy store $\phi_k$ as $\theta_k$, we have \[\modelerror(k,m_k)=\Theta(4^{-m_k/d})\]
% \end{lemma}
% \begin{proof}
%     We consider the following lossy compress method. When storing $\phi_k$ as $\theta_k$ using $m_k$ bits, we randomly rounding every coordinate of $\phi_k$ to $m_k/d$ bits in binary, which can be represented as:
% \[
% \theta_k^{(i)} =
% \begin{cases}
% \frac{\lfloor \phi_k^{(i)} \cdot 2^{m_k/d} \rfloor}{2^{m_k/d}} & \text{with probability } p = \phi_k^{(i)} \cdot 2^{m_k/d} - \lfloor \phi_k^{(i)} \cdot 2^{m_k/d} \rfloor, \\
% \frac{\lceil \phi_k^{(i)} \cdot 2^{m_k/d} \rceil}{2^{m_k}/d} & \text{with probability } 1 - p.
% \end{cases}
% \]
% where $\phi_k^{(i)}$ and $\theta_k^{(i)}$ represents the $i_{th}$ coordinate of $\phi_k$ and $\theta_k$. Note that we ignore the integer limit of bits here.

% Denote $p^{(i)}=\phi_k^{(i)} \cdot 2^{m_k/d} - \lfloor \phi_k^{(i)} \cdot 2^{m_k/d} \rfloor$.

% We know that this random rounding method is a unbiased estimation of $\phi_k$ and the variance \[\text{Var}(\theta_k^{(i)}-\phi_k^{(i)})=\dfrac{p^{(i)}(1-p^{(i)})}{4^{-m_k/d}}.\]

% Then by applying the conclusion of \Cref{lem: KL}, we know that 
% \[\mathbb{E}_{\theta_k}\left[\DKL (P_{\phi_k}\|P_{\theta_k})\right]=\Theta(4^{-m_k/d})\sum\limits_{i=1}^dp^{(i)}(1-p^{(i)}).\]

% Taking expectation with respect to $\phi_k$, we complete the proof.

% \end{proof}

% \zhixuan{Some discussion like rate-distortion of Gaussian}
% \jian{justify the above assumption}
The first assumption states that when the rate is small, the distortion cannot drop below a threshold, which is a natural property of rate-distortion theory for general sources.

The second assumptions can be formally justified within the framework of rate-distortion theory when the distortion function is the mean squared error (MSE) \citep{cover2006elements}. Although the distortion function in our setting is not MSE but rather the KL divergence between the probability distributions induced by the parameters, \Cref{lem:KL and MSE} shows that the two distortion measures exhibit similar behavior in a local neighborhood.

\begin{itemize}
% \item \textbf{Convexity of the Distortion-Rate Function}  
% The distortion-rate function \( D(R) \) is convex in the rate \( R \), thus its derivative \( D'(R) \) is monotonically increasing.

    \item \textbf{Distortion Rate Function of Gaussian Distribution:}
    Let $X \sim \mathcal{N}(0, \sigma^2)$ be a zero-mean Gaussian source, and let the distortion measure be the mean squared error (MSE).
The distortion-rate function $D(R)$ is defined as the minimum achievable distortion under a given rate $R$:

\[
D(R) = \min_{p(\hat{x}|x): \mathbb{I}(X; \hat{X}) \leq R} \mathbb{E}[(X - \hat{X})^2]
\]
For the Gaussian source with MSE distortion, the closed-form solution is:

\[
D(R) = \sigma^2 \cdot 2^{-2R}, \quad R \geq 0
\]
\item \textbf{Distortion-Rate Upper Bound via Maximum Entropy of Gaussian Distribution}  
Among all real-valued random variables with fixed variance, the Gaussian distribution achieves the maximum differential entropy:
\[
h(X) \leq \frac{1}{2} \log(2\pi e \sigma^2),
\]
with equality if and only if \( X \sim \mathcal{N}(0, \sigma^2) \). This implies that under a fixed distortion constraint and mean squared error (MSE) as the distortion measure, the Gaussian source requires the highest rate among all sources with the same variance.

Therefore, the distortion-rate function of any real-valued source with variance \( \sigma^2 \) must satisfy:
\[
D(R) \leq \sigma^2 \cdot 2^{-2R}, \quad R \geq 0,
\]
with equality achieved only when the source is Gaussian. This provides a universal lower bound on achievable distortion under MSE for a given rate \( R \).



%     \item \textbf{Distortion Rate Function from the Perspective of Numerical Precision}
%     We consider the case of storing a bounded real number using \( m \) bits. For simplicity, we assume that the real number lies within the interval \( (0,1) \). Representing it with \( m \) bits corresponds to retaining \( m \) binary digits in its fractional part. There are generally two types of rounding methods: deterministic but biased methods, and stochastic but unbiased methods. Biased methods, such as truncation, introduce an absolute MSE error on the order of \( \Theta(4^{-m}) \). In contrast, unbiased methods, such as stochastic rounding, introduce an estimation variance on the order of \( \Theta(4^{-m}) \). This exponential form can naturally extends to the case of storing a \( d \)-dimensional real-valued vector. 
%     \item \textbf{KL Divergence as the Distortion Measure}
%     While both of the above examples use MSE as the distortion measure, it is important to note that for a sufficiently regular parametric family, MSE and KL divergence behave similarly in a local neighborhood.
%     When two probability distributions \( p \) and \( q \) are close, the KL divergence behaves similarly to a weighted mean squared error. Specifically,
% \[
% D_{\mathrm{KL}}(p \| q) = \sum_i p_i \log \frac{p_i}{q_i},
% \]
% a second-order Taylor expansion around \( q_i = p_i \) yields:
% \[
% D_{\mathrm{KL}}(p \| q) \approx \sum_i \left( \frac{q_i - p_i}{p_i} + \frac{(q_i - p_i)^2}{2p_i^2} \right) p_i = \sum_i  \frac{(q_i - p_i)^2}{2p_i}
% \]
% This approximation shows that, in a local neighborhood, KL divergence and MSE have similar curvature and gradient behavior.
\end{itemize}

% Based on this lemma, we can solve the optimization problem \eqref{model scaling} and get the following result:

  
% \begin{theorem}
%  Under \Cref{ass:model scaling} and \Cref{ass:powerlaw}, the optimal value of the optimization problem under the given model size constraint satisfies:
% \[\testmodelerror(\capacity) = \Theta(\capacity^{-1/\alpha+1}).\] 
% % \jian{say something for each class?}
% % \zhixuan{we can not We cannot derive a rigorous theoretical conclusion for each class unless we assume that the rate-distortion function is exactly an exponential function}
% Moreover, if we further assume that $D_k(R)=a_kb_k^{-R}$ for some constant $b_{\min}\le b_k\le b_{\max},a_{\min}\le a_k\le a_{\max}$, we have \[p_k\ratedistortion_k(m_k^*)=\Theta(\min\{k^{-1/\alpha},C^{-1/\alpha}\}).\]
% where $m_k^*$ is the solution of the optimization problem.
% \end{theorem}
\begin{theorem}[Resteatement of \Cref{thm:model scaling}]
Under \Cref{ass:model scaling} and \Cref{ass:powerlaw}, the optimal value of the optimization problem under the given model size constraint satisfies:
\[\testmodelerror(\capacity) = \Theta(\capacity^{-1/\alpha+1}).\] 
% \jian{say something for each class?}
% \zhixuan{we can not We cannot derive a rigorous theoretical conclusion for each class unless we assume that the rate-distortion function is exactly an exponential function}
%\jian{Zhixuan:state the loss ,including the irreducible part}
Moreover, if we further assume that $D_k(R)=a_kb_k^{-R}$ for some constant $b_{\min}\le b_k\le b_{\max},a_{\min}\le a_k\le a_{\max}$,  
the contribution of the $k$-th cluster is
\[
p_k\ratedistortion_k(m_k^*)=\Theta(\min\{k^{-1/\alpha},C^{-1/\alpha}\}).
\]
where $m_k^*$ is the solution of the optimization problem.\eqref{app:model scaling}.
\end{theorem}
\begin{proof}[Proof of \Cref{thm:model scaling}]
% Denote the solution of the optimization problem as $m_i^*$. 
We first provide a lower bound for the optimal value.
By \Cref{ass:model scaling}, we have \( D_k(m_k) \ge c_4 \) whenever \( m_k \le c_3 \).
Since \( \sum\limits_{k=1}^{\infty} m_k \le C \), there can be at most \( \lfloor C / c_3 \rfloor \) indices \( k \) such that \( D_k(m_k) < c_4 \).
Due to the monotonicity of $p_k$, we know that 
\begin{align}
    \sum\limits_{k=1}^{\infty}p_kD_k(m_k)\ge \sum\limits_{k>\lfloor C / c_3 \rfloor}^{\infty}p_kc_4=\Omega(C^{-1/\alpha+1}).
    \label{eq:model scaling Omega}
\end{align}

% We consider the optimization problem as following:
% \begin{align}
%     \text{minimize} \quad & \sum_{i=1}^{\infty} p_i c_i b_i^{-m_i}, \label{app:model scaling general} \\
%     \text{subject to} \quad & \sum_{i=1}^{\infty} m_i \le C, \quad m_i \ge 0 \ \text{for all } i \in \mathbb{N}^+. \nonumber
% \end{align}


% Consider the following optimization problem:
% \begin{align}
%     \text{minimize} \quad & \sum_{i=1}^{\infty} p_i c_{\min} b_{\min}^{-m_i}, \label{app:model scaling lower bound} \\
%     \text{subject to} \quad & \sum_{i=1}^{\infty} m_i \le C, \quad m_i \ge 0 \ \text{for all } i \in \mathbb{N}^+. \nonumber
% \end{align}
% Denote the solution of \eqref{app:model scaling lower bound} as $m^{1}_i$ and the optimal value as $R_{\min}(C)$.

% Then we know that 
% \begin{align*}
%     R_{\min}(\capacity)=&\sum\limits_{i=1}^{\infty} p_i c_{\min}b_{\min}^{-m_i^{1}}
%     \\ \le & \sum\limits_{i=1}^{\infty} p_i c_{\min}b_{\min}^{-m_i^{0}}
%     \\ \le & \sum\limits_{i=1}^{\infty} p_i R_i(m_i^0)=\testmodelerror(\capacity)
% \end{align*}

Next, we solve the optimization problem \eqref{appsec:model scaling} under assumption $D_k(R)=a_kb_k^{-R}$ using KKT condition.
First consider the Lagrangian dual:
\begin{align*}
    \mathcal{L}(k, \lambda, \mu) = \sum_{k=1}^{\infty} p_k \ratedistortion_k(m_k) + \lambda \left(\sum_{k=1}^{\infty} m_k - C\right) - \sum_{k=1}^{\infty} \mu_k m_k. 
\end{align*}
We list the KKT conditions below:
\begin{align}
\label{stationary}
    & \text{Stationarity:} \quad p_k \ratedistortion_k'(m_k^*) + \lambda - \mu_k = 0, \quad \forall k \geq 1. \\
    & \text{Primal feasibility:} \quad \sum_{k=1}^{\infty} m_k^* \leq C, \quad m_k^* \geq 0, \quad \forall k \geq 1. \\
    & \text{Dual feasibility:} \quad \lambda \geq 0, \quad \mu_k \geq 0, \quad \forall k \geq 1. \\
    & \text{Complementary slackness:} \quad \lambda \left( \sum_{k=1}^{\infty} m_k^* - C\right) = 0, \quad \mu_k m_k^* = 0, \quad \forall k \geq 1.
\end{align}
Due to monotonicity, we know that $\sum\limits_{k=1}^{\infty} m_k^* = C$. 
By the Stationary condition \eqref{stationary}, we know that
\begin{align*}
   \ratedistortion'_k(m_k^*)= \dfrac{\mu_k-\lambda}{p_k}.
\end{align*}
Due to the convexity of $\ratedistortion_k$\citep{cover2006elements}, we know that $\ratedistortion_k'$ is monotonically increasing. Therefore, the inverse of $\ratedistortion_k'$ exists, and we denote it as $g_k$.

% \zhixuan{citation}

For $m_k^*>0$, we have $\mu_k=0$, thus we know that 
\[m_k^*=\max\left\{g_k\left(\dfrac{-\lambda}{p_k}\right),0\right\}.\]

Next, we estimate the value of $\lambda$. 
Since we have assumed that $D_k(R)=a_kb_k^{-R}$ for some $b_{\min}\le b_k\le b_{\max},a_{\min}\le a_k\le a_{\max}$, so there exists $r_1',r_2',d_1,d_2$ such that 
\[
-r_1'\ln(d_1x)\le g_k(-x)\le -r_2'\ln(d_2x).
\]
Then, we know that \[
\max\left\{0,r_1\ln \dfrac{d_1k^{-1/\alpha}}{\lambda}\right\} \le m_k^*\le  \max\left\{0,r_2\ln \dfrac{d_2k^{-1/\alpha}}{\lambda}\right\}.
\] 
where $r_1=r_1'\ln \zeta(1/\alpha), r_2=r_2'\ln\zeta(1/\alpha)$. 

Assume $l_1$ is the maximum integer such that $d_1l_1^{-1/\alpha}> \lambda$. 
We know that $d_1(l_1+1)^{-1/\alpha}\le \lambda< d_1l_1^{-1/\alpha}.$ Therefore, we have
\begin{align*}
    C  & \ge \sum\limits_{k=1}^{\infty} \max\left\{0,r_1\ln \dfrac{d_1k^{-1/\alpha}}{\lambda}\right\}
    =r_1 \sum\limits_{k=1}^{l_1} \ln\dfrac{k^{-1/\alpha}}{\lambda} \\
    & \ge r_1\sum\limits_{k=1}^{l_1} \ln \dfrac{k^{-1/\alpha}}{l_1^{-1/\alpha}}
     =\dfrac{r_1}{\alpha}
     \left(\ln l_1^{l_1} - \ln l_1!\right).
\end{align*}
By Stirling's formula, we know that $\ln l_1^{l_1} - \ln l_1!=\Theta(l_1)$. So we know that 
\[
C\ge \Theta(l_1)=\Theta(\lambda^{-1/\alpha}).
\]
Similarly, assume $l_2$ is the maximum integer such that $d_2l_2^{-1/\alpha}> \lambda$. 
We know that $d_2(l_2+1)^{-1/\alpha}\le \lambda< d_2l_2^{-1/\alpha}.$ Therefore, we have
\begin{align*}
    C  & \le \sum\limits_{k=1}^{\infty} \max\left\{0,r_2\ln \dfrac{d_2k^{-1/\alpha}}{\lambda}\right\}
    =r_2 \sum\limits_{k=1}^{l_2} \ln\dfrac{k^{-1/\alpha}}{\lambda} \\
    & \le r_2\sum\limits_{k=1}^{l_2} \ln \dfrac{k^{-1/\alpha}}{(l_2+1)^{-1/\alpha}}
     =\dfrac{r_2}{\alpha}
     \left(\ln (l_2+1)^{l_2+1} - \ln l_2!\right).
\end{align*}
By Stirling's formula, we know that $\ln l_2^{l_2} - \ln l_2!=\Theta(l_2)$. So we know that 
\[
C\ge \Theta(l_2)=\Theta(\lambda^{-1/\alpha}).
\]
Since we know that $\Theta(l_1)=\Theta(\lambda^{-1/\alpha})=\Theta(l_2)$, we have that
\[
C=\Theta(l_1)=\Theta(l_2)=\Theta(\lambda^{-1/\alpha}).
\]
For $k\le l_1$, we have that
$$
m_k^*\ge \max\left\{0,r_1\ln \dfrac{d_1k^{-1/\alpha}}{\lambda}\right\}\ge r_1\ln \dfrac{d_1k^{-1/\alpha}}{\lambda}>0.
$$
For $k> l_2$, we have that
$$
m_k^*\le \max\left\{0,r_2\ln \dfrac{d_2k^{-1/\alpha}}{\lambda}\right\}=0.
$$
Thus $m_k^*=0$ for $k > l_2$.
Hence, we have that  \[p_k\ratedistortion_k(m_k^*)=\Theta(\min\{k^{-1/\alpha},C^{-1/\alpha}\}).\]

Finally, we prove that $\testmodelerror(C)\ge \Theta(M^{-1/\alpha+1})$. 
Since we have assumed that $D_k(R)=a_kb_k^{-R}$ for some $b_{\min}\le b_k\le b_{\max},a_{\min}\le a_k\le a_{\max}$, we know that there exists constant $d_3,d_4$ such that 
% \jian{typo?? not finished??}
\[d_3\ratedistortion_k'(m^*_k)\le\ratedistortion_k(m^*_k)\le d_4\ratedistortion_k'(m^*_k)\]
Thus we know that 
\begin{align}
\label{eq:1}
    \sum\limits_{k=1}^{\infty}p_k \ratedistortion_k(m_k^*) \notag
 & \ge d_3\sum\limits_{k\le l_2}\lambda + c_3\sum\limits_{k>l_2} p_k\notag \\
 & = \Theta(\lambda l_2) + \Theta(l_2^{-1/\alpha+1}) = \Theta(M^{-1/\alpha+1}).
\end{align}
Similarly, 
% we can prove $\testmodelerror(C)\le \Theta(C^{-1/\alpha+1})$.
\begin{align}
\label{eq:2}
    \sum\limits_{k=1}^{\infty}p_k \ratedistortion_k(m_k^*)
 &\le d_4\sum\limits_{k\le l_1}\lambda + c_4\sum\limits_{k>l_1} p_k
 \\&= \Theta(\lambda l_1) + \Theta(l_1^{-1/\alpha+1}) = \Theta(C^{-1/\alpha+1})
\end{align}
Combining \Cref{eq:1} and \Cref{eq:2} together, we have that 
\[ 
\sum\limits_{k=1}^{\infty}p_k \ratedistortion_k(m_k^*)=\Theta(C^{-1/\alpha+1}).
\]
Note that the above solution is carried out under a strong assumption $D_k(R)=a_kb_k^{-R}.$
Thus under \Cref{ass:model scaling}, we have $\testmodelerror(C)\le \Theta(C^{-1/\alpha+1})$. Combining this and \eqref{eq:model scaling Omega}, we have that 
\[
\testmodelerror(C)= \Theta(C^{-1/\alpha+1}).
\]
\end{proof}

% \subsection{General Scaling Law}
% We combine the content from the \Cref{sec:data scaling} and \Cref{sec:model scaling}  to derive a power law for error with respect to data size $N$ and model size $M$, which has the same form with the scaling law in LLMs.
% \begin{assumption} 
% \label{assumption: f}
% We make the following assumptions about $\ferror(k,m,t)$:
%     \begin{itemize}
%         \item $\ferror(k,t,m)=f_{\text{data}}(k,t)+f_{\text{model}}(k,m)$, 
%         \item $f_{\text{data}}(k,0)\ge C_1, (t+1)f_{\text{data}}(k,t)\le C_2$ for all $k,t \in \mathcal{N}^+$, 
%         \item $f_{\text{model}}(k,m)=\alpha_k 2^{-m}$ where $\alpha_k \in [\alpha_{\min},\alpha_{\max}]$ for all $k\in \mathcal{N}^+$ and some constant $\alpha_{\min},\alpha_{\max}>0$.
%     \end{itemize}
% \end{assumption}

% \begin{theorem}
% \label{thm: general scaling law}
%       Denote $R(N,M)$ as the optimal value of the following optimization problem:
%       \begin{align*}
%     \text{minimize} \quad & \mathbb{E}_{k,t_k}[\ferror(k,m_k,t_k)] = \sum\limits_{k=1}^{\infty} p_k \mathbb{E}_{t_k}[\ferror(k,m_k,t_k)], \nonumber \\
%     \text{subject to} \quad & \sum\limits_{k=1}^{\infty} m_k \le M,  \quad  m_k \ge 0 \ \ \text{for all} \ k\in \mathcal{N}^+.
% \end{align*}
% Under \Cref{assumption: f}, we have 
% \begin{align*}
%     R(N,M) = \Theta(N^{-1+1/\alpha}) + \Theta(M^{-\alpha+1}).
% \end{align*}
      
% \end{theorem}










% From the information theory perspective, suppose the original entropy of the $i$-th knowledge is $h_i$, and after using $k_i$ bits to encode it, the remaining entropy is $h_i - k_i$. In other words, $g(k_i) = h_i - k_i$. The solution to this optimization problem is evident and corresponds to a loss of $\Theta(n^{-1+1/\alpha})$.
% \subsection{General Scaling Law}
% In the previous two sections, we analyzed the data scaling law in the case of an infinitely large model and the model scaling law in the case of infinitely many data. In this section, we attempt to consider both data and model limitations simultaneously.

% \subsection{Discussion}
% The main reason for the Scaling Law, as discussed here, is that when learning World Knowledge, the frequency of occurrence of World Knowledge itself follows the Scaling Law.

% Based on our analysis, solving the issue of insufficient data with synthetic data or data generated by large models is not feasible. This is because synthetic data is completely useless in learning World Knowledge.



% scaling law

% use stochastic grammar to encode knowledge

% knowledge follows generalize heap's law.

% submodularity of entropy function 

% to better compress (minimize the cross entropy)
% first learn the grammar, then the knowledge
% (according to the frequency).

% then we get scaling law

% three aspects

% (1) as data (knowledge) increases

% (2) as model size increases

% (3) as learning progress (e.g., learning rates), harder and harder
% knowledge are learned.

% Universal coding for stationary and ergodic process/
% individual sequence/nonstationary process
% --mixture of stationary process--nonstationary entropy rate

% (this relates to ICL)
% fact like task vector, task superposition, skill composition.
\section{Fine-Tuning}
\label{sec:instruction}

Fine‐tuning an LLM is typically used for two broad purposes: (1) instruction‐following or (2) knowledge injection. Here, we focus on a standard instruction fine‐tuning scenario. Consider, for instance, in our experiment, a model that has been pretrained on a bioS (multi+permute) dataset \citep{allen2024physics}
and has already learned factual knowledge about specific individuals. During fine‐tuning, our goal is to fine-tune the model to be able to answer the questions about these individuals, and thus the fine‐tuning dataset consists of question–answer pairs with facts drawn from the same distribution as the pretraining data.








% Within our theoretical framework, this setting can be viewed as generating data from the same knowledge model $\knwModel$ but with a {\em different syntax encoder} $\instructionModel$ (e.g. the question–answer format). 
We view the generation of pretraining data and instruction fine-tuning data as a two-stage data generation process. 
The first stage, corresponding to the generation of pretraining data, follows the data model introduced in \Cref{app:datageneration}. 
In the second stage, i.e., the instruction fine-tuning stage, the knowledge model and the generation process of knowledge elements in the  data remain unchanged; however, the syntactic elements are produced by a different syntax model (e.g., the question–answer format), which we denote by \( \instructionModel \).

Let \( X_{1:N} \) denote the pretraining corpus, which is generated by
$(\synModel,\knwModel)$ and \( X_{N+1:N+n} \) the instruction fine-tuning corpus, generated by $(\instructionModel,\knwModel)$. 
The full corpus is denoted by \( X_{1:N+n} \).
% We now analyze the redundancy of the instruction fine-tuning corpus within the theoretical framework developed in \Cref{sec:Bayesian Framework}. 
% This redundancy can be expressed as the difference between the redundancy of the full corpus \( X_{1:N+n} \) and that of the pretraining corpus \( X_{1:N} \).
% Specifically, the redundancy of the full corpus \( X_{1:N+n} \) is given by the mutual information 
% \( \mathbb{I}(X_{1:N+n}; \synModel, \instructionModel, \knwModel) \), 
% while the redundancy of the pretraining corpus is \( \mathbb{I}(X_{1:N}; \synModel, \knwModel) \). 
% Therefore, the redundancy of the instruction fine-tuning corpus \( X_{N+1:N+n} \) is:
% \[
% \mathbb{I}(X_{1:N+n}; \synModel, \instructionModel, \knwModel) - \mathbb{I}(X_{1:N}; \synModel, \knwModel).
% \]
Following the same reasoning as in Section~\ref{sec:baye data law}, 
the Bayesian redundancy of the full corpus \( X_{1:N+n} \) 
is given by the mutual information  
$$
\mathbb{I}(X_{1:N+n}; \synModel, \instructionModel, \knwModel). 
$$ 
In the following, we present a (somewhat heuristic) decomposition of the redundancy 
over the full corpus. 
The $\approx$ in the derivation below requires certain independence 
assumption, which may not hold exactly in real world,
and bounding the approximation error is left as a future work.
% \begin{align*}
%         &\mathbb{I}(X_{1:N+n}; \synModel, \instructionModel, \knwModel)
%     \\ \overset{(1)}{=}& \mathbb{I}(X_{1:N}; \synModel, \instructionModel, \knwModel)+\mathbb{I}(X_{N+1:N+n}; \synModel, \instructionModel, \knwModel|X_{1:N})
%     \\ \overset{(2)}{=}& \mathbb{I}(X_{1:N}; \synModel, \knwModel)+\mathbb{I}(X_{N+1:N+n};  \instructionModel, \knwModel|X_{1:N})+\mathbb{I}(X_{N+1:N+n};\synModel |X_{1:N},\instructionModel, \knwModel)
%     \\ \overset{(3)}{=}& \mathbb{I}(X_{1:N}; \synModel, \knwModel)+\mathbb{I}(X_{N+1:N+n};  \instructionModel, \knwModel|X_{1:N})
%     \\ \overset{(4)}{=}& \mathbb{I}(X_{1:N}; \synModel, \knwModel)+\mathbb{I}(X_{N+1:N+n};   \knwModel|X_{1:N}) + \mathbb{I}(X_{N+1:N+n}; \instructionModel|X_{1:N}, \knwModel)
%     \\ \overset{(5)}{=}& \mathbb{I}(X_{1:N}; \synModel, \knwModel)+\mathbb{I}(K_{N+1:N+n};   \knwModel|X_{1:N}) + \mathbb{I}(S_{N+1:N+n}; \knwModel|X_{1:N},K_{N+1:N+n})
%     \\+&\mathbb{I}(S_{N+1:N+n}; \instructionModel|X_{1:N}, \knwModel)+\mathbb{I}(K_{N+1:N+n}; \instructionModel|X_{1:N}, \knwModel,S_{N+1:N+n})
%     \\ \overset{(6)}{=}& \mathbb{I}(X_{1:N}; \synModel, \knwModel)+\mathbb{I}(K_{N+1:N+n};   \knwModel|X_{1:N}) + \mathbb{I}(S_{N+1:N+n}; \instructionModel|X_{1:N}, \knwModel)
%     \\ \overset{(7)}{=}& \mathbb{I}(X_{1:N}; \synModel, \knwModel)+\mathbb{I}(K_{N+1:N+n};   \knwModel|X_{1:N}) + \mathbb{I}(S_{N+1:N+n}; \instructionModel|\knwModel)
% \end{align*}
% Equations (1), (2), (4), and (5) apply the chain rule of mutual information.  
% Equation (3) applies the conditional independence \( X_{N+1:N+n} \perp \synModel \mid X_{1:N}, \instructionModel, \knwModel \).  
% Equation (6) applies the conditional independencies \( S_{N+1:N+n} \perp \knwModel \mid X_{1:N}, K_{N+1:N+n} \) and \( K_{N+1:N+n} \perp \instructionModel \mid X_{1:N}, \knwModel, S_{N+1:N+n} \).  
% Equation (7) applies the conditional independencies \( S_{N+1:N+n} \perp X_{1:N} \mid \knwModel \) and \( \instructionModel \perp X_{1:N} \mid \knwModel \).

As in our data generation process, there is a one-to-one mapping between
a sentence $X$ and its syntax-knowledge element pair $(\Synelem,\Knwelem)$.
Hence, using the chain rule of mutual information, we can write
$$
    \mathbb{I}(X_{1:N+n}; \synModel, \instructionModel, \knwModel) \notag
    = \mathbb{I}(\Knwelem_{1:N+n}; \synModel, \instructionModel, \knwModel)
    +\mathbb{I}(\Synelem_{1:N+n}; \synModel, \instructionModel, \knwModel\mid\Knwelem_{1:N+n}) 
$$
Also note that the first and second terms in the RHS can be decomposed as
$$
\mathbb{I}(\Knwelem_{1:N+n}; \synModel, \instructionModel, \knwModel)
=\mathbb{I}(\Knwelem_{1:N+n}; \knwModel)+\mathbb{I}(\Knwelem_{1:N+n}; \synModel,\instructionModel\mid\knwModel),
$$
$$
\mathbb{I}(\Synelem_{1:N+n}; \synModel, \instructionModel, \knwModel\mid\Knwelem_{1:N+n}) 
=
\mathbb{I}(\Synelem_{1:N+n}; \synModel, \instructionModel\mid\Knwelem_{1:N+n})+
\mathbb{I}(\Synelem_{1:N+n}; \knwModel\mid\synModel, \instructionModel,\Knwelem_{1:N+n}). 
$$
Note that
$\mathbb{I}(\Knwelem_{1:N+n}; \synModel,\instructionModel\mid\knwModel)=0$
since the knowledge elements are generated independently of the syntax model,
and 
$\mathbb{I}(\Synelem_{1:N+n}; \knwModel\mid\synModel, \instructionModel,\Knwelem_{1:N+n})=0$
since $\knwModel$ -- $\Knwelem$ -- $\Synelem$ form a Markov chain.
Hence, we can see that
\begin{align}
    \mathbb{I}(X_{1:N+n}; \synModel, \instructionModel, \knwModel)
     =& \mathbb{I}(\Knwelem_{1:N+n}; \knwModel)+\mathbb{I}(\Synelem_{1:N+n}; \synModel,\instructionModel\mid\Knwelem_{1:N+n}) \notag
    \\ 
    \smash{\overset{\scriptstyle (1)}{\approx}}& \mathbb{I}(\Knwelem_{1:N+n}; \knwModel)+\mathbb{I}(\Synelem_{1:N+n}; \synModel,\instructionModel)  \notag
    \\ 
    =& \mathbb{I}(\Knwelem_{1:N+n}; \knwModel)+\mathbb{I}(\Synelem_{1:N}; \synModel,\instructionModel)+\mathbb{I}(\Synelem_{N+1:N+n}; \synModel,\instructionModel\mid\Synelem_{1:N}) \notag
    \\ 
    =& \mathbb{I}(\Knwelem_{1:N+n}; \knwModel)+\mathbb{I}(\Synelem_{1:N}; \synModel)+\mathbb{I}(\Synelem_{1:N}; \instructionModel\mid\synModel) \notag \\
    +&\mathbb{I}(\Synelem_{N+1:N+n}; \instructionModel\mid\Synelem_{1:N})+\mathbb{I}(\Synelem_{N+1:N+n}; \synModel\mid\instructionModel,\Synelem_{1:N})  \notag
    \\ 
    \smash{\overset{\scriptscriptstyle (2)}{\approx}}& \mathbb{I}(\Knwelem_{1:N+n}; \knwModel)+\mathbb{I}(\Synelem_{1:N}; \synModel)+\mathbb{I}(\Synelem_{N+1:N+n}; \instructionModel) \notag
\end{align}
In the above, approximation $\smash{\overset{\scriptscriptstyle (1)}{\approx}}$ relies on the assumption \( \Synelem_{1:N+n} \perp \Knwelem_{1:N+n} \),
and $\smash{\overset{\scriptscriptstyle (2)}{\approx}}$ requires \( \Synelem_{N+1:N+n} \perp \Synelem_{1:N} \) and \( \Synelem_{N+1:N+n} \perp \synModel \mid \instructionModel, \Synelem_{1:N} \), 
both of which can be induced by the same assumption \( \Synelem_{1:N+n} \perp \Knwelem_{1:N+n} \).
Under the same assumption, the redundancy of the pretraining corpus can also be decomposed as follows:
\begin{align*}
    \mathbb{I}(X_{1:N}; \synModel, \knwModel) 
     =& \mathbb{I}(\Knwelem_{1:N}; \synModel, \knwModel)+\mathbb{I}(\Synelem_{1:N}; \synModel, \knwModel\mid\Knwelem_{1:N}) \notag
    \\
%    =& \mathbb{I}(\Knwelem_{1:N}; \knwModel)+\mathbb{I}(\Knwelem_{1:N}; \synModel\mid\knwModel)+\mathbb{I}(\Synelem_{1:N}; \synModel\mid\Knwelem_{1:N})+\mathbb{I}(\Synelem_{1:N}; \knwModel\mid\synModel, \Knwelem_{1:N}) \\
    =& \mathbb{I}(\Knwelem_{1:N}; \knwModel)+\mathbb{I}(\Synelem_{1:N}; \synModel\mid\Knwelem_{1:N})
    \\ 
    \approx & \mathbb{I}(\Knwelem_{1:N}; \knwModel)+\mathbb{I}(\Synelem_{1:N}; \synModel).  
\end{align*}
The redundancy of the instruction fine-tuning corpus can be written as the difference between the redundancy of the full corpus and the redundancy of the pretraining corpus:
\begin{align}
    &\mathbb{I}(X_{1:N+n}; \synModel, \instructionModel, \knwModel)-\mathbb{I}(X_{1:N}; \synModel, \knwModel) \notag
     \approx \underbrace{\mathbb{I}(\Knwelem_{1:N+n}; \knwModel)-\mathbb{I}(\Knwelem_{1:N}; \knwModel)}_{\redundency_{\text{knw}}}+\underbrace{\mathbb{I}(\Synelem_{N+1:N+n}; \instructionModel)}_{\redundency_{\text{ins}}}
\end{align}
The two terms in the above equation 
correspond to the redundancy of knowledge elements $\redundency_{\text{knw}}$, and the redundancy of syntactic elements $\redundency_{\text{ins}}$ in the instruction fine-tuning corpus, respectively.

% \begin{align*}
%     &\mathbb{I}(X_{1:N}; \synModel, \knwModel)
%     \\=& \mathbb{I}(\Knwelem_{1:N}; \synModel, \knwModel)+\mathbb{I}(\Synelem_{1:N}; \synModel,  \knwModel|\Knwelem_{1:N})
%     \\=& \mathbb{I}(\Knwelem_{1:N}; \knwModel)+\mathbb{I}(S_{1:N}; \synModel|\Knwelem_{1:N})
% \end{align*}



% \begin{align*}
%     &\mathbb{I}(X_{1:N+n}; \synModel, \instructionModel, \knwModel) - \mathbb{I}(X_{1:N}; \synModel, \knwModel)
%     \\ =& \mathbb{I}(\Knwelem_{1:N+n}; \synModel, \instructionModel, \knwModel)+\mathbb{I}(S_{1:N+n}; \synModel, \instructionModel, \knwModel |\Knwelem_{1:N+n}) \\-& \left(\mathbb{I}(\Knwelem_{1:N}; \synModel, \knwModel)+\mathbb{I}(S_{1:N}; \synModel, \knwModel |\Knwelem_{1:N})\right)
% \end{align*}



% Under additional independence assumptions, we can decompose the redundancy of the instruction fine-tuning corpus \( X_{N+1:N+n} \).
Following the same arguments in the proof of Theorem~\ref{thm:data_scaling_Mutual_Information}, we can see that
\( \mathbb{I}(\Knwelem; \knwModel) = \Tilde{O}(N^{\alpha}) \),
and hence for simplicity, we further assume that \( \mathbb{I}(\Knwelem; \knwModel) = cN^{\alpha} \) for some constant \( c > 0 \).
Then,
the average redundancy of the instruction fine-tuning corpus satisfies:
\begin{align*}
  &  \dfrac{1}{n}\redundency_{\text{knw}}=\mathbb{I}(\Knwelem_{1:N+n}; \knwModel)-\mathbb{I}(\Knwelem_{1:N}; \knwModel) = \dfrac{(N+n)^{\alpha}-N^{\alpha}}{n}\overset{(1)}{=}O(N^{\alpha-1}),
 \\   
 & \dfrac{1}{n}\redundency_{\text{ins}}=\mathbb{I}(S_{N+1:N+n}; \instructionModel)= \Tilde{O}(n^{-1}).
\end{align*}
where $\overset{(1)}{=}$ relies on the natural assumption that the size of pretrain corpus $N$ is much larger than the size of instruction fine-tuning corpus $n$.
Similar to Theorem~\ref{thm:data_scaling_Mutual_Information}, the primary effect of fine‐tuning in this scenario is that the model learns the new syntax to reduce the fine‐tuning loss (i.e., redundancy) with fast rate $\Tilde{O}(n^{-1})$, hence requiring relative fewer samples, while allowing the model to retain the factual knowledge it has already acquired
($N$ is large and the redundancy for the knowledge elements
$O(N^{\alpha-1})$ is quite small).

On the other hand, consider the setting where our objective is to {\em inject 
new knowledge} via fine‐tuning.
Suppose the fine‐tuning data uses a drastically different syntax or format from that used in pretraining, as well as completely new knowledge.
Now, the Bayesian redundancy is equal to
\begin{align}
\mathbb{I}(X_{1:N+n}; \synModel, \instructionModel, \knwModel, \knwModelNew)
\label{eq:miins}
\end{align}
where $\knwModelNew$ is the new knowledge model.
If the new syntax component is very different from the original one,
and the knowledge component is also quite new (i.e.,
nearly independent from the pretraining), the mutual information
\eqref{eq:miins} can be approximately decomposed into 
$$
\mathbb{I}(X_{1:N+n}; \synModel, \instructionModel, \knwModel, \knwModelNew)
\approx
\mathbb{I}(X_{1:N}; \synModel,  \knwModel)
+\mathbb{I}(X_{N+1:N+n}; \instructionModel, \knwModelNew). 
$$
The first term corresponds to the redundancy of the pretraining phase,
and the second term corresponds to finetuning and the redundancy (per sentence) can be bounded
as $\widetilde{O}(c_1/n+c_2/n^{1-\alpha})$ according to the same argument
in Theorem~\ref{thm:data_scaling_Mutual_Information}. As $n$ is much smaller
than $N$, we can see that the fine-tuning would experience a significant perplexity shift (even the entropy, which is the irreducible part of the loss, may also shift).
Moreover, the additional syntax learning step not only slows down the learning, but also risks occupying the model's limited capacity, leading to forgetting of the pretrained knowledge, especially when the model's capacity is constrained.
See the experimental results in Appendix~\ref{app:additionalexp}.
Therefore, two practical recommendations are in order (consistent with good practices in prior empirical works \citep{zhang2025unveiling, grattafiori2024llama, yang2024qwen2}).

\begin{itemize}
\item Knowledge Injection: When injecting new knowledge during fine-tuning, avoid drastically different formats or syntaxes that deviate significantly from the original pretraining data, especially in the capacity constrained setting. By preserving familiar structures, the model can focus on absorbing new facts rather than first learning an unfamiliar syntax. Moreover, 
it is beneficial to adopt a gradual approach that mixes the pretrained data with a portion of new knowledge during finetuning, which
helps mitigate large perplexity shifts and makes the optimization more stable \citep{guptacontinual,ibrahimsimple}.

\item Instruction Fine‐Tuning. For instruction fine‐tuning (e.g., adapting a model to question–answer formats), use a knowledge set that closely aligns with the data distribution seen during pretraining. 
This ensures that most of the fine‐tuning effort is spent on learning new syntax or style. Moreover, if we use completely a new set of knowledge in the SFT stage, the model only focuses on compressing the SFT dataset, which slows down the learning and may lead to forgetting of pretrained knowledge as well.
\end{itemize}

\section{Experimental Setting and Additional Experimental Results}
\label{app:experiment}


\subsection{Experiment Setting}


Following the experimental setting from \citep{AllenZhu-icml2024-tutorial, allen2024physics}, we generate profiles for 400,000 individuals. Each profile contains six attributes: date of birth, birth city, university, major, employer, and employer city. These attributes are used to populate a diverse set of templates (each type of information/instruction has 50 different templates), forming both pretraining and instruction fine-tuning datasets, as shown in Table~\ref{tab:datasets}. In the instruction tuning, we use data instances with odd indices (e.g., 1, 3, 5, 7, ...) as the training set and those with even indices (e.g., 2, 4, 6, 8, ...) as the test set.

\begin{table}[h]
    \centering
    \caption{Examples of Pretraining and Instruction Fine-Tuning Data}
    \begin{tabular}{c p{4cm}}
        \toprule
        \textbf{Dataset Type} & \multicolumn{1}{c}{\textbf{Example}} \\
        \hline
        Pretraining & ``Gracie Tessa Howell was born in Camden, NJ. He studied Biomedical Engineering and worked at UnitedHealth Group. He entered the world on April 15, 2081, and is employed in Minnetonka. He is an alumnus/alumna of Buena Vista College.'' \\
        \hline
        Instruction Fine-Tuning & ``Q: What area of study did Gracie Tessa Howell focus on? A: Biomedical Engineering'' \\
        \bottomrule
    \end{tabular}
    
    \label{tab:datasets}
\end{table}

As discussed in our theoretical section, the occurrence frequency of individuals in the pretraining/instruction fine-tuning dataset follows a
power-law distribution, formulated as Equation~\ref{eq:powerlaw}:
% \jian{indicate how it corresponds to the parameters in PY model.}

\begin{equation}
    P(i) = \frac{(i + \text{bias})^{-a}}{\sum_{j=1}^{N} (j + \text{bias})^{-a}}
\label{eq:powerlaw}
\end{equation}
where $i$ denotes the index of the individual, $N$ is the total number of individuals, and $a$ is the exponent parameter. Note that the frequency of data occurrence generated by a power-law distribution exhibits asymptotic behavior similar to that produced by the 
 $\PYCRP(1/a,\beta)$ (See \Cref{lem:PYCRP fre}).


We set $\text{bias} = 1000$ and vary the parameter $a$ over $\{0, 1.05, 1.20, 1.35, 1.50\}$. If not specified, we choose $a=1.20$. The final pretraining dataset contains 1.45B tokens (tokenized using the GPT-2 tokenizer), and the instruction fine-tuning dataset consists of 258M tokens. Pretraining is conducted for 4 epochs, totaling 5.8B tokens processed—almost twice the amount suggested by the Chinchilla scaling law.
\paragraph{Model Configuration}

We conduct experiments with RoPE-encoded GPT-like models \cite{modded_nanogpt_2024} of various sizes. The model configurations are detailed in Table~\ref{tab:model_sizes}.

\paragraph{Training Procedure}

For training, we use a sequence length of 512 and batch size of 128. We apply a warmup ratio of 0.05 and a warmdown ratio of 0.9. Pretraining runs for 4 epoch with a learning rate of 0.0003, while fine-tuning runs for 1 epochs with a reduced learning rate of 0.00003. We employ weight decay of 0.1 and bf16 precision. To enhance parallelism, multiple sequences are packed into 512-token sequences, but cross-sequence attention is masked out.




\subsection{Experiment Results}
\label{app:additionalexp}

\paragraph{Data Heterogeneity} \label{para:data_hetero}
\begin{figure}[htbp]
\centering
\begin{minipage}{0.48\linewidth}
  \centering
  \includegraphics[width=\linewidth]{property_by_model_size.png}
\end{minipage}
\hfill
\begin{minipage}{0.48\linewidth}
  \centering
  \includegraphics[width=\linewidth,scale=0.95]{property_by_model_size.jpg}
\end{minipage}
\caption{
Accuracy of different properties across varying model sizes. The left panel shows the results under the uniform setting, while the right panel corresponds to the power-law setting. In both cases, model accuracy across properties exhibits heterogeneity—i.e., models prioritize certain properties over others. However, the trend differs: under the power-law setting, accuracy improves gradually across all properties as model size increases, whereas the uniform setting displays a sharp phase transition—accuracy remains low until a critical model size is reached, after which it rapidly improves.
}
\label{fig:property_capacity_frequency}
\end{figure}


We evaluate the accuracy of various properties under two different data distributions—uniform and power-law—to examine the heterogeneous nature of model learning. For each setting, we train models of varying sizes and measure their accuracy across multiple properties, revealing how capacity constraints influence what types of knowledge are learned first.

As shown in Figure~\ref{fig:property_capacity_frequency}, both settings demonstrate that models with limited capacity exhibit selective learning behavior—some properties are prioritized over others. Properties with lower entropy, such as Major, are easier to learn, as they contain less variability and thus can be compressed more efficiently by small models. In contrast, high-entropy properties require more capacity to capture and generalize. A notable example is the property \textit{company city}. In both settings, this attribute starts with a non-trivial accuracy of approximately 7\%, even when other properties remain close to zero. This is due to New York appearing with 7\% frequency in the dataset. Small models, unable to generalize meaningfully, default to outputting the most frequent token—an effect more visible in the uniform setting where each person appears the same number of times in the pretraining corpus, i.e. no person is inherently easier. 

However, the dynamics of this heterogeneity differ across the two distributions. Under the \textbf{power-law setting} (right), model accuracy for each property increases gradually as capacity grows. This suggests a smooth transition in learning, where dominant and compressible properties are acquired first, followed by rarer or more complex ones as capacity permits. In contrast, the \textbf{uniform setting} (left) reveals a more abrupt phase transition: most properties remain near-zero in accuracy until the model reaches a critical capacity threshold, after which multiple properties are rapidly acquired. This implies that, for capacity-constrained models, structuring knowledge injection with a skewed distribution—rather than uniform—may ensure partial acquisition of salient knowledge, rather than risking a complete failure to learn under uniform allocation.



\paragraph{Instruction Fine-Tuning}
\label{para:ift}

We investigate how different instruction tuning strategies affect knowledge retention and acquisition under varying model capacities. We use uniform distribution in this experiment. Based on a shared pretrained model, we inject new knowledge using both Supervised Fine-Tuning (SFT) and Continued Pretraining (CPT), mixing in 80\% of the original pretraining data to mitigate catastrophic forgetting following \cite{chen2020recall}. The underlying pretrained model is trained up to 50\%, 90\%, 100\%, and 120\% of its maximum capacity.



\begin{figure}[htbp]
  \centering
  \includegraphics[width=0.6\linewidth,scale=0.95]{sft.jpg}
  \caption{
    Training loss comparison between SFT and CPT. Due to format differences, SFT initially exhibits a higher loss.
  }
  \label{fig:sft_loss}
\end{figure}

\begin{table}[h]
  \centering
  \caption{
    Accuracy on old and new knowledge after instruction tuning. New knowledge is learned equally well by both methods, but old knowledge is better retained with CPT, especially when the model is near or beyond capacity.
  }
  \label{tab:sft_cpt_accuracy}
  \begin{tabular}{lcccc}
    \toprule
    & \multicolumn{4}{c}{Accuracy on Old Knowledge (\%)} \\
    \cmidrule(lr){2-5}
    & 50\% & 90\% & 100\% & 120\% \\
    \midrule
    SFT (on old) & 100.0 & 93.25 & 89.9 & 78.5 \\
    SFT (on new) & 100.0 & 100.0 & 100.0 & 100.0 \\
    CPT (on old) & 99.6  & 96.5  & 94.5  & 85.1 \\
    CPT (on new) & 100.0 & 100.0 & 100.0 & 100.0 \\
    \bottomrule
  \end{tabular}
\end{table}
While both SFT and CPT achieve perfect accuracy on new knowledge across all settings, they differ in how well they preserve prior knowledge. As shown in Table~\ref{tab:sft_cpt_accuracy}, SFT suffers from more severe forgetting, particularly in models trained at or beyond capacity (100\% and 120\%), where old accuracy drops significantly. This degradation is likely due to format-induced syntactic overhead in SFT, which consumes model capacity and leads to abrupt loss spikes (Figure~\ref{fig:sft_loss}). In contrast, CPT introduces new knowledge more seamlessly, resulting in more stable loss and better retention of previously learned content. Interestingly, for under-capacity models (e.g., 50\%), the difference is negligible—suggesting that forgetting primarily emerges when the model’s capacity is saturated.

\begin{nips}
    \begin{table}[ht]
    \centering
    \caption{Model Configurations with Parameter Counts}
    \begin{tabular}{lcccc}
        \toprule
        \textbf{Model Size} & \textbf{Layers} & \textbf{Heads} & \textbf{Emb Dim} & \textbf{Params (M)} \\
        \hline
        6xs2  & 4  & 4  & 64   & 3.4  \\
        6xs   & 3  & 4  & 80   & 4.3  \\
        5xs   & 3  & 4  & 96   & 5.2  \\
        5xs2  & 3  & 4  & 112  & 6.1  \\
        5xs1  & 3  & 4  & 128  & 7.0  \\
        4xs   & 4  & 4  & 128  & 7.2  \\
        4xs2 & 4  & 4  & 144  & 8.2  \\
        4xs1  & 4  & 4  & 192  & 11.4 \\
        3xs  & 5  & 4  & 192  & 11.9 \\
        3xs2 & 5  & 4  & 224  & 14.3 \\
        3xs1 & 5  & 4  & 256  & 16.8 \\
        xxs  & 6  & 4  & 256  & 17.6 \\
        xxs3 & 6  & 4  & 288  & 20.5 \\
        xxs2 & 6  & 8  & 352  & 26.6 \\
        xxs4 & 6  & 8  & 416  & 33.4 \\
        xxs1 & 6  & 8  & 448  & 37.0 \\
        xs   & 7  & 8  & 448  & 39.4 \\
        xs3  & 8  & 8  & 448  & 41.8 \\
        xs2  & 8  & 8  & 512  & 50.9 \\
        xs1  & 9  & 8  & 512  & 54.1 \\
        s    & 10 & 8  & 512  & 57.2 \\
        s3   & 10 & 8  & 704  & 94.9 \\
        s2   & 10 & 12 & 768  & 109.4 \\
        s1   & 11 & 12 & 768  & 116.5 \\
        base & 12 & 12 & 768  & 123.6 \\
        m5   & 12 & 16 & 896  & 160.7 \\
        m4   & 12 & 16 & 1024 & 202.5 \\
        m3   & 13 & 16 & 1024 & 215.1 \\
        m2   & 14 & 16 & 1024 & 227.7 \\
        m1   & 15 & 16 & 1024 & 240.3 \\
        m    & 16 & 16 & 1024 & 252.8 \\
        \bottomrule
    \end{tabular}
    \label{tab:model_sizes}
\end{table}
\end{nips}

\begin{arxiv}

\begin{table}[ht]
    \centering
    \scalebox{0.8}{
    \begin{tabular}{lcccc}
        \toprule
        \textbf{Model Size} & \textbf{Layers} & \textbf{Heads} & \textbf{Emb Dim} & \textbf{Params (M)} \\
        \hline
        6xs2  & 4  & 4  & 64   & 3.4  \\
        6xs   & 3  & 4  & 80   & 4.3  \\
        5xs   & 3  & 4  & 96   & 5.2  \\
        5xs2  & 3  & 4  & 112  & 6.1  \\
        5xs1  & 3  & 4  & 128  & 7.0  \\
        4xs   & 4  & 4  & 128  & 7.2  \\
        4xs2 & 4  & 4  & 144  & 8.2  \\
        4xs1  & 4  & 4  & 192  & 11.4 \\
        3xs  & 5  & 4  & 192  & 11.9 \\
        3xs2 & 5  & 4  & 224  & 14.3 \\
        3xs1 & 5  & 4  & 256  & 16.8 \\
        xxs  & 6  & 4  & 256  & 17.6 \\
        xxs3 & 6  & 4  & 288  & 20.5 \\
        xxs2 & 6  & 8  & 352  & 26.6 \\
        xxs4 & 6  & 8  & 416  & 33.4 \\
        xxs1 & 6  & 8  & 448  & 37.0 \\
        xs   & 7  & 8  & 448  & 39.4 \\
        xs3  & 8  & 8  & 448  & 41.8 \\
        xs2  & 8  & 8  & 512  & 50.9 \\
        xs1  & 9  & 8  & 512  & 54.1 \\
        s    & 10 & 8  & 512  & 57.2 \\
        s3   & 10 & 8  & 704  & 94.9 \\
        s2   & 10 & 12 & 768  & 109.4 \\
        s1   & 11 & 12 & 768  & 116.5 \\
        base & 12 & 12 & 768  & 123.6 \\
        m5   & 12 & 16 & 896  & 160.7 \\
        m4   & 12 & 16 & 1024 & 202.5 \\
        m3   & 13 & 16 & 1024 & 215.1 \\
        m2   & 14 & 16 & 1024 & 227.7 \\
        m1   & 15 & 16 & 1024 & 240.3 \\
        m    & 16 & 16 & 1024 & 252.8 \\
        \bottomrule
    \end{tabular}
    }
    \caption{Model Configurations with Parameter Counts}
    \label{tab:model_sizes}
\end{table}
\end{arxiv}
% \section{Preliminaries}
% \label{appendix:preliminary}

% \jian{provide references and citations}

% \begin{definition}[Fisher Information]
% Let $X$ be a random variable with probability density function (PDF) $f(x; \theta)$, then the Fisher information is given by:
% \begin{align*}
%     I(\theta) = \mathbb{E}_X \left[ \left( \frac{\partial}{\partial \theta} \log f(X; \theta) \right)^2 \right]=- \mathbb{E}_X \left[  \frac{\partial^2}{\partial \theta^2} \log f(X; \theta)  \right].
% \end{align*}
% \end{definition}




% \begin{lemma}
% \label{lem: Fisher}
%     If $\log f(X; \theta)$ is twice differentiable with respect to $\theta$, we have
%     \[I(\theta)=- \mathbb{E} \left[  \frac{\partial^2}{\partial \theta^2} \log f(X; \theta)  \right]\]
% \end{lemma}
% \begin{proof}[Proof of \Cref{lem: Fisher}]
% \begin{align*}
%     \frac{\partial^2}{\partial \theta^2} \log f(X; \theta) &= \frac{\frac{\partial^2}{\partial \theta^2} f(X; \theta)}{f(X; \theta)} - \left( \frac{\frac{\partial}{\partial \theta} f(X; \theta)}{f(X; \theta)} \right)^2
% \\&= \frac{\frac{\partial^2}{\partial \theta^2} f(X; \theta)}{f(X; \theta)} - \left( \frac{\partial}{\partial \theta} \log f(X; \theta) \right)^2
% \end{align*}

% Taking expectation with respect to $x$ on both sides we know that 

% \begin{align*}
%      \mathbb{E} \left[  \frac{\partial^2}{\partial \theta^2} \log f(X; \theta)  \right] &= \mathbb{E}\left[\frac{\frac{\partial^2}{\partial \theta^2} f(X; \theta)}{f(X; \theta)}\right] - \mathbb{E}\left[\left( \frac{\partial}{\partial \theta} \log f(X; \theta) \right)^2\right]
%      \\ & =  \frac{\partial^2}{\partial \theta^2} \int f(x; \theta) dx  - \mathbb{E}\left[\left( \frac{\partial}{\partial \theta} \log f(X; \theta) \right)^2\right]
%      \\ & = - \mathbb{E}\left[\left( \frac{\partial}{\partial \theta} \log f(X; \theta) \right)^2\right]
% \end{align*}


 


    

% \end{proof}

% \begin{definition}[Cramér-Rao Regular Family]
% \label{def:CR regular}
% A family of probability distributions $\{ f(x; \theta) \}$ parameterized by $\theta \in \Theta \subseteq \mathbb{R}^k$ is said to belong to the \textbf{Cramér-Rao regular family} if it satisfies the following regularity conditions:
% \begin{enumerate}
%     \item \textbf{Differentiability:} The probability density (or mass) function $f(x; \theta)$ is differentiable with respect to the parameter $\theta$:
%     \begin{align*}
%         \frac{\partial}{\partial \theta} f(x; \theta)
%     \end{align*}
%     exists and is well-defined.
%     \item \textbf{Support Condition:} The support of $f(x; \theta)$ does not depend on $\theta$, meaning:
%     \begin{align*}
%         \text{Supp}(f(x; \theta)) \quad \text{is independent of } \theta.
%     \end{align*}
%     \item \textbf{Interchangeability of Differentiation and Integration:} The order of integration and differentiation with respect to $\theta$ can be interchanged:
%     \begin{align*}
%         \frac{\partial}{\partial \theta} \int f(x; \theta) \, dx = \int \frac{\partial}{\partial \theta} f(x; \theta) \, dx = 0.
%     \end{align*}
%     \item \textbf{Finite Fisher Information:} The Fisher information matrix for $\theta$ is finite:
%     \begin{align*}
%         I(\theta) = \mathbb{E} \left[ \left( \frac{\partial}{\partial \theta} \log f(X; \theta) \right)^2 \right] < \infty.
%     \end{align*}
% \end{enumerate}
% \end{definition}

% \begin{theorem}[Cramér-Rao bound]
% Suppose $\{ p_{\theta}(x), \theta \in \Theta \}$ is a Cramer-Rao regular family, where $\Theta \subseteq \mathbb{R}^k$. Its Fisher information $I(\theta)$ is a nonsingular matrix. 
%     Then for any unbiased estimator $\hat{\boldsymbol{\theta}}$ based on $n$ independent samples. Assume $\hat{\boldsymbol{\theta}}$ has finite second moments, the covariance matrix satisfies:
% \begin{align*}
%     \text{Cov}(\hat\theta) \succeq (nI(\theta))^{-1},
% \end{align*}
% \end{theorem}








% \subsection{Properties of Pitman-Yor Mixture Model}
% \begin{theorem}[\cite{pitman2006combinatorial}]
% \label{thm:PYP rank fre} 
% The mixture probability of $\PYMM(\alpha,\beta)$ can be represented as a stick-breaking process:
% \[p_k=W_k\prod\limits_{i=1}^{k-1}(1-W_i),\]
% where $W_i \sim \text{Beta}(1 - \alpha, \beta + i \alpha)$.
% \end{theorem}


\section{Omitted Proofs}
\label{app:Omitted Proofs}
\subsection{Properties of Pitman-Yor Mixture Model}
\label{subsec:PitmanYor}
\begin{lemma}[Theorem 4.3 in \cite{pitman2006combinatorial}]
\label{thm:PYP rank fre} 
The mixing weights $p=(p_1,p_2,\cdots)\sim\PYCRP(\alpha,\beta)$ can be represented as a stick-breaking process:
\[p_j=W_j\prod\limits_{i=1}^{j-1}(1-W_i),\]
where $W_i$ are independent and $W_i \sim \text{Beta}(1 - \alpha, \beta + i \alpha)$.
\end{lemma}
\begin{lemma}[Theorem 3.13 in \cite{pitman2006combinatorial}]
\label{lem:PYCRP fre}
Let \( p = (p_1, p_2, \dots) \sim \mathrm{PYCRP}(\alpha, \beta) \) be the sequence of mixing weights drawn from a Pitman–Yor process. Then, the following limit almost surely exists:
\[
S_{\alpha,\beta} = \lim_{i \to \infty} i^{1/\alpha} p_i.
\]
\end{lemma}


\begin{lemma}
\label{lem:L_j}
Denote the remaining stick length of the stick-breaking process in \Cref{thm:PYP rank fre} as \[\Len_j:=\prod_{i=1}^j(1-W_i)=\sum\limits_{i=j+1}^{\infty}p_i.\]
Then the expectation of the remaining stick length satisfies:
    \[\mathbb{E}[\Len_n]=\Theta(n^{1-1/\alpha}).\]
\end{lemma}
\begin{proof}
Since $W_i \sim \text{Beta}(1 - \alpha, \beta + i \alpha)$, we know that \[\mathbb{E}[W_i]=\frac{\beta + i\alpha}{1 + \beta + (i-1)\alpha}.\]
Moreover, $W_i$ are independent. Thus we have
\begin{align*}
\mathbb{E}[\Len_n] 
&= \prod_{i=1}^{n} \frac{\beta + i\alpha}{1 + \beta + (i-1)\alpha} 
= \frac{\prod_{i=1}^{n} (\beta + i\alpha)}{\prod_{i=1}^{n} (1 + \beta + (i-1)\alpha)} \\
&= \frac{\alpha^n \cdot \prod_{i=1}^{n} \left( \frac{\beta}{\alpha} + i \right)}{\alpha^n \cdot \prod_{i=0}^{n-1} \left( \frac{1 + \beta}{\alpha} + i \right)} 
= \frac{\Gamma\left(n + 1 + \frac{\beta}{\alpha} \right)}{\Gamma\left(1 + \frac{\beta}{\alpha} \right)} \cdot \frac{\Gamma\left( \frac{1 + \beta}{\alpha} \right)}{\Gamma\left( n + \frac{1 + \beta}{\alpha} \right)} \\
&= C \cdot \frac{\Gamma\left(n + 1 + \frac{\beta}{\alpha} \right)}{\Gamma\left( n + \frac{1 + \beta}{\alpha} \right)}, \quad \text{where } C = \frac{\Gamma\left( \frac{1 + \beta}{\alpha} \right)}{\Gamma\left(1 + \frac{\beta}{\alpha} \right)}.
\end{align*}
Then by Stirling's approximation, we know that that $\mathbb{E}(\Len_j)=\Theta(n^{1-1/\alpha}).$
\end{proof}
\subsection{KL-divergence of Mixtures}
\begin{lemma}
\label{lem:mixture KL}
Let $\phi=\sum\limits_{i=1}^{n}p_i\delta_{\phi_i},\theta=\sum\limits_{i=1}^{n}q_i\delta_{\theta_i}$
be two mixtures. Then we have
\[\DKL(P_{\phi}(x)||P_{\theta}(x))\le \sum\limits_{i=1}^np_i\log \dfrac{p_i}{q_i}+\sum\limits_{i=1}^np_i \DKL(P_{\phi_i}(x)||P_{\theta_i}(x)).\]
\end{lemma}
\begin{proof}
By the log-sum inequality, we know that 
\begin{align*}
    \DKL(P_{\phi}(x)||P_{\theta}(x)) &= \sum\limits_{x} \left(\sum_{i=1}^n p_i P_{\phi_i}(x) \right) \log\dfrac{ \sum_{i=1}^n p_i P_{\phi_i}(x) }{ \sum_{i=1}^nP_{\theta_i}(x)}
    \\& \le \sum\limits_{x} \sum_{i=1}^n \left(p_i P_{\phi_i}(x)  \log\dfrac{ p_iP_{\phi_i}(x) }{ q_iP_{\theta_i}(x)}\right)
    \\&= \sum\limits_{i=1}^n p_i\log\dfrac{p_i}{q_i}+\sum\limits_{i=1}^np_i \DKL(P_{\phi_i}(x)||P_{\theta_i}(x))
\end{align*}
\end{proof}

\subsection{KL-divergence and Fisher Information}
\begin{definition}[Fisher Information]
Let $X$ be a random variable with probability density function (PDF) $P_{\theta}(x)$. The Fisher information is given by:
\begin{align*}
    J(\theta) = \mathbb{E}_X \left[ \left( \frac{\partial}{\partial \theta} \log P_{\theta}(X) \right)^2 \right]=- \mathbb{E}_X \left[  \frac{\partial^2}{\partial \theta^2} \log P_{\theta}(X) \right].
\end{align*}
\end{definition}
\begin{lemma}[Exercise 11.7 in \cite{cover2006elements}]
\label{lem:KL and MSE}
For a parametric family $\{P_\theta(x)\}$ that
\[
\lim_{\theta' \to \theta} \frac{\DKL(P_\theta \| P_{\theta'})}{(\theta - \theta')^TJ(\theta)(\theta-\theta')}  = \frac{1}{2}.
\]
\end{lemma}
\begin{proof}
We provide a proof for completeness. We perform a second-order Taylor expansion of the KL divergence with respect to its second argument \( \theta' \), around the  first term \( \theta \). 

\begin{align*}
    D_{\text{KL}}(P_{\theta}||P_{\theta'}) &= \mathbb{E}_{x\sim P_{\theta}}\left[\log \dfrac{P_{\theta}(x)}{P_{\theta'}(x)}\right]
    \\ &=-\mathbb{E}_{x\sim P_{\theta}}\left[(\theta-\theta')\nabla_{\theta} \log P_{\theta} + \dfrac{1}{2}(\theta-\theta')^T\nabla^2_{\theta} \log P_{\theta}(\theta-\theta') \right] + o(|\theta-\theta'|^2)
    \\ & = -\mathbb{E}_{x\sim P_{\theta}}\left[(\theta-\theta')\nabla_{\theta} \log P_{\theta}\right] - \dfrac{1}{2}\mathbb{E}_{x\sim P_{\theta}}\left[(\theta-\theta')^T\nabla^2_{\theta} \log P_{\theta}(\theta-\theta') \right] + o(|\theta-\theta'|^2)
    \\ & = -(\theta-\theta')\mathbb{E}_{x\sim P_{\theta}}\left[\nabla_{\theta} \log P_{\theta}\right] - \dfrac{1}{2}(\theta-\theta')^T\mathbb{E}_{x\sim P_{\theta}}\left[\nabla^2_{\theta} \log P_{\theta} \right](\theta-\theta') + o(|\theta-\theta'|^2)
    \\ & = -(\theta-\theta')\mathbb{E}_{x\sim P_{\theta}}\left[\nabla_{\theta} \log P_{\theta}\right] + \dfrac{1}{2}(\theta-\theta')^TJ(\theta)(\theta-\theta') + o(|\theta-\theta'|^2)
\end{align*}

The first-order term $(\theta-\theta')\mathbb{E}_{x\sim P_{\theta}}\left[\nabla_{\theta} \log P_{\theta}\right]=0$ since the KL divergence attains its minimum value of zero at \( \theta' = \theta \).

 % \begin{align*}
 %     \mathbb{E}_{x\sim P_{\theta}}\left[(\theta-\theta')^T\nabla^2_{\theta} \log P_{\theta}(\theta-\theta') \right]&= \mathbb{E}_{x\sim P_{\theta}}\left[(\theta-\theta')^T\nabla^2_{\theta} \log P_{\theta}(\theta-\theta') \right]
 %     \\ &= (\theta-\theta')^T\mathbb{E}_{x\sim P_{\theta}}\left[\nabla^2_{\theta} \log P_{\theta} \right](\theta-\theta')
 %     \\ &=
 %     % \\ &= tr\left((\theta-\theta')^T(\theta-\theta')\mathbb{E}_{x\sim P_{\theta}}\left[\nabla^2_{\theta} \log P_{\theta} \right]\right)
 %     % \\ & =tr((\theta-\theta')^T(\theta-\theta')I(\theta))
 % \end{align*}

% \begin{align*}
%     \mathbb{E}_{\theta'}\left[D_{\text{KL}}(P_{\theta}||P_{\theta'})\right]& = \mathbb{E}_{\theta'}\left[(\theta-\theta')\mathbb{E}_{x\sim P_{\theta}}\left[\nabla_{\theta} \log P_{\theta}\right] \right]+\mathbb{E}_{\theta'}\left[tr((\theta-\theta')^T(\theta-\theta')I(\theta)) \right] + o(|\theta-\theta'|^2)
%     \\ &= tr(\Sigma I(\theta))+o(tr(\Sigma))
% \end{align*}
% Since $I(\theta)$ is nonsingular, so it is positive definite. Thus $tr(\Sigma I(\theta))=\Omega(tr(\Sigma))$
Thus we know that 
\[
\lim_{\theta' \to \theta} \frac{\DKL(P_\theta \| P_{\theta'})}{(\theta - \theta')^TJ(\theta)(\theta-\theta')}  = \frac{1}{2}.
\]

\end{proof}
\begin{remark}
\label{justify KL and MSE}
The conclusion of this lemma indicates that the  \Cref{ass: mutual information} is essentially analogous to requiring the Fisher Information of the distribution family to be bounded. Such assumptions are common in related literature. For instance, Conditions 1 and 4 in \citep{clarke1990information} impose similar requirements. Notably, \citep{clarke1990information} also points out that these assumptions are widely satisfied within a class of exponential family distributions. 
\end{remark}

% Recall \Cref{loss dec} \begin{align*}
%     L^{(n)} &= \mathbb{E}\left[-\sum\limits_{i=1}^{l}\log P_{\theta_S^{(n)}}(v'_{i}|v'_{0:i-1})\right]
%       + \mathbb{E}\left[-\log P_{\theta_K^{(n)}}(y_n|x_{z_n}) \right] \nonumber
%      \\& = H(P_{\theta^*}) + D_{\text{KL}}(P_{\theta^*}\|P_{\theta_S^{(n)}}) \nonumber
%      \\& + \mathbb{E}_{z_n} \left[H(P_{\theta_{x_{z_n}}})\right] + \mathbb{E}_{z_n} \left[D_{\text{KL}}(P_{\theta_{x_{z_n}}}\|P_{\theta_K^{(n)}}(\cdot|x_{z_n})) \right]
% \end{align*}
% \begin{proof}[Proof of \Cref{loss dec}]
%     \begin{align*}
%     L^{(n)} &= \mathbb{E}\left[-\sum\limits_{i=1}^l\log P^{(n)}(v_{i}|v_{0:i-1}) \right] \\
%      & = \mathbb{E}\left[-\sum\limits_{i=1,i\ne j}^{l}\log P^{(n)}(v_{j}|v_{0:i-1})\right]
%      \\& + \mathbb{E}\left[-\log P^{(n)}(v_j|v_{0:j-1}) \right] \\
%      & = \mathbb{E}\left[-\sum\limits_{i=1}^{l}\log P_{\theta_S^{(n)}}(v_{i}|v_{0:i-1})\right]
%      \\& + \mathbb{E}\left[-\log P_{\theta_S^{(n)}}(v_y|v_0\cdots v_{l-1}) \right]
%      \\& + \mathbb{E}\left[-\log P_{\theta_K^{(n)}}(y_n|x_{z_n}) \right]
%      \\& = \mathbb{E}\left[-\sum\limits_{i=1}^{l}\log P_{\theta_S^{(n)}}(v'_{i}|v'_{0:i-1})\right]
%      \\& + \mathbb{E}\left[-\log P_{\theta_K^{(n)}}(y_n|x_{z_n}) \right]
%      \\& = H(P_{\theta^*}) + D_{\text{KL}}(P_{\theta^*}\|P_{\theta_S^{(n)}}) 
%      \\& + \mathbb{E}_{z_n} \left[H(P_{\theta_{x_{z_n}}})\right] + \mathbb{E}_{z_n} \left[D_{\text{KL}}(P_{\theta_{x_{z_n}}}\|P_{\theta_K^{(n)}}(\cdot|x_{z_n})) \right]
% \end{align*}
% \end{proof}


% \subsection{Proofs of \Cref{sec:loss}}
% \begin{lemma}
% \label{lem: KL}
%      Let $\hat{\theta}$ be an unbiased estimator of $\theta$ and $\text{Cov}(\hat\theta)=\Sigma$. Assuming that the Fisher information $I(\theta)$ is a nonsingular matrix, then we have
%     \[\lim\limits_{\Sigma\rightarrow 0} \dfrac{\mathbb{E}_{\hat\theta}\left[D_{\text{KL}}(P_{\theta}||P_{\hat{\theta}})\right]}{tr(\Sigma I(\theta))}=\dfrac{1}{2}\]

% \end{lemma}
% \begin{proof}[Proof of \Cref{lem: KL}]
% \begin{align*}
%     D_{\text{KL}}(P_{\theta}||P_{\hat{\theta}}) &= \mathbb{E}_{x\sim P_{\theta}}\left[\log \dfrac{P_{\theta}(x)}{P_{\hat\theta}(x)}\right]
%     \\ &=\mathbb{E}_{x\sim P_{\theta}}\left[(\theta-\hat\theta)\nabla_{\theta} \log P_{\theta} + \dfrac{1}{2}(\theta-\hat\theta)^T\nabla^2_{\theta} \log P_{\theta}(\theta-\hat\theta) \right] + o(|\theta-\hat\theta|^2)
%     \\ & = \mathbb{E}_{x\sim P_{\theta}}\left[(\theta-\hat\theta)\nabla_{\theta} \log P_{\theta}\right] + \dfrac{1}{2}\mathbb{E}_{x\sim P_{\theta}}\left[(\theta-\hat\theta)^T\nabla^2_{\theta} \log P_{\theta}(\theta-\hat\theta) \right] + o(|\theta-\hat\theta|^2)
% \end{align*}

%  \begin{align*}
%      \mathbb{E}_{x\sim P_{\theta}}\left[(\theta-\hat\theta)^T\nabla^2_{\theta} \log P_{\theta}(\theta-\hat\theta) \right]&= \mathbb{E}_{x\sim P_{\theta}}\left[tr((\theta-\hat\theta)^T\nabla^2_{\theta} \log P_{\theta}(\theta-\hat\theta) )\right]
%      \\ &= \mathbb{E}_{x\sim P_{\theta}}\left[tr((\theta-\hat\theta)^T(\theta-\hat\theta)\nabla^2_{\theta} \log P_{\theta} )\right]
%      \\ &= tr\left((\theta-\hat\theta)^T(\theta-\hat\theta)\mathbb{E}_{x\sim P_{\theta}}\left[\nabla^2_{\theta} \log P_{\theta} \right]\right)
%      \\ & =tr((\theta-\hat\theta)^T(\theta-\hat\theta)I(\theta))
%  \end{align*}

% \begin{align*}
%     \mathbb{E}_{\hat\theta}\left[D_{\text{KL}}(P_{\theta}||P_{\hat{\theta}})\right]& = \mathbb{E}_{\hat\theta}\left[(\theta-\hat\theta)\mathbb{E}_{x\sim P_{\theta}}\left[\nabla_{\theta} \log P_{\theta}\right] \right]+\mathbb{E}_{\hat\theta}\left[tr((\theta-\hat\theta)^T(\theta-\hat\theta)I(\theta)) \right] + o(|\theta-\hat\theta|^2)
%     \\ &= tr(\Sigma I(\theta))+o(tr(\Sigma))
% \end{align*}
% Since $I(\theta)$ is nonsingular, so it is positive definite. Thus $tr(\Sigma I(\theta))=\Omega(tr(\Sigma))$

% Then we know that \[\lim\limits_{\Sigma\rightarrow 0} \dfrac{\mathbb{E}_{\hat\theta}\left[D_{\text{KL}}(P_{\theta}||P_{\hat{\theta}})\right]}{tr(\Sigma I(\theta))}=\dfrac{1}{2}.\]

% \end{proof}


% \begin{theorem}[Restatement of \Cref{thm:R_S lower bound}]
%     Let  $\hat\theta$ be an consistent and unbiased estimator of $\theta$ based on the observations $T_0,T_1,\cdots,T_{n-1}$ where $T_i$ are i.i.d. drawn from $P_{\theta^*}, \theta^*\in \mathbb{R}^d$. Then we have
%     \[ \lim\limits_{n\rightarrow \infty} 2n\mathbb{E}[D_{\text{KL}}(P_{\theta^*}||P_{\hat\theta})]\ge d.\]
% \end{theorem}
% \begin{proof}[Proof of \Cref{thm:R_S lower bound}]
%     We use Cramér-Rao bound and \Cref{lem: KL} to prove this theorem.

%     Assume that $\text{Cov}(\hat\theta)=\Sigma$. 

%     Since $\hat\theta$ is an consistent estimator, we know that $\Sigma\rightarrow 0$ when $n\rightarrow \infty$.

%     Thus we can apply the conclusion of \Cref{lem: KL}.

%     \[ \lim\limits_{n\rightarrow \infty} \dfrac{\mathbb{E}[D_{\text{KL}}(P_{\theta^*}||P_{\hat\theta})]}{tr(\Sigma I(\theta^*))} = \dfrac{1}{2}.\]

%     By Cramér-Rao bound, we know that 
%     \[\Sigma \succeq (nI(\theta^*))^{-1}. \]

%     \[tr(\Sigma I(\theta^*))\ge tr(I(\theta^*)^{-1}I(\theta^*))=d\]

%     Thus we have\[ \lim\limits_{n\rightarrow \infty} 2n\mathbb{E}[D_{\text{KL}}(P_{\theta^*}||P_{\hat\theta})]\ge d.\]
% \end{proof}



% \begin{theorem}[Restatement of \Cref{thm: MLE}]
%     Let  $\hat\theta$ be the MLE estimator of $\theta^*$ based on the observations $T_1,T_2,\cdots,T_{n-1}$ where $T_i$ are i.i.d. drawn from $P_{\theta^*}, \theta^*\in \mathbb{R}^d$. Then we have
%     \[\lim\limits_{n\rightarrow \infty}2n \mathbb{E}[D_{\text{KL}}(P_{\theta^*}||P_{\hat{\theta}})]=d.\]
% \end{theorem}
% \begin{proof}[Proof of \Cref{thm: MLE}]
%     By asymptotic normality of MLE, we know that 
% \[
% \sqrt{n}(\hat{\theta} - \theta^*) \xrightarrow{d} \mathcal{N}(0, I(\theta^*)^{-1}),
% \]
% Reusing the logic of \Cref{lem: KL}, we know that 
% \[ \lim\limits_{n\rightarrow \infty} \dfrac{n\mathbb{E}[D_{\text{KL}}(P_{\theta^*}||P_{\hat\theta})]}{tr(I(\theta^*)^{-1} I(\theta^*))} = \dfrac{1}{2}.\]
% Thus
% \[\lim\limits_{n\rightarrow \infty}2n \mathbb{E}[D_{\text{KL}}(P_{\theta^*}||P_{\hat{\theta}})]=d.\]
% \end{proof}


% \begin{theorem}[Restatement of \Cref{thm: R_K bound}]
% Assuming that for any unseen $x, D_{\text{KL}}(P_{\theta_x}\|P_{\theta_K^{(n)}}(\cdot|x))\ge C_1$ for some constant $C_1$. Then we have
%     \begin{align*}
%         \mathbb{E}_{z_{0:n}}\left[R_K^{(n)}\right]&=\mathbb{E}_{z_{0:n}} \left[D_{\text{KL}}(P_{\theta_{x_{z_n}}}\|P_{\theta_K^{(n)}}(\cdot|x_{z_n})) \right]
%         \\&\ge \Omega\left(\dfrac{1}{n^{\beta/1+\beta}}\right).
%     \end{align*}
% \end{theorem}
% \begin{proof}[Proof of \Cref{thm: R_K bound}]
%     Recall that $z_i$ represents the index of the table $i-\text{th}$ customer sits. Denote $K(n)$ as the number of occupied tables in PYCRP when $n+1-\text{th}$ customer comes.

%     \begin{align*}
%         \mathbb{E}_{z_{0:n}} \left[D_{\text{KL}}(P_{\theta_{x_{z_n}}}\|P_{\theta_K}^{(n)}(\cdot|x_{z_n}))\right] &= \mathbb{E}_{z_{0:n-1}}\left[\mathbb{E}_{z_n} \left[D_{\text{KL}}(P_{\theta_{x_{z_n}}}\|P_{\theta_K}^{(n)}(\cdot|x_{z_n}))\right]\right]
%         \\ &\ge C_1\mathbb{E}_{z_{0:n-1}}\left[\text{Pr}(x_{z_n} \text{ was unseen})\right]
%         \\ & = C_1\mathbb{E}_{z_{0:n-1}}\left[\dfrac{\alpha+\beta K(n)}{n+\alpha}\right]
%         \\ & = \dfrac{C_1(\alpha+\beta\mathbb{E}_{z_{0:n-1}}\left[ K(n)\right])}{n+\alpha}
%     \end{align*}

% Notice that 
% \begin{align*}
%     \mathbb{E}_{z_{0:n-1}}\left[K(n)\right]&=\mathbb{E}_{z_{0:n-2}}\left[K(n-1)+\dfrac{\alpha+\beta K(n-1)}{n-1+\alpha}\right]
%     \\ & = \dfrac{\alpha}{n-1+\alpha} + \dfrac{n-1+\alpha+\beta}{n-1+\alpha}\mathbb{E}_{z_{0:n-2}}\left[K(n-1) \right]
%     \\ & =\cdots
%     \\ & = \Theta(\log n) + \prod\limits_{k=1}^{n-1} \dfrac{k+\alpha+\beta}{k+\alpha}
%     \\ & = \Theta(n^\beta)
% \end{align*}

% Thus we know that \begin{align*}
%         \mathbb{E}_{z_{0:n}}\left[R_K^{(n)}\right]\ge \Omega\left(\dfrac{1}{n^{1-\beta}}\right).
%     \end{align*}
% \end{proof}


\subsection{\texorpdfstring{Proofs of \Cref{sec:baye data law}}{Proofs of \ref{sec:baye data law}}}
\begin{theorem}[Restatement of \Cref{thm:data_scaling_Mutual_Information}]
Under the Bayesian sequential prediction framework and \Cref{ass: mutual information}, the averaged optimal Bayesian redundancy (per sentence) of the hierarchical data model $\dataModel$  satisfies:
    \begin{align*}
    % \mathcal{L}_{N,S}:&=\dfrac{\mathbb{I}(T_{1:N}; \synModel)}{N} = \Tilde{O}(d_SN^{-1}),\\
    \inf_{\Model}\dfrac{1}{N}\redundency_N(\Model,\datafamily)=\dfrac{1}{N}\mathbb{I}(X_{1:N}; \dataModel) = 
    \widetilde{O}\left(\frac{\knwdimension}{N^{1-\alpha}}+\frac{ n_s\syndimension}{N}\right).
    \end{align*}
\end{theorem}
\begin{proof} 
% We first present a lemma that provides an upper bound on the mutual information using the KL-covering number in \cite{duchi2024lecture}.
% We first define the KL-covering number and the $L_2$-norm covering number
We use the {\em index of resolvability} approach 
similar to \citep{barron1991minimum} to prove the theorem, by constructing a covering of the parameter space in terms of 
KL-divergence. 
However, we cannot directly apply the conclusion, as we are unable to explicitly construct a KL-divergence covering for $\dataprobability$. Instead, we can construct a set of probability distributions $\mathbb{Q}$ such that, for every specific $\dataModel \in \datafamily$, there exists a $Q \in \mathbb{Q}$ satisfying $\DKL(\dataModel \,\|\, Q) \leq f(\dataModel)$, where $f(\dataModel)$ is a value dependent on $\dataModel$.



Define the KL-covering number of the probability family $\mathcal{P}$ and the $L_2$ norm covering number of the parameter space $\Phi$ as:
\[
N_{\text{kl}}(\epsilon, \mathcal{P}) := \inf \left\{ n \in \mathbb{N} \mid \exists Q_i, i = 1, \ldots, n, \sup_{P \in \mathcal{P}} \min_i D_{\text{kl}}(P \| Q_i) \leq \epsilon^2 \right\}.
\]



% \begin{lemma}[Lemma ??? in \cite{duchi2024lecture}]
% Assume that \( X_1, \ldots, X_n \) are drawn i.i.d. from \( P_v \) conditional on \( V = v \). Let \( N_{\text{kl}}(\epsilon, \mathcal{P}) \) denote the KL-covering number of a collection \( \mathcal{P} \) containing the distributions (over a single observation) \( P_v \) for all \( v \in \mathcal{V} \). Then
% \[
% I(V; X_1, \ldots, X_n) \leq \inf_{\epsilon \geq 0} \{ n \epsilon^2 + \log N_{\text{kl}}(\epsilon, \mathcal{P}) \}.
% \]
% \end{lemma}


% We first prove $\dfrac{\mathbb{I}(T_{1:N}; \synModel)}{N} = \Tilde{O}(d_SN^{-1})$.

% Define the $L_2$ norm covering number of the parameter space $\Phi$ as:
\[
N_{L_2}(\epsilon, \Phi) := \inf \left\{ n \in \mathbb{N} \mid \exists \phi_i, i = 1, \ldots, n, \sup_{\phi \in \Phi} \min_i \|\phi-\phi_i\|_2\le \epsilon \right\}.
\]

We first construct KL-divergence coverings for the probability families $\synprobability$ and $\knwprobability$, as well as a KL-divergence covering for the $n$-dimensional discrete simplex. We subsequently construct a point-wise approximate covering set $\mathbb{Q}$ for $\dataprobability$ based on these coverings.
By \Cref{ass: mutual information}, we know that 
$$
N_{\text{kl}}(\epsilon, P_{\synfamily})\le N_{L_2}(\epsilon/\synLip, \synfamily) \quad \text{ and }\quad 
N_{\text{kl}}(\epsilon, P_{\knwfamily})\le N_{L_2}(\epsilon/\knwLip, \knwfamily).
$$
Then we consider the $L_2$ norm covering number of the parameter space $\synfamily$. 
We simply divide each coordinate into \( \lceil \frac{2\synLip\sqrt{\syndimension}}{ \epsilon} \rceil \) equal parts. It is easy to see that this results in an \( L_2 \) covering of the parameter space with a resolution of approximately \( \frac{\epsilon}{\synLip} \).
 Thus we know that 
\[
N_{\text{kl}}(\epsilon, P_{\boldsymbol{\synfamily}})\le N_{L_2}(\epsilon/\synLip, \synfamily)\le \left( \dfrac{2\synLip\sqrt{\syndimension}}{\epsilon} \right) ^{\syndimension}.
\]
Similarly, we know that 
\[
N_{\text{kl}}(\epsilon, P_{\boldsymbol{\knwfamily}})\le \left( \dfrac{2\knwLip\sqrt{\knwdimension}}{\epsilon} \right) ^{\knwdimension}.
\]
Next, we consider a KL-covering of the \( n \)-dimensional simplex.

Denote $\mathcal{P}_n=\{(p_1,\cdots,p_n);p_i\ge 0 ,\sum\limits_{i=1}^np_i=1\}$.
Consider the discretization of the simplex $\mathcal{Q}_{n,\epsilon}=\{(q_1,\cdots,q_n); \lfloor n/\epsilon^2 \rfloor q_i  \in \mathcal{N}^+,\sum\limits_{i=1}^n q_i=1\}$. 

We prove that $\mathcal{Q}_{n,\epsilon}$ is an $\epsilon^2$-KL-covering 
Then for any $(p_1,p_2,\cdots,p_n)\in \mathcal{P}_n$, there exists $(q_1,\cdots,q_n)\in \mathcal{Q}_{n,\epsilon}$ such that $|p_i-q_i|\le \dfrac{\epsilon^2}{n}$.
Therefore, we know that \begin{align*}
    \sum\limits_{i=1}^n p_i\log \dfrac{p_i}{q_i}&\le \sum\limits_{i=1}^n p_i\dfrac{p_i-q_i}{q_i}
   =\sum\limits_{i=1}^n \dfrac{(p_i-q_i)^2}{q_i}
    \le \sum\limits_{i=1}^n \dfrac{\epsilon^2}{n} \le \epsilon^2.
\end{align*}
Note that a simple upper bound for $|\mathcal{Q}_{n,\epsilon}|$ is $|\mathcal{Q}_{n,\epsilon}|\le \lfloor n/\epsilon^2\rfloor^n$.
% Denote $Q_{\text{syn}}=\dfrac{1}{\syncover}\sum\limits_{i=1}^{\syncover}\left(Q_{\text{syn}}^{(i)}\right)^N$ and $Q_{\text{knw}}=\dfrac{1}{\knwcover}\sum\limits_{i=1}^{\knwcover}\left(Q_{\text{knw}}^{(i)}\right)^N$.

For notational simplicity, we denote
\[
\syncover = \left(\frac{2\synLip\sqrt{\syndimension}}{\epsilon}\right)^{\syndimension}, \quad 
\knwcover = \left(\frac{2\knwLip\sqrt{\knwdimension}}{\epsilon}\right)^{\knwdimension},
\]
as upper bounds on the $\epsilon-$KL-covering numbers of $\synprobability$ and $\knwprobability$, respectively.

Assume that distributions 
\[
\{Q_{\text{syn}}^{(1)}(\Synelem), \dots, Q_{\text{syn}}^{(\syncover)}(\Synelem)\}
\] 
form an $\epsilon$-KL-covering of $\mathcal{P}_{\synfamily}$, and distributions 
\[
\{Q_{\text{knw}}^{(1)}(\Knwelem), \dots, Q_{\text{knw}}^{(\knwcover)}(\Knwelem)\}
\] 
form an $\epsilon$-KL-covering of $\mathcal{P}_{\knwfamily}$.
Let $m = \lfloor N^{\alpha} \rfloor$ denote the truncation point used subsequently in the truncated estimation of $\knwprobability$.
Then we construct the covering of $\dataprobability$.
Denote by $\mathcal{Q}$ the set of joint distributions over knowledge element $\Knwelem$ and syntax element $\Synelem$.
% which form is a KL-covering of $\mathcal{P}_{\boldsymbol{\Phi}_{\text{data}}}$. 
A distribution $Q(\Knwelem,\Synelem)\in \mathcal{Q}$ if and only if it satisfies all of the following conditions:

\begin{itemize}
    \item There exist $q = (q_1,\dots,q_{m+1})\in \mathcal{Q}_{m+1,\epsilon}$ and indices $k_1,\dots,k_{m}$, such that
    \[
    Q(\Knwelem)=\sum\limits_{i=1}^{m} q_i Q_{\text{knw}}^{(k_i)}+q_{m+1}Q_u,
    \]
    where $Q_u$ is the uniform distribution over the support of knowledge element $\knwset$.

    \item There exist indices $s_1,\dots,s_{n_s}$, such that for each $i\in\{1,\dots,n_s\}$,
    \[
    Q(\Synelem\mid \Knwelem\in \mathbb{K}_i)=Q_{\text{syn}}^{(s_i)},
    \]
    where $\mathbb{K}_i$ is the set of knowledge corresponding to the syntax model $\synModel^{(i)}$.
\end{itemize}


Next we prove that, for any $P_{\dataModel}\in \mathcal{P}_{\datafamily}$, we have that
\[
\min_{Q\in \mathcal{Q}}\DKL(P_{\dataModel}(\Knwelem,\Synelem)||Q(\Knwelem,\Synelem))\le \Len_m\log \mathbb{|K|}+3\epsilon^2,
\]
where $\Len_j=\sum\limits_{i=j+1}^{\infty} p_j$ is the remaining stick length of $\knwModel$. (See \Cref{lem:L_j}.)


% In order to prove this conclusion, we need the following lemma:
% \begin{lemma}[KL Divergence of Mixture Models]
% \label{lem:mixture KL}
% Let $\phi=\sum\limits_{i=1}^{n}p_i\delta_{\phi_i},\theta=\sum\limits_{i=1}^{n}q_i\delta_{\theta_i}$. Then for all $\theta$, we have
% \[\DKL(P_{\phi}(x)||P_{\theta}(x))\le \sum\limits_{i=1}^np_i\log \dfrac{p_i}{q_i}+\sum\limits_{i=1}^np_i \DKL(P_{\phi_i}(x)||P_{\theta_i}(x)).\]
% \end{lemma}
% \begin{proof}
% Denote $z$ as the component indicator variable: $P(z=k)=p_k$.

% Then by the chain rule of KL-divergence, we know that
% \begin{align*}
%     \DKL(P_{\Phi}(x)||P_{\Theta}(x))&=\DKL(P_{\Phi}(x,z)||P_{\Theta}(x,z))-\mathbb{E}_{x}\DKL(P_{\Phi}(z|x)||P_{\Theta}(z|x))
%     \\ & \le \DKL(P_{\Phi}(x,z)||P_{\Theta}(x,z))
%     \\& = \DKL(P_{\Phi}(z)||P_{\Theta}(z)) + \mathbb{E}_{z}\DKL(P_{\Phi}(x|z)||P_{\Theta}(x|z))
%     \\& = \sum\limits_{k=1}^{n}p_k\log \dfrac{p_k}{q_k}+\sum\limits_{k=1}^{n}p_k \DKL(P_{\phi_k}(x)||P_{\theta_k}(x)).
% \end{align*}
% By the log-sum inequality, we know that 
% \begin{align*}
%     \DKL(P_{\phi}(x)||P_{\theta}(x)) &= \sum\limits_{x} \left(\sum_{i=1}^n p_i P_{\phi_i}(x) \right) \log\dfrac{ \sum_{i=1}^n p_i P_{\phi_i}(x) }{ \sum_{i=1}^nP_{\theta_i}(x)}
%     \\& \le \sum\limits_{x} \sum_{i=1}^n \left(p_i P_{\phi_i}(x)  \log\dfrac{ p_iP_{\phi_i}(x) }{ q_iP_{\theta_i}(x)}\right)
%     \\&= \sum\limits_{i=1}^n p_i\log\dfrac{p_i}{q_i}+\sum\limits_{i=1}^np_i \DKL(P_{\phi_i}(x)||P_{\theta_i}(x))
% \end{align*}
% \end{proof}
We rewrite $\knwModel$ as:
\[
\knwModel=\sum\limits_{i=1}^m p_i\delta_{\phi_i}+\Len_m\delta_{\phi^{\prime}},
\]
where $\phi'=\dfrac{1}{\Len_m}\sum\limits_{i=m+1}^{\infty}p_i\delta_{\phi_i}$.
By the definition of $\mathcal{Q}$, there exists a distribution $Q^*\in\mathcal{Q}$ such that
\[
Q^*(\Knwelem)=\sum\limits_{i=1}^{m} q^*_i Q_{\text{knw}}^{(k^*_i)}+q_{m+1}^* Q_u,\quad
Q^*(\Synelem\mid \Knwelem\in \mathbb{K}_i)=Q_{\text{syn}}^{(s^*_i)},
\]
and $Q^*$ satisfies the following three conditions:
\begin{itemize}
    \item $\DKL(P_{\phi_i}(\Knwelem)||Q_{\text{knw}}^{(k^*_i)}(\Knwelem))\le \epsilon^2$ for all $1\le i\le m$.
    \item $\sum\limits_{i=1}^{m}p_i\log \dfrac{p_i}{q^*_i}+\Len_m\log\dfrac{\Len_m}{q^*_{m+1}}\le \epsilon^2.$
    \item $\DKL(P_{\synModel^{(i)}}(\Synelem)||Q_{\text{syn}}^{(s^*_i)}(\Synelem))\le \epsilon^2$ for all $1\le i\le n_s$.
\end{itemize}
% Since $Q_{\text{knw}}^{(1)}(K),\cdots,Q_{\text{knw}}^{(\knwcover)}(K)$ is a 
% $\epsilon$ KL-covering of $\mathcal{P}_{\knwfamily}$, we know that there exists $k_1,\cdots,k_m$ such that $\DKL(P_{\phi_i}(K)||Q_{\text{knw}}^{(k_i)}(K))\le \epsilon$ for all $1\le i\le m$.

% Since $\mathcal{Q}_{m+1,\epsilon}$ is a $\epsilon$-KL-covering of $\mathcal{P}_{m+1}$, we know that there exists $q=(q_1,\cdots,q_{m+1})\in \mathcal{Q}_{m+1,\epsilon}$ such that $\sum\limits_{i=1}^{m}p_i\log \dfrac{p_i}{q_i}+L_m\log\dfrac{L_m}{q_{m+1}}\le \epsilon.$

% Since $Q_{\text{syn}}^{(1)}(S),\cdots,Q_{\text{syn}}^{(\syncover)}(S)$ is a $\epsilon$-KL-covering of $\mathcal{P}_{\synfamily}$, we know that there exists $s_1,\cdots,s_{n_s}$ such that $\DKL(P_{\synModel^{(i)}}(S)||Q_{\text{syn}}^{(s_i)}(S))\le \epsilon$ for all $1\le i\le n_s$.


Then by the chain rule of KL-divergence and \Cref{lem:mixture KL}, we know that 
\begin{align*}
    \DKL(P_{\dataModel}(\Knwelem,\Synelem)||Q^*(\Knwelem,\Synelem))=&\DKL(P_{\knwModel}(\Knwelem)||Q^*(\Knwelem))+ \mathbb{E}_K\left[\DKL\left(P_{\synModel}(\Synelem|\Knwelem)||Q^*(\Synelem|\Knwelem)\right)\right]
    \\ 
    \le& \sum\limits_{i=1}^m p_i\log \dfrac{p_i}{q^*_i}+\Len_m\log\dfrac{\Len_m}{q^*_{m+1}}+\sum\limits_{i=1}^mp_i\DKL(P_{\phi_i}(\Knwelem)||Q_{\text{knw}}^{(k^*_i)}(\Knwelem)) \\
   + &\Len_m\DKL(P_{\phi'}(\Knwelem)||Q_u(\Knwelem)) 
    + \sum\limits_{i=1}^{n_s}P(\Knwelem\in \mathbb{K}_i)\DKL(P_{\synModel^{(i)}}(\Synelem)||Q_{\text{syn}}^{(s^*_i)}(\Synelem))
    \\  
    \le & \epsilon^2+ \sum\limits_{i=1}^m p_i\epsilon^2 + \Len_m\log |\mathbb{K}|+\sum\limits_{i=1}^{n_s}P(\Knwelem\in \mathbb{K}_i)\epsilon^2
    \\ \le & 3\epsilon^2+\Len_m\log |\mathbb{K}| 
\end{align*}
Subsequently, we derive an upper bound for the associated covering number. We know that $|\mathcal{Q}|\le |\mathcal{Q}_{m+1,\epsilon}|\knwcover^m\syncover^{n_s}$. That is to say
\begin{align*}
    \log |\mathcal{Q}|&\le \log |\mathcal{Q}_{m+1,\epsilon}|+m\log \knwcover+n_s\log\syncover
    \\ &\le 2 m\log \dfrac{m}{\epsilon^2}+m\knwdimension\log\dfrac{\knwLip\knwdimension}{\epsilon}+n_s\syndimension\log\dfrac{\synLip\syndimension}{\epsilon}.
\end{align*}
Finally, we prove the conclusion in this theorem.
Consider $Q_0=\dfrac{1}{|\mathcal{Q}|}\sum\limits_{Q\in \mathcal{Q}}Q^N$.
We can bound the redundancy as follows:
\begin{align*}
\inf_{\Model} \redundency_N(\Model,\datafamily)
   \le & \redundency_N(Q_0,\datafamily)
   =\int_{\datafamily} \mathcal{P}(\dataModel)\DKL(P^N_{\dataModel}||Q_0) d\dataModel \\ =& \int_{\datafamily} \mathcal{P}(\dataModel)\mathbb{E}\left[\log \dfrac{ P^N_{\dataModel}}{\sum\limits_{Q\in \mathcal{Q}}Q^N}+\log |\mathcal{Q}|\right]  d\dataModel 
    \\ \le& \log |\mathcal{Q}|+ \int_{\datafamily} \mathcal{P}(\dataModel)\mathbb{E}\left[\log \dfrac{ P_{\dataModel}^N}{\max\limits_{Q\in \mathcal{Q}} 
    Q^N}\right]d\dataModel
    \\ \le& \log |\mathcal{Q}|+ N\int_{\datafamily} \mathcal{P}(\dataModel)\min_{Q\in \mathcal{Q}}\mathbb{E}\left[\log \dfrac{ P_{\dataModel}}{
    Q}\right]d\dataModel
    \\ =& \log |\mathcal{Q}|+ N\int_{\datafamily} \mathcal{P}(\dataModel)\min_{Q\in \mathcal{Q}}\DKL(P_{\dataModel}||Q)d\dataModel
    \\ \le & \log |\mathcal{Q}|+ N\int_{\datafamily} (\Len_m\log|\mathbb{K}|+3\epsilon^2) d\dataModel
    \\ = & \log |\mathcal{Q}|+ 3N\epsilon^2 + N\mathbb{E}[\Len_m]\log |\mathbb{K}|
\end{align*}
Plugging in the bound of $|\mathcal{Q}|$
and $\mathbb{K}$, we can see the above is 
upper bounded by
\[N\mathbb{E}[\Len_m]\log |\mathbb{K}| + 3N\epsilon^2 +  2m\log \dfrac{m}{\epsilon^2}+2m\knwdimension\log\dfrac{\knwLip\knwdimension}{\epsilon}+n_s\syndimension\log\dfrac{\synLip\syndimension}{\epsilon}.
\]
% We bound the two terms in the above expression separately, starting with the first term $\int_{\knwfamily} \mathcal{P}(\knwModel)\DKL(P^N_{\knwModel}(K)||Q(K)) d\knwModel$. 
Choosing $\epsilon=N^{-1}$ and applying the result in \Cref{lem:L_j}, we know that 
\begin{align*}
    \inf_{\Model} \redundency_N(\Model,\datafamily)&\le O(N^{\alpha}(\log |\knwset|+\knwdimension(\log N+\log \knwLip+\log \knwdimension)))
\\ &+O(n_s\syndimension(\log \synLip+\log \syndimension+\log N)).
\end{align*}
Therefore, we know that 
\begin{align*}
    \inf_{\Model}\dfrac{1}{N}\redundency_N(\Model,\datafamily)=\dfrac{1}{N}\mathbb{I}(X_{1:N}; \dataModel) = 
    \widetilde{O}\left(\frac{\knwdimension}{N^{1-\alpha}}+\frac{ n_s\syndimension}{N}\right).
\end{align*}
This completes the proof.
\end{proof}



% Note that a simple upper bound for $|Q_{\epsilon}|$ is $|Q_{\epsilon}|\le \lfloor n/\epsilon\rfloor^n$.

% Assume that $Q_1,Q_2,\cdots,Q_{m}$ is a $\epsilon$-covering of the parametric family $\knwfamily$ where $m \le \left(\dfrac{2L\sqrt{d_k}}{\epsilon}\right)^{d_k}$.

% Let $Q(K)=\dfrac{1}{|Q_{\lfloor n^{\alpha}\rfloor,\epsilon}|}\sum\limits_{(q_1,\cdots,q_{\lfloor n^{\alpha}\rfloor})\in Q_{\lfloor n^{\alpha}\rfloor,\epsilon}} \left(\sum\limits_{i=1}^n q_i \dfrac{1}{\knwcover} \sum\limits_{j=1}^{\knwcover} Q_{\text{knw}}^{(j)}\right)^N$.

% Let $Q(K)= \dfrac{\sum\limits_{j=1}^{\knwcover} \left(0.5Q_u+0.5Q_{\text{knw}}^{(j)}\right)^N}{\knwcover+1} $ where $Q_u$ is the uniform distribution over all possible knowledge $K$.

% \begin{align*}
%     \DKL(P^N_{\knwModel}(K)||Q(K))  &= \log (\knwcover)+\mathbb{E}_{K\sim P}\left[\log \dfrac{P^N_{\knwModel}}{\sum_{i=1}\left(Q^{(i)}_{\text{knw}}\right)^N} \right]
%     \\ &\le \log + \mathbb{E}_{K\sim P}\left[\log \dfrac{P^N_{\knwModel}}{\sum_{i=1}Q^{(i)}_{\text{knw}}} \right]
% \end{align*}

% Then by \Cref{lem:mixture KL}, we know that
% \begin{align*}
%    & \int_{\knwModel} \mathcal{P}(\knwModel)\DKL(P^N_{\knwModel}(K)||Q(K)) d\knwModel
%     \\ \le & \int_{\knwModel} \mathcal{P}(\knwModel)\DKL(P^N_{\knwModel}(K)||Q(K)) d\knwModel
% \end{align*}

% \end{proof}


% \subsection{Proofs of \Cref{sec: scaling law}}
% \begin{theorem}[Restatement of \Cref{thm: data scaling lower bound}]
%     Under (\uppercase\expandafter{\romannumeral 1}) in \Cref{ass:data law}, the total expected redundancy $\testdataerror$ satisfies
%     \begin{align*}
%     \testdataerror(N) = \Omega(N^{-1+\alpha}).
% \end{align*}
% Under (\uppercase\expandafter{\romannumeral 2}) in \Cref{ass:data law}, the total expected redundancy $\testdataerror$ satisfies
%     \begin{align*}
%     \testdataerror(N) = O(N^{-1+\alpha}).
% \end{align*}
% \end{theorem}






% \begin{theorem}[Restatement of \Cref{thm: general scaling law}]
%       Denote $R(N,M)$ as the optimal value of the following optimization problem:
%       \begin{align*}
%     \text{minimize} \quad & \mathbb{E}_{k,t_k}[\ferror(k,m_k^*,t_k)] = \sum\limits_{k=1}^{\infty} p_k \mathbb{E}_{t_k}[\ferror(k,m_k^*,t_k)], \nonumber \\
%     \text{subject to} \quad & \sum\limits_{k=1}^{\infty} m_k^* \le M,  \quad  m_k^* \ge 0 \ \ \text{for all} \ k\in \mathcal{N}^+.
% \end{align*}
% Under \Cref{assumption: f}, we have 
% \begin{align*}
%     R(N,M) = \Theta(N^{-1+1/\alpha}) + \Theta(M^{-\alpha+1}).
% \end{align*}
% \end{theorem}
% \begin{proof}[Proof of \Cref{thm: general scaling law}]
%     By \Cref{assumption: f},  we know that 
%     \[\mathbb{E}_{k,t_k}[\ferror(k,m_k^*,t_k)]=\mathbb{E}_k[f_{\text{model}}(k,m_k^*)]+\mathbb{E}_{k,t_k}[f_{\text{data}}(k,t_k)].\]

%     Notice that $\mathbb{E}_{k,t_k}[f_{\text{data}}(k,t_k)$ does not depend on $m_k^*$. 
    
%     By \Cref{thm: data scaling lower bound} and \Cref{thm: data scaling upper bound}, we know that $\mathbb{E}_{k,t_k}[f_{\text{data}}(k,t_k)]=\Theta(N^{-1+1/\alpha})$.

%     Then this optimization problem reduces to the same optimization problem as \eqref{model scaling}, plus a $\Theta(N^{-1+1/\alpha})$ term that does not depend on the optimization variables $m_k^*$.

%     By \Cref{thm:model scaling}, we know that 
%     \begin{align*}
%     R(N,M) = \Theta(N^{-1+1/\alpha}) + \Theta(M^{-\alpha+1}).
% \end{align*}
% \end{proof}





\end{document}
